\input{PreamblePub}
\setcounter{page}{237}
\def\hindsection{4}

% Chp. 4 - N3,J5,G4; P11,C6,V3; N23,J7; V19,E2,C7(partly); G11,N19,V15; J10,Jyo

\newcommand{\xinclude}[1]{} %
\newcommand{\xexclude}[1]{#1} %
\TeXtoTEI{xinclude}{}
\TeXtoTEI{xexclude}{ab}
% \newcommand{\orgx}[2]{#1}
\newcommand{\orgvnr}[1]{}
\newcommand{\myfnx}[1]{}

\begin{document}
\thispagestyle{empty}
%\begin{center}
\section*{Chapter 4}
%\extramarksright{Chapter 4}
%\end{center}
\bigskip
\begin{otherlanguage}{iast}
\begin{ekdosis}


\teimute{\setcounter{saved@poemline}{1}}
%<*vs1a>
\begin{ava}[hp04_001a]
\app{\lem[nolem]{}
	\rdg[wit={N3,J5,P11,C6}]{\only}}% G4 broken
atha samādhiḥ/ 
\end{ava}
%</vs1a>

\vspace{-.5ex}
\avacite{1a}\bigskip

\iffalse
\startaltnormal
%<*vs0-1>
\begin{alttlg}[hp04_000_1]
\tl{\app{\lem[nolem]{}
	\rdg[wit={Gamma,Delta1}]{\also}}%
\pada{\app{\lem[wit={ceteri}]{namaḥ}
	\rdg[wit={N23,E2,V3}]{oṃ namaḥ}% +K3,C7
	} śivāya gurave}
\pada{nādabindu%
	\app{\lem[wit={Gamma,Delta,G5,J10,C6,Jyo}]{kalātmane}
		\rdg[wit={G11,Zeta2,P11,V3}]{layātmane}% +M3,F; °tmano? V3
		}/}\\+}
\tl{
\pada{\app{\lem[wit={ceteri}]{nirañjanapadaṃ}
		\rdg[wit={N23}]{\om}
		\rdg[wit={V3}]{nirañjanaṃ padaṃ}
		}
	\app{\lem[wit={Gamma,Delta1,G11,V15,J10,Jyo}]{yāti}
		\rdg[wit={G5,N19,Pi}]{yānti}}}
\pada{\app{\lem[wit={J7,V19,Epsi,Zeta2,Pi,Jyo}]{nityaṃ}
		\rdg[wit={N23}]{aharniśaṃ}
		\rdg[wit={E2}]{yatra}
		\rdg[wit={J10}]{yato}
		} % +K3,C7
	\app{\lem[wit={V19,Epsi,V15,P11,V3,Jyo}]{yatra}
		\rdg[wit={Gamma,N19}]{yatna}% yatta? N23; +F
		\rdg[wit={E2}]{nityaṃ}% +K3,C7
		\rdg[wit={J10}]{yogī}
		\rdg[wit={C6}]{ca yat}
%		\rdg[wit={Jyo}]{tatra}
		}%
	\app{\lem[wit={Gamma,Delta,V15,J10,V3,Jyo}]{parāyaṇaḥ}
		\rdg[wit={Epsi,N19,P11,C6}]{parāyaṇāḥ}%##
		}//\versenr}%
	\xinclude{\note[type=anmkg, labelb={note0},nonum,lem=\texteng{Note:}]{For the conventions specific to this edition, please refer to the introduction §\,5.3.4. The grey-scaled verses are usually only found in manuscript groups \textepsilon, \textzeta, \texteta, \textpi, and \textchi. Notes about the omission or inclusion of these verses are only provided when there is a deviation from this pattern.}}
	\\!}
\end{alttlg}
%</vs0-1>

%<*vs0-2>
\begin{alttlg}[hp04_000_2]
\tl{\app{\lem[nolem]{}
	\rdg[wit={Gamma,Delta1}]{\also}}%
\pada{\app{\lem[wit={ceteri}]{athedānīṃ} % °dāniṃ V15
		\rdg[wit={N23}]{athekṣanīṃ}
		\rdg[wit={V3}]{athodānī}
		}
	pravakṣyāmi} % vakṣāmi N19,N23,V19,V3
\pada{samādhikrama%
	\app{\lem[wit={Epsi,Zeta2,J10,Pi,Jyo},alt={°m uttamam}]{\skp{°}m uttamam}
		\rdg[wit={Gamma,Delta}]{lakṣaṇam}}/}\\+}
\tl{
\pada{mṛtyughnaṃ % ghaṃ N23
	\app{\lem[wit={Gamma,E2,Epsi,Pi}]{tu}
		\rdg[wit={V19}]{su}
		\rdg[wit={Zeta2,J10,Jyo}]{ca}
		} sukhopāyaṃ} % muṣopāpaṃ N23, mukho° N19
\pada{brahmānandakaraṃ 
	\app{\lem[wit={ceteri}]{param}
		\rdg[wit={G5}]{sadā}}//\versenr}\\!}% brahmanāṃda J7ac; para J10ac
\end{alttlg}
%</vs0-2>

%\startaltrecension
%\teimute{\small}
% \begin{alttlg}[hp04_000_3]
% \tl{%\vin
% \pada{\app{\lem[wit={G11,V15,Jyo}]{rājayogaḥ}
		% \rdg[wit={G5,N19,J10,C6}]{rājayoga}% +F
		% }
	% \app{\lem[wit={Epsi,J10,C6,Jyo}]{samādhiś ca}% +G7
		% \rdg[wit={Zeta2}]{samādhiḥ syād}% samādhi V15; +F,N22,P6
		% }}
% \pada{\app{\lem[wit={ceteri}]{unmanī}
	% \rdg[wit={G11}]{py unmanī}} ca manonmanī/}\\+}
% \tl{%\vin
% \pada{\app{\lem[wit={V15,J10}]{amaraugho}
		% \rdg[wit={G11}]{amaraughā}
		% \rdg[wit={C6}]{amaraughi}
		% \rdg[wit={N19}]{avaraubhū}
		% \rdg[wit={Jyo}]{amaratvaṃ}
		% \rdg[wit={G5}]{aromaro}}
	% \app{\lem[wit={G11,N19,J10,C6,Jyo},alt={layas}]{laya\skp{s}}
		% \rdg[wit={V15}]{layes}
		% \rdg[wit={G5}]{yas tat}}%
	% \app{\lem[wit={G11,Zeta2,C6,Jyo},alt={tattvaṃ}]{\skm{s }tattvaṃ}
		% \rdg[wit={J10}]{tatra}
		% \rdg[wit={G5}]{tulyaḥ}}}
% \pada{\app{\lem[wit={G11,Zeta2,J10,Jyo}]{śūnyāśūnyaṃ} % °śūnya N19
		% \rdg[wit={C6}]{śūnyāc chūnyaṃ}
		% \rdg[wit={G5}]{śūnyāt śūnya}}
		% paraṃ padam//\versenr}\label{synonym3}
	% \sgwit{G11,G5,N19,V15,J10,C6,Jyo} \anm{\textleftarrow\ \ref{A1}}\\!}
% \end{alttlg}

% \begin{alttlg}[hp04_000_4]
% \tl{%\vin
% \pada{amanaskaṃ tathādvaitaṃ}
% \pada{nirālambaṃ nirañjanam/}\\+}
% \tl{%\vin
% \pada{jīvanmuktiś ca % jīvamuktiś N19, °muktaś G5
	% \app{\lem[wit={Epsi,N19,J10,C6,Jyo}]{sahajaṃ}
% %		\rdg[wit={Jyo}]{sahajā}
		% \rdg[wit={V15}]{\om}}}
% \pada{\app{\lem[wit={V15,C6}]{turyaṃ} % +J11
		% \rdg[wit={G5}]{tulyaṃ}
		% \rdg[wit={N19}]{turyai}
		% \rdg[wit={Jyo}]{turyā}
		% \rdg[wit={G11}]{turīyaṃ}
		% \rdg[wit={J10}]{muktiś}}
	% \app{\lem[wit={J10pc,Jyo}]{cety ekavācakāḥ}
		% \rdg[wit={J10ac},alt={°kaḥ}]{cety ekavācakaḥ}
		% \rdg[wit={G5,C6}]{cety ekavācakam}% +F; caitye° C6
		% \rdg[wit={V15}]{cittaikavācakam}% +N22,P6
		% \rdg[wit={N19}]{ciṃtaikavācakam}
		% \rdg[wit={G11}]{caikavācakaṃ}
		% }//\versenr}\label{synonym4}
% \sgwit{G11,G5,N19,V15,J10,C6,Jyo} \anm{\textleftarrow\ \ref{A2}}%
% \myfn{\getsiglum{C6} has these verses on synonyms both here and at \ref{A1}/\ref{A2}, but \getsiglum{P11} has them at the latter place only.}\\!}
% \end{alttlg}
%\endaltrecension
%\teimute{\normalsize}

%<*vs0-3>
\begin{alttlg}[hp04_000_3]
\tl{\app{\lem[nolem]{}
	\rdg[wit={Gamma,Delta1}]{\also}}%
\pada{salile saindhavaṃ
	\app{\lem[wit={ceteri},alt={yadvat}]{yadva\skp{t}}
		\rdg[wit={N19}]{tadvat}}}%
\pada{t sāmyaṃ \app{\lem[wit={Gamma,Delta,J10,C6,Jyo}]{bhajati}
		\rdg[wit={Epsi,Zeta2}]{bhavati}% +F
		\rdg[wit={P11}]{ttadgati}
		\rdg[wit={V3}]{bhajata}
		} yogataḥ/}\\+} % °ta N19
\tl{
\pada{\app{\lem[wit={ceteri}]{tathā}
		\rdg[wit={J10}]{yathā}
		\rdg[wit={V3}]{athā}
		}%
	\app{\lem[wit={ceteri},alt={°tmamanasor}]{\skp{°}tmamanaso\skp{r}} % °saur N23
		\rdg[wit={J10}]{tmānamanor}
		}r aikyaṃ} % ma om. P7
\pada{samādhiḥ % °dhi P11, °dhir G11,G5,J10,Jyo
	\app{\lem[wit={ceteri}]{so}
		\rdg[wit={Epsi,J10,Jyo}]{a}
		\rdg[wit={P11}]{sā}
		}%
	\app{\lem[wit={ceteri}]{'bhidhīyate}
		\rdg[wit={N23}]{vidhīyate}
		}//\versenr}\label{salile}\\!}
\end{alttlg}
%</vs0-3>
%\myfn{In the following, not all of the differences in the verse order of \getsiglum{Pi} and \getsiglum{Jyo} are noted.}% \getsiglum{Pi} follow the order of \getsiglum{Gamma} (or of \getsiglum{Delta}?) in the beginning and the end (after 4.72). The middle part is a kind of mix of \getsiglum{Gamma} and \getsiglum{N19,V15}. The verse order of \getsiglum{Jyo} is similar to that of \getsiglum{N19,V15}, but with many small differences.}%


%\startaltrecension
%\teimute{\small}
%<*vs0-4>
\begin{alttlg}[hp04_000_4]
\tl{\app{\lem[nolem]{}
 \rdg[wit={Pi}]{\NotIn}}%
\pada{\app{\lem[wit={G5,Zeta2}]{yat samatvaṃ}
	\rdg[wit={G11}]{tat samatvaṃ}
	\rdg[wit={J10,Jyo}]{tat samaṃ ca}% jat J10pc
	}
	dvayo%r
	\app{\lem[wit={Epsi},alt={atra}]{\skm{r }atra}
	\rdg[wit={Zeta2}]{eva}
	\rdg[wit={J10,Jyo}]{aikyaṃ}
	}}
\pada{jīvātmaparamātmanoḥ/}\\+} % ātmā G11
\tl{%\vin
\pada{\app{\lem[wit={G5,Zeta2,J10}]{samastanaṣṭa}
	\rdg[wit={G11}]{samastaṃ naṣṭa}
	\rdg[wit={Jyo}]{pranaṣṭasarva}}%
\app{\lem[wit={Epsi,V15}]{saṃkalpaḥ}
	\rdg[wit={N19,J10}]{saṃkalpa}
	\rdg[wit={Jyo}]{saṃkalpaṃ}}}
\pada{samādhiḥ 
\app{\lem[wit={G11,Zeta2,J10,Jyo}]{so}
		\rdg[wit={G5}]{sā}
		}'bhidhīyate//\versenr}% 
\myfn{After this verse, \getsiglum{J10} has another similar verse:
\vspaceinnote\\
\devnote{karpūraṃ salile yadvat saindhavaṃ salile yathā/
tathātmamanasor aikyaṃ samādhiḥ so'bhidhīyate//} (cf. \ref{karpura}ab and \ref{salile}cd)}\\!}
\end{alttlg}
%</vs0-4>
%\endaltrecension
%\teimute{\normalsize}


%<*vs0-5>
\begin{alttlg}[hp04_000_5]
\tl{\app{\lem[nolem]{}
	\rdg[wit={Gamma,Delta1}]{\also}}%
\pada{rājayogasya
	\app{\lem[wit={ceteri}]{māhātmyaṃ}
		\rdg[wit={J7}]{māhatmyaṃ}
		\rdg[wit={V15}]{mahā}}}
\pada{ko vā jānāti tattvataḥ/}\\+} % ki N19; cā P11; jānaṃti N19
\tl{
\pada{\app{\lem[wit={ceteri},alt={jñānān}]{jñānā\skp{n}}
		\rdg[wit={V19}]{jñān}
		\rdg[wit={V15,J10}]{jñāna}
		\rdg[wit={Jyo}]{jñānaṃ}
		}%
	\app{\lem[wit={Gamma,E2,C6,Jyo},alt={muktiḥ}]{\skm{n }muktiḥ}% +K1
		\rdg[wit={V19,Epsi,Zeta2,J10,P11,V3}]{mukti}}
	\app{\lem[wit={Epsi}]{sthirā}% +M3
		\rdg[wit={Gamma,E2,J10,C6,Jyo}]{sthitiḥ}
		\rdg[wit={V19}]{sthiti<<ḥ>>}
		\rdg[wit={N19,V3}]{sthite}% +F,K1
		\rdg[wit={V15}]{s tato}
		\rdg[wit={P11}]{sthitai}% sthitā P6
		}
	\app{\lem[wit={Zeta2,P11,C6,Jyo},alt={siddhir}]{siddhi\skp{r}}% +F,K1
		\rdg[wit={Gamma,Delta,Epsi}]{siddhā}
		\rdg[wit={J10,V3}]{siddhi}
		}}%
\pada{r guru\app{\lem[wit={ceteri}]{vākyena} % gurur P11
		\rdg[wit={N23}]{vākyāt <<pra>>}}
	\app{\lem[wit={ceteri}]{labhyate}
		\rdg[wit={J10}]{sidhyati}}//\versenr}\\!}
\end{alttlg}
%</vs0-5>

%<*vs0-6>
\begin{alttlg}[hp04_000_6]
\tl{\app{\lem[nolem]{}
	\rdg[wit={Gamma,Delta1}]{\also}}%
\pada{durlabho viṣayatyāgo} % durlabha C6; viṣayāt* yogo N23
\pada{durlabhaṃ tattvadarśanam/}\\+} % labha V3
\tl{
\pada{durlabhā sahajāvasthā} % dullā G11, tu sahāvasthā G5
\pada{sadguroḥ karuṇāṃ vinā//\versenr}% guro P11,V3; karuṇā P11,N19
\\!}
\end{alttlg}
%</vs0-6>

%<*vs0-7>
\begin{alttlg}[hp04_000_7]
\tl{\app{\lem[nolem]{}
	\rdg[wit={Gamma,Delta1}]{\also}}%
\pada{vividhair % vicitrair G5
	āsanaiḥ % āsanai P11, āsanaḥ V15
	kumbhai}% °bhai N23,C6
\pada{\app{\lem[wit={G5,C6,Jyo},alt={vicitraiḥ}]{\skm{r }vicitraiḥ}% +K3,M3,F
		\rdg[wit={Gamma,Delta,Zeta2,J10,P11,V3}]{vicitra}
		\rdg[wit={G11}]{citraiś ca}}
	\app{\lem[wit={Delta,Epsi,J10,Pi,Jyo}]{karaṇair api}
		\rdg[wit={N23}]{kalaṇair api}
		\rdg[wit={J7}]{karuṇair api}
		\rdg[wit={Zeta2}]{karaṇair atha}}/}\\+}
\tl{
\pada{\app{\lem[wit={ceteri},alt={prabuddhāyām}]{prabuddhāyā\skp{m}}% pravudhā° V19, °yam G5
		\rdg[wit={N19}]{pradhadhāyām}}%
	\app{\lem[wit={ceteri},alt={ādi}]{\skm{m }ādi}
		\rdg[wit={V15}]{idaṃ}
		\rdg[wit={Jyo}]{mahā}}%
	\app{\lem[wit={ceteri}]{śaktau} % proktau J10ac
		\rdg[wit={N23}]{śaktiḥ}}}
\pada{prāṇaḥ śūnye % prāṇaṃ C6, prāṇa V15; sampra° G5
	\app{\lem[wit={N23,Delta,Epsi,J10,C6}]{vilīyate}
		\rdg[wit={J7}]{vidhīyate}
		\rdg[wit={Zeta2,P11,V3,Jyo}]{pralīyate}% +F
		}//\versenr}\\!}
\end{alttlg}
%</vs0-7>

%<*vs0-8>
\begin{alttlg}[hp04_000_8]
\tl{\app{\lem[nolem]{}
	\rdg[wit={Gamma,Delta1}]{\also}}%
\pada{\app{\lem[nolem]{\skp{pāda a}}
	\rdg[wit={C6}]{\om}}%
\app{\lem[wit={ceteri}]{utpanna}
		\rdg[wit={N23}]{ut<<pan>>na}
		\rdg[wit={V19}]{utpannā}
		}%
	\app{\lem[wit={ceteri}]{śaktibodhasya}
		\rdg[wit={N23}]{śaktibodhaḥ syāt}
		\rdg[wit={V15}]{śaktibodhaś ca}}}
\pada{\app{\lem[nolem]{\skp{pāda b}}
	\rdg[wit={C6}]{\om}}%
\app{\lem[wit={ceteri}]{tyakta}
		\rdg[wit={N23}]{prakṣa}}%
	niḥśeṣakarmaṇaḥ/}\\+}
\tl{
\pada{\app{\lem[wit={ceteri}]{yoginaḥ}
	\rdg[wit={C6}]{yogināṃ}} sahajāvasthā}
\pada{svaya%m
	\app{\lem[wit={Epsi,V15,J10,P11,V3},alt={eva prakāśate}]{\skm{m }eva prakāśate}% +K1; prakāśyate N22,P6
		\rdg[wit={Gamma,Delta,C6,Jyo}]{eva prajāyate}
		\rdg[wit={N19}]{eva prakāśayet}
		}//\versenr}\\!}
\end{alttlg}
%</vs0-8>

%<*vs0-9>
\begin{alttlg}[hp04_000_9]
\tl{\app{\lem[nolem]{}
	\rdg[wit={Gamma,Delta1}]{\also}}%
\pada{suṣumṇā\app{\lem[wit={ceteri}]{vāhini}% °<<m>>ṇā N23
		\rdg[wit={N23,G5,N19,V3}]{vāhinī}
		\rdg[wit={V19}]{vāhi}}
	\app{\lem[wit={ceteri}]{prāṇe}
		\rdg[wit={V3}]{prāṇa}}}
\pada{\app{\lem[wit={Epsi,V15,P11}]{śūnyaṃ}
		\rdg[wit={Gamma,Delta,C6,Jyo}]{śūnye}% +K1,N22,P6
		\rdg[wit={N19}]{śūnyā}
		\rdg[wit={J10}]{śūnya}% śūnyaṃ J10pc
		\rdg[wit={V3}]{śūne}
		} 
	\app{\lem[wit={ceteri}]{viśati}
		\rdg[wit={P11}]{vasati}}
	\app{\lem[wit={G11,P11,V3,Jyo}]{mānase} % +M1,M3
		\rdg[wit={Gamma,Delta,G5,Zeta2,C6}]{mārute}% mārutai N23; +K1,N22,P6,G7
		\rdg[wit={J10}]{mārutaḥ}
		}\marma/}\\+} 
\tl{
\pada{\app{\lem[wit={Gamma,Delta,G11}]{tathā}
		\rdg[wit={G5,Zeta2,J10,Pi,Jyo}]{tadā}}
	\app{\lem[wit={ceteri}]{samasta}
		\rdg[wit={J10,Jyo}]{sarvāṇi}}karmāṇi}
\pada{\app{\lem[wit={ceteri}]{nirmūlayati}
		\rdg[wit={N23}]{nirmūlaṃ yāti}
		\rdg[wit={V19,V15}]{nimūlayati}
		\rdg[wit={G5}]{nirmalaṃ yāti}} % °ḷayati V15
	\app{\lem[wit={G11,N19,J10,Pi}]{marmavit}% K1,M1,G7
		\rdg[wit={N23,G5,V15}]{karmavit}% P6,P22,M3,F
		\rdg[wit={J7}]{karmakṛt}
		\rdg[wit={Delta,Jyo}]{yogavit}
		}//\versenr}\\!}
\end{alttlg}
%</vs0-9>

%<*vs0-10>
\begin{alttlg}[hp04_000_10]
\tl{
\pada{\app{\lem[nolem]{\skp{pāda a}}
	\rdg[wit={Gamma,Delta1}]{\also}
	}%
	\app{\lem[wit={G5,V15,V3}]{amaraugha}% +F
		\rdg[wit={Gamma}]{amano nir}
		\rdg[wit={Delta}]{amalo nir}
		\rdg[wit={G11}]{amareśa}
		\rdg[wit={N19,P11}]{amarogha}% +M3?
		\rdg[wit={J10,Jyo}]{amarāya}% +C2
		\rdg[wit={C6}]{amaraughi}
		} % verse om. K1,P6
	\app{\lem[wit={Epsi,Zeta2,J10,Pi,Jyo}]{namas tubhyaṃ}
	\rdg[wit={Gamma}]{manāḥ śūnyaṃ}
	\rdg[wit={Delta}]{malaḥ śūnyaṃ}}%
	}
\pada{so'pi \app{\lem[wit={G11,N19,C6,V3,Jyo}]{kālas tvayā}% trayā V3
	\rdg[wit={G5}]{kālasya vā}
	\rdg[wit={V15}]{kāla tvayā}
	\rdg[wit={J10}]{kālantayā}
	\rdg[wit={P11}]{kālaṃ tvayā}
	}
\app{\lem[wit={G11,Zeta2,J10,Pi}]{hataḥ}% hata N19
	\rdg[wit={G5}]{hakaḥ}
	\rdg[wit={Jyo}]{jitaḥ}
	}/}\\+}
\tl{
\pada{patitaṃ \app{\lem[wit={Epsi,Zeta2,Pi,Jyo}]{vadane}
	\rdg[wit={J10}]{pavane}
	} yasya} % yasyā P11
\pada{\app{\lem[nolem]{\skp{pāda d}}
	\rdg[wit={Gamma,Delta1}]{\also}}%
	jagad etac carācaram//\versenr} % yagad V19; carāccaraṃ J7, °care P11
	\\!}
\end{alttlg}
%</vs0-10>

%<*vs0-11>
\begin{alttlg}[hp04_000_11]
\tl{\app{\lem[nolem]{}
	\rdg[wit={Gamma,Delta1}]{\also}}%
\pada{citte
	\app{\lem[wit={ceteri},alt={samatvam}]{samatva\skp{m}}
		\rdg[wit={N23}]{samatyam}
		\rdg[wit={Zeta2}]{śamatvam}
		}m āpanne}
\pada{\app{\lem[wit={J7,Delta,Epsi,N19,Jyo}]{vāyau}
		\rdg[wit={N23,V3}]{vāyor}
		\rdg[wit={V15}]{vāyo}
		\rdg[wit={J10,C6}]{vāyur}
		\rdg[wit={P11}]{vāyu}
		}
	\app{\lem[wit={ceteri}]{vrajati}
		\rdg[wit={N23}]{javati}} madhyame/}\\+}
\tl{
\pada{\app{\lem[nolem]{\skp{pāda c}}
		\rdg[wit={Gamma}]{\om}}%
	\app{\lem[wit={Epsi,N19}]{tadāmaraugha}
		\rdg[wit={V19}]{eṣā naulīti}
		\rdg[wit={E2}]{eṣā naulī ca}% +C7
		\rdg[wit={V15}]{tadāmaroḷi}
		\rdg[wit={J10}]{tathāmarolī}
		\rdg[wit={P11,V3}]{eṣāmaraugha}
		\rdg[wit={C6}]{saivāmarolī}
		\rdg[wit={Jyo}]{tadāmarolī}
		}% 
	\app{\lem[wit={V19,G5,N19,J10,Pi,Jyo}]{vajrolī}% +G5,M3,K3,C7
		\rdg[wit={E2}]{vajrī ca}
		\rdg[wit={G11}]{vajrolīs}
		\rdg[wit={V15}]{vajrolis}
		}}
\pada{\app{\lem[nolem]{\skp{pāda d}}
		\rdg[wit={Gamma}]{\om}}%
	\crux\app{\lem[wit={Epsi,Zeta2},alt={tadāśā jīvite 'pi ca}]{tadāśā jīvite'pi ca}% +M3
		\rdg[wit={Delta}]{sadā cābhimateti ca}
		\rdg[wit={J10}]{sahajolī mato pi ca}
		\rdg[wit={Pi}]{sadā me bhimateti ca} % bhimate cita V3; timateti vaḥ P11
		\rdg[wit={Jyo}]{sahajolī prajāyate}}\crux//\versenr}
		\\!}
\end{alttlg}
%</vs0-11>
% G7 tadha manā vajroci tādhāśrājiṃtasya ca (*12 om.)
% M3 tato maśā .. vajraulī tadāśājīvitepi ca
% K1 eṣāmarolī vajrolī sadā abhinayāti ca
% P6 eṣāmaroli vajroli sadā abhinayaṃti ca (*12 om.)
% N22 eṣāmaroli vajroli sadā savitayiti ca (*12 om.)
% P11 eṣāmaraughavajrolī sadā me timateti vaḥ


%<*vs0-12>
\begin{alttlg}[hp04_000_12]
\tl{\app{\lem[nolem]{}
	\rdg[wit={Gamma,Delta1}]{\also}}%
\pada{jñānaṃ 
	\app{\lem[wit={ceteri}]{kuto}
		\rdg[wit={G11}]{tato}} manasi
	\app{\lem[wit={Gamma,Delta,J10,Pi}]{jīvati devi yāvat} % jāvat V19, yā<<va>>t N23
		\rdg[wit={G11,N19}]{jīvati devi tāvat}
		\rdg[wit={G5}]{jīvati tepi tāvat}
		\rdg[wit={V15}]{jīvati durvikalpe}
		\rdg[wit={Jyo}]{saṃbhavatīha tāvat}
		}}\\+}
\tl{
\pada{\app{\lem[wit={ceteri},alt={prāṇo 'pi}]{prāṇo'pi}
		\rdg[wit={G11,V15,C6}]{prāṇe pi}
		\rdg[wit={G5}]{prāṇeha}
		} jīvati mano
	\app{\lem[wit={ceteri}]{mriyate} % sriyate N23, mrīyate P7
		\rdg[wit={J7,V19}]{mṛyate}
		\rdg[wit={G5}]{priyate}
		\rdg[wit={V15}]{miyata}
		}
	\app{\lem[wit={ceteri}]{na}
		\rdg[wit={N19}]{ca}}
	\app{\lem[wit={Pi}]{tāvat}
		\rdg[wit={ceteri}]{yāvat}% +F
		}/}\\+}
\tl{
\pada{\app{\lem[wit={ceteri}]{prāṇo}
		\rdg[wit={Delta}]{prāṇaṃ}}
	\app{\lem[wit={ceteri}]{mano}
		\rdg[wit={Epsi,N19}]{'pi ca}
		} dvayam idaṃ % dvayām N23
	\app{\lem[wit={ceteri}]{vilayaṃ}
		\rdg[wit={V15}]{na vilī}}
	\app{\lem[wit={P11,C6}]{prayāti}% +K1,P6,F,source
		\rdg[wit={N23}]{jayed  yo}
		\rdg[wit={J7}]{naved yo}
		\rdg[wit={Delta,Jyo}]{nayed yo}
		\rdg[wit={G11}]{nayet taṃ}% nayete M3; gate cen M1; nayeta? G3
		\rdg[wit={G5}]{na yattat}
		\rdg[wit={N19}]{na yāvat}
		\rdg[wit={V15}]{yate tra}
		\rdg[wit={J10}]{na yāti}% +G7,K3
		\rdg[wit={V3}]{prajāti}
		}}\\+}
\tl{
\pada{mokṣaṃ
	\app{\lem[wit={ceteri}]{sa}
		\rdg[wit={V15}]{na}
		\rdg[wit={C6}]{ca}
		} gacchati % gacchatiti V19
		naro na kathaṃci%d % narā N23
	\app{\lem[wit={ceteri},alt={anyaḥ}]{\skm{d }anyaḥ}
		\rdg[wit={G5}]{anyam}
		\rdg[wit={J10}]{anyat}
		\rdg[wit={V3}]{anya}}//\versenr}\\!}
\end{alttlg}
%</vs0-12>


%<*vs0-13>
\begin{alttlg}[hp04_000_13]
\tl{\app{\lem[nolem]{}
	\rdg[wit={Gamma,Delta1}]{\also}}%
\pada{\app{\lem[wit={ceteri}]{rasasya}% +P6,M3
		\rdg[wit={J7,Zeta2}]{rasaś ca}} % +G7
	\app{\lem[wit={ceteri}]{manasaś caiva} % caivaṃ G11,J10
		\rdg[wit={N23}]{manasaiva caṃ}
		\rdg[wit={V3}]{manaś caiva}
		}}
\pada{\app{\lem[wit={ceteri}]{cañcalatvaṃ}
		\rdg[wit={N23}]{calatvaṃ ca}
		\rdg[wit={N19}]{vaṃcatvaṃ ca}} svabhāvataḥ/}\\+} % °ta P11
\tl{
\pada{\app{\lem[wit={G5,V15}]{rasabandhe}% +M3
		\rdg[wit={N23}]{rase baddho}
		\rdg[wit={J7,Delta,J10,Pi,Jyo}]{raso baddho}
		\rdg[wit={G11}]{rasabaddhe}% +G7; rasabaddho? F
		\rdg[wit={N19}]{rase bandhe}
		} 
	mano\app{\lem[wit={V15}]{bandhe}% +M3
		\rdg[wit={J7,N23,Delta,N19,J10,V3,Jyo}]{baddhaṃ} % varddhaṃ N23
		\rdg[wit={G11}]{baddhe}
		\rdg[wit={G5}]{dhatte}
		\rdg[wit={P11}]{baddhaḥ}% +F
		\rdg[wit={C6}]{baddho}
		}}
\pada{\app{\lem[wit={ceteri}]{kiṃ}
		\rdg[wit={N19}]{tan}}
	na sidhyati bhūtale//\versenr}\\!} % siddhyaṃti V3
\end{alttlg}
%</vs0-13>
%<*vs0-14>
\begin{alttlg}[hp04_000_14]
\tl{\app{\lem[nolem]{}
	\rdg[wit={Gamma,Delta1}]{\also}}%
\pada{mūrchito % mūrchato E2
	\app{\lem[wit={Gamma,Delta,Zeta2,C6,Jyo}]{harate}
		\rdg[wit={Epsi,J10,P11,V3}]{harati}% +F
		}
	\app{\lem[wit={ceteri}]{vyādhiṃ}% vyādhin P11
		\rdg[wit={G5,Jyo}]{vyādhīn}
		\rdg[wit={J10,V3}]{vyādhi}
		}}
\pada{mṛto
	\app{\lem[wit={ceteri}]{jīvayati}
		\rdg[wit={V15}]{jīvayate}}
		svayam/}\\+}
\tl{
\pada{baddhaḥ % badho P11, bandhaḥ G5
	\app{\lem[wit={ceteri}]{khecaratāṃ}
		\rdg[wit={V19}]{khacatāṃ}}
	\app{\lem[wit={ceteri}]{dhatte}
		\rdg[wit={N23,N19}]{dhartte}
		\rdg[wit={V3}]{yāti}}}
\pada{\app{\lem[wit={ceteri}]{raso vāyuś ca}
		\rdg[wit={J10}]{sa jīveśvara}
		\rdg[wit={V3}]{vāyuś ca}
		}
	\app{\lem[wit={Delta,C6}]{bhairavi}
		\rdg[wit={Gamma,Epsi,Zeta2}]{bhairavī}
		\rdg[wit={J10}]{seśvaraḥ}
		\rdg[wit={P11}]{tad dvayaṃ}
		\rdg[wit={V3},post=\texteng{(tathā \textapp{for missing} raso)}]{bhairavī tathā}
		\rdg[wit={Jyo}]{pārvati}
		}//\versenr}\\!}
\end{alttlg}
%</vs0-14>

\endaltnormal
\fi


%<*vs1>
\begin{tlg}[hp04_001]
\tl{
\pada{\app{\lem[wit={ceteri}]{indriyāṇāṃ} % °nāṃ V19, °ṇā N3
		\rdg[wit={N19}]{indriyāṇi}} mano nātho} % nāthe G11
\pada{\app{\lem[wit={N3,J5,Epsi,Pi,Jyo}]{manonāthas tu}
		\rdg[wit={G4}]{manonāthasu}
		\rdg[wit={N23,Delta,V15,J10}]{manonāthaś ca}
		\rdg[wit={J7}]{manaso nātha}
		\rdg[wit={N19}]{manonāthaḥ su}
		} mārutaḥ/}\\+} % mārute P11
\tl{
\pada{mārutasya layo
	\app{\lem[wit={ceteri},alt={nāthas/nāthaḥ/nātho}]{nātha\skp{s/nāthaḥ/nātho}}
		\rdg[wit={J7}]{nāthāḥ}}}%
\pada{\app{\lem[wit={N3,J5,Epsi,Zeta2,J10,V3},alt={taṃ nāthaṃ layam āśrayet}]{\skm{s }taṃ nāthaṃ layam āśrayet}% K1A; nātha N3
		\rdg[wit={G4}]{tan nātho laya\,+\,+\,+}
		\rdg[wit={Gamma,E2,C6,Jyo}]{sa layo nādam āśritaḥ}% K1B
		\rdg[wit={V19}]{layo dasamāśrayaḥ}
		\rdg[wit={P11}]{laya nātha niraṃjanāṃ}% tasya nātho nirañjanaḥ F
		}//\versenr}
		\orgvnr{1}\\!}
\end{tlg}
%</vs1>

%\avacite{1a}
\commcite\newpage


\iffalse
\startaltnormal
%<*vs1-1>
\begin{alttlg}[hp04_001_1]
\tl{\app{\lem[nolem]{}
	\rdg[wit={Gamma,Delta1}]{\also}}%
\pada{\app{\lem[wit={Epsi,V15,Pi,Jyo},alt={so 'yam evāstu}]{so'yam evāstu}
		\rdg[wit={Gamma,Delta}]{ayam eva tu}% +F
		\rdg[wit={N19}]{soyamo vāstu}
		\rdg[wit={J10}]{svayam evāstu}% +K1,P6
		} % evaṃ N23
	\app{\lem[wit={ceteri}]{mokṣākhyo} % °kṣyo N19, °syo P11, °khyā G5
		\rdg[wit={J10}]{vā mokṣaḥ}
		}}
\pada{\app{\lem[wit={Epsi,V15,Pi,Jyo}]{māstu vāpi}% +P11; astu K1,P6
		\rdg[wit={N23}]{aya vāpi}
		\rdg[wit={J7}]{'stu vāpi sa}
		\rdg[wit={V19}]{yas tu vāpi}% +K3,C7
		\rdg[wit={E2}]{yas tu vyāpi}
		\rdg[wit={N19}]{māstu kapi}
		\rdg[wit={J10}]{sosti vāpi}
		}
	matāntare/}\\+} % matātare J7, matāṃbare P11, mātā° J10ac
\tl{
\pada{manaḥprāṇa% mana P11
	\app{\lem[wit={Gamma,Epsi,V15,P11,C6}]{layānando} % layāṃnado N23; layo? V15
		\rdg[wit={Delta}]{layo nādo}
		\rdg[wit={N19}]{layānanda}
		\rdg[wit={J10}]{m apānaṃ vā}
		\rdg[wit={V3}]{layāna} % 1 syllable too short
		\rdg[wit={Jyo}]{laye kaścid}
		}}
\pada{\app{\lem[wit={G11,Zeta2,P11,C6}]{mayi}% +K1,M1,M3,P6
		\rdg[wit={Gamma,Delta}]{nāpi}
		\rdg[wit={G5}]{bhuvi}
		\rdg[wit={J10}]{layaḥ}
		\rdg[wit={V3}]{māpi}
		\rdg[wit={Jyo}]{āna}
		}
	\app{\lem[wit={ceteri},alt={kaścit/°cid}]{kaści\skp{t/°cid}}
		\rdg[wit={V19}]{kviṃcid}
		\rdg[wit={Jyo}]{ndaḥ saṃ}}%
	\app{\lem[wit={G11,Zeta2,J10,P11,C6,Jyo},alt={pravartate}]{\skm{t }pravartate}% °vattate G11
		\rdg[wit={N23}]{vibhedyate}
		\rdg[wit={J7,Delta}]{vibhidyate}
		\rdg[wit={G5}]{pravartatām}
		\rdg[wit={V3}]{pravartate na}
		}//\versenr}\\!}
\end{alttlg}
%</vs1-1>
\endaltnormal
\fi


%<*vs2>
\begin{tlg}[hp04_002]
\tl{
\pada{\app{\lem[wit={J7,Delta,Epsi,V15,J10,V3},alt={pranaṣṭo-/praṇaṣṭocchvāsa}]{pranaṣṭocchvāsa}% ṇaṣṭa Delta,G11,+F (pranaṣṭa is richtig; see PW)
		\rdg[wit={N3,Jyo}]{praṇaṣṭaśvāsa}% na Jyo ##?
		\rdg[wit={J5}]{praṇaṣṭabhyāsa}
		\rdg[wit={N23}]{prabhṛṣṭo\,\_\,sa}
		\rdg[wit={N19}]{pranaṣṭauśvāsa}
		\rdg[wit={P11}]{pranaṣṭosvāsa}
		\rdg[wit={C6}]{pranaṣṭaḥ svā<<sa>>}
		}%
	\app{\lem[wit={N3,Epsi,V15,Jyo}]{niśvāsaḥ}
		\rdg[wit={J5,V3}]{niśvāsa}
		\rdg[wit={J7}]{niśvāsāḥ}
		\rdg[wit={N23}]{niśvāsā}
		\rdg[wit={Delta,C6ac}]{niḥśvāsa}% niḥsvāsa V19
		\rdg[wit={N19,J10,P11,C6pc}]{niḥśvāsaḥ} % niḥñcāsaḥ N19, °svāsaḥ P11
		}}
\pada{\app{\lem[wit={ceteri}]{pradhvasta}% praddhasta N19; +M3
		\rdg[wit={G11}]{prabhṛṣṭa}
		\rdg[wit={J10}]{pranaṣṭa}% °naṣṭaḥ V17, prā° N26.
		}%
	\app{\lem[wit={ceteri}]{viṣaya}
		\rdg[wit={G11}]{viṣayā}
		\rdg[wit={N19}]{viṣaga}}%
	\app{\lem[wit={N3,J5,V19,Epsi,J10,C6,V3,Jyo}]{grahaḥ}
		\rdg[wit={Gamma,E2}]{grahāḥ}
		\rdg[wit={N19}]{hvaraḥ}
		\rdg[wit={V15}]{jvaraḥ}
		\rdg[wit={P11}]{grataḥ}
		}/}\\+}
\tl{
\pada{\app{\lem[wit={Epsi,C6,V3,Jyo}]{niśceṣṭo}
		\rdg[wit={N3}]{niḥśceṣṭo}
		\rdg[wit={J5}]{niścaiṣṭo}
		\rdg[wit={Gamma,Delta,V15}]{niśceṣṭā}
		\rdg[wit={N19}]{nidyeṣṭo}
		\rdg[wit={J10}]{niścalo}
		\rdg[wit={P11}]{niḥśreṣṭo}
		}
	\app{\lem[wit={N23,Epsi,Zeta2,J10,Pi,Jyo}]{nirvikāraś ca} % nirvikā<<ra>>ś N23, nivikāraś V15
		\rdg[wit={N3}]{nirvikāras tu}
		\rdg[wit={J5}]{nivikalpas tu}
		\rdg[wit={J7,Delta}]{nirvikārāś ca}
		}}
\pada{\app{\lem[wit={ceteri}]{layo}
		\rdg[wit={Gamma,E2}]{layaṃ}
		\rdg[wit={V19}]{laye}
		}
	\app{\lem[wit={ceteri}]{jayati}
		\rdg[wit={Gamma,Delta}]{yānti ca}
		\rdg[wit={G5}]{layati}}
	\app{\lem[wit={N3,J5,Epsi,Zeta2,Pi,Jyo}]{yoginām}% °naṃ N3
		\rdg[wit={Gamma,Delta,J10}]{yoginaḥ}}//\versenr}
		\orgvnr{2}\\!}
\end{tlg}
%</vs2>
\commcite\newpage

%<*vs3>
\begin{tlg}[hp04_003]
\tl{\app{\lem[nolem]{}
	\rdg[wit={E2}]{\om}}%
\pada{\app{\lem[wit={ceteri}]{ucchinna}
		\rdg[wit={N3,G11,V15}]{ucchinnaḥ}
		\rdg[wit={V19}]{ucchūna}% +K3,C7
		}%
	sarva\app{\lem[wit={ceteri}]{saṃkalpo}
		\rdg[wit={V19}]{saṃkalpe}}}
\pada{\app{\lem[wit={ceteri}]{niḥśeṣāśeṣa}% ni<<ḥ>>śeṣā° J10
		\rdg[wit={J5,V3}]{niḥśeṣośeṣa}
		\rdg[wit={Gamma}]{niḥśeṣagata}
		}%
	\app{\lem[wit={ceteri}]{ceṣṭitaḥ}
		%\rdg[wit={K3,C7}]{veṣṭitaḥ}
		\rdg[wit={V15}]{varjitaḥ}
		\rdg[wit={C6}]{ceṣṭitam}
		}/}\\+}
\tl{
\pada{\app{\lem[wit={N3,J5,V19,J10,V3,Jyo}]{svāvagamyo}
		\rdg[wit={G4,Epsi,P11}]{svāvagamya}
		\rdg[wit={Gamma}]{svāgate cā}
		\rdg[wit={N19}]{svāgamyo}
		\rdg[wit={V15}]{svānugamyo}
		\rdg[wit={C6}]{sovagamyo}
		}
		layaḥ ko'pi} % pi added in margin C7; laya kaupi P11
\pada{\app{\lem[wit={Alpha,G11,C6}]{jayatāṃ vāgagocaraḥ}% °ra J5
		\rdg[wit={Gamma,V19}]{manovācām agocaraḥ}% +K3,C7
		\rdg[wit={G5}]{jagatāṃ vāgagocaraḥ}
		\rdg[wit={N19}]{japatāṃ vāgagocara}
		\rdg[wit={V15}]{jāyatāṃ vāgagocaraḥ}
		\rdg[wit={J10,V3,Jyo}]{jāyate vāgagocaraḥ} % agocara V3
		\rdg[wit={P11}]{jāyatāṃ cāpi gaucaraḥ}
		}//\versenr}
		\orgvnr{3}\\!}
\end{tlg}
%</vs3>
\commcite\newpage


%<*vs4>
\begin{tlg}[hp04_004]
\tl{\app{\lem[nolem]{}
	\rdg[wit={E2}]{\om}}%
\pada{yatra % yabha N23, yanna N19
	\app{\lem[wit={ceteri},alt={dṛṣṭir}]{dṛṣṭi\skp{r}}
		\rdg[wit={N3,G5,V15,J10}]{dṛṣṭi}
		\rdg[wit={C6}]{vṛṣṭir}
		}r layas tatra} % yayas? G11
\pada{bhūtendriya% bhute° N19
	\app{\lem[wit={N3,J5,Epsi,V15,V3}]{sanātanaḥ}
		\rdg[wit={Gamma,V19,J10,C6,Jyo}]{sanātanī}% +K3,C7
		\rdg[wit={N19}]{sanātanaṃ}
		\rdg[wit={P11}]{sanātana}
		}/}\\+}
\tl{
\pada{\app{\lem[wit={N3,Gamma,V19},alt={syāc chaktir/°tiḥ}]{syāc chakti\skp{r}}% °kti<<ḥ>> N23; +K3,C7
		\rdg[wit={J5}]{syāt saktir}
		\rdg[wit={Epsi,N19,J10,Pi,Jyo}]{sā śaktir}% śakti P11,G11; +F
		\rdg[wit={V15}]{sa śaktir}}%
	\app{\lem[wit={N3,J5,Epsi,J10,Pi,Jyo},alt={jīva}]{\skm{r }jīva}
		\rdg[wit={Gamma,V19}]{sarva}% +K3,C7
		\rdg[wit={Zeta2}]{bhāva}}%
	\app{\lem[wit={ceteri}]{bhūtānāṃ}
		\rdg[wit={N23}]{bhūtānī}
		\rdg[wit={N19}]{bhūnāṃ}}}
\pada{\app{\lem[wit={N3,G4,Gamma,J10,C6,V3},alt={dṛṣṭir}]{dṛṣṭi\skp{r}}
		\rdg[wit={J5,V19,Epsi,P11}]{dṛṣṭi}% +K3,C7
		\rdg[wit={Zeta2}]{dṛṣṭe}
		\rdg[wit={Jyo}]{dve a}}%
	\app{\lem[wit={Epsi,N19,P11,V3},alt={lakṣ(y)e layaṃ gatā}]{\skm{r }lakṣye layaṃ gatā}% lakṣe P11,V3,G5,N19
		\rdg[wit={N3}]{lakṣe layaṃ gatāḥ}
		\rdg[wit={J5}]{lakṣe la(!) gatā}
		\rdg[wit={G4}]{lakṣy[e] layaṃ gataḥ}
		\rdg[wit={N23}]{lakṣana saṃgatā}
		\rdg[wit={J7}]{lakṣe na saṃgatā}
		\rdg[wit={V19}]{lakṣeṇa saṃgatā}%lakṣyeṇa K3,C7
		\rdg[wit={V15}]{lakṣaṃ layaṃ gatau}
		\rdg[wit={J10,Jyo}]{lakṣye layaṃ gate}
		\rdg[wit={C6}]{gacchel layaṃ gate}}//\versenr}%
		%\myfnx{After this verse, \getsiglum{Jyo} has \ref{layo}.}
		\label{yatradrsti}
		\orgvnr{4}
		\\!}
\end{tlg}
%</vs4>
\commcite\newpage


%<*vs5>
\begin{tlg}[hp04_005]
\tl{
\pada{\app{\lem[nolem]{\skp{pāda a}}
		\rdg[wit={J5,V3}]{\om}}%
	vedaśāstra\app{\lem[wit={N3,G4,Epsi,Zeta2,J10,P11,C6,Jyo}]{purāṇāni}
		\rdg[wit={N23}]{purāṇādyāḥ}
		\rdg[wit={J7}]{puraṇādyāḥ}
		\rdg[wit={V19}]{purāṇaiś ca}
		\rdg[wit={E2}]{purāṇaughāḥ}% +K3,C7
		}}
\pada{\app{\lem[nolem]{\skp{pāda b}}
		\rdg[wit={J5,V3}]{\om}}%
	\app{\lem[wit={ceteri}]{sāmānya}% °yaṃ G11, °nyā E2
		\rdg[wit={C6}]{samāni}}% 
	gaṇikā iva/}% gaṇivā V19
		\\+}
\tl{
\pada{\app{\lem[nolem]{\skp{pāda c}}
	\rdg[wit={V3}]{\om}}%
\app{\lem[wit={ceteri}]{ekaiva}
		\rdg[wit={E2}]{idaṃ tu}
		} śāmbhavī % sāṃ° J10
	\app{\lem[wit={ceteri}]{mudrā}
		\rdg[wit={V15}]{māyā}
		\rdg[wit={J10}]{vidyā}}}
\pada{\app{\lem[nolem]{\skp{pāda d}}
	\rdg[wit={V3}]{\om}}%
\app{\lem[wit={N3,J5,Gamma,G5,P11,C6,Jyo}]{guptā kulavadhūr iva}% vadhū iva P7
		\rdg[wit={Delta}]{sarvatantreṣu gopitā rakṣaṇīyā prayatnena guptā kulavadhūr iva}
		\rdg[wit={G11,Zeta2},post=\texteng{(cf.\,\ref{gopita}d)}]{sarvatantreṣu gopitā}
		\rdg[wit={J10}]{gopyā kulavadhūr iva}
		}//\versenr}\label{vedasastra}
		\orgvnr{5}\\!}
\end{tlg}
%</vs5>
\commcite\newpage

% atha śāmbhavī (Jyo)


%<*vs6>
\begin{tlg}[hp04_006]
\tl{\app{\lem[nolem]{}
	\rdg[wit={Zeta2}]{\om}}% \anm{eye-skip?}
\pada{anta\app{\lem[wit={J5,Gamma,G5,J10,C6ac,V3,Jyo},alt={lakṣ(y)aṃ}]{\skm{r }lakṣyaṃ} % antalakṣaṃ P7; °lakṣaṃ V3,N23,J10; °lakṣyaṃ Jyo,C6ac
		\rdg[wit={N3,G11,P11,C6pc}]{lakṣ(y)a}
		\rdg[wit={V19}]{lakṣā} % lakṣyā K3
		\rdg[wit={E2}]{lakṣyo}
		} % y C7,P11,C6
	\app{\lem[wit={ceteri},alt={bahir}]{bahi\skp{r}}
		\rdg[wit={J10}]{mano}}%
	\app{\lem[wit={ceteri},alt={dṛṣṭir}]{\skm{r }dṛṣṭi\skp{r}}
		\rdg[wit={J5,V19,G11,J10,V3}]{dṛṣṭi}}}%
\pada{\app{\lem[wit={ceteri},alt={nimeṣonmeṣa}]{\skm{r }nimeṣonmeṣa}
		\rdg[wit={N23,P11}]{nirmiṣonmeṣa}% °ṣya N23
		}\app{\lem[wit={ceteri}]{varjitā}
		\rdg[wit={E2,P11}]{varjjitaḥ}}/}\\+}
\tl{
\pada{\app{\lem[wit={N3,Epsi,P11,C6,Jyo}]{eṣā sā}% sāṃ P11
		\rdg[wit={J5}]{eṣāsau}
		\rdg[wit={Gamma,V19}]{saiṣā tu}% +K3,C7
		\rdg[wit={E2}]{eṣā vai}
		\rdg[wit={J10}]{eṣā tu}
		\rdg[wit={V3}]{eṣā hi}
		}
		śāmbhavī mudrā} % sāṃbhavī J10
\pada{\app{\lem[wit={ceteri}]{sarvatantreṣu}% sarve J10
		\rdg[wit={V19}]{sarvatantreṣu śastreṣu}
		%\rdg[wit={K3,C7}]{sarvaśāstreṣu}
		\rdg[wit={Jyo}]{vedaśāstreṣu}}
		gopitā//\versenr}\label{gopita}
		\orgvnr{6}\\!}
\end{tlg}
%</vs6>
\commcite\newpage
		%\myfn{The omission was probably caused by haplography. \getsiglum{L2} inserts the following passage after 4.20a: \devnote{eṣā tu śāṃbhavī vidyā gopā ca kulavadhūr iva/ aṃtarlakṣamanodṛṣṭi nimeṣonmeṣonmeṣa(!)varjito/}}\\!

%<*vs7>
\begin{tlg}[hp04_007]
\tl{
\pada{anta\app{\lem[wit={N3,Delta,Epsi,J10,P11,C6,Jyo},alt={lakṣya}]{\skm{r}lakṣya}
		\rdg[wit={J5,Gamma,Zeta2,V3}]{lakṣa}}%
		vilīnacittapavano yogī
	\app{\lem[wit={ceteri}]{yadā}
		\rdg[wit={Alpha,N19}]{sadā}% +F
		\rdg[wit={G5}]{yamā}
		\rdg[wit={J10}]{yathā}
		} vartate}\\+} % ##?
\tl{
\pada{\app{\lem[wit={ceteri}]{dṛṣṭyā}
		\rdg[wit={J10}]{dṛṣṭvā}
		\rdg[wit={P11}]{dṛṣyā}
		\rdg[wit={V3}]{dṛśyā}}
	niścala\app{\lem[wit={ceteri}]{tārayā}
		\rdg[wit={N23}]{tāra}
		\rdg[wit={P11}]{tālayā}
		}
	% <ba>hir N23; bahidharaḥ! G11
	bahi\app{\lem[wit={Alpha,Epsi,V15,J10,Pi,Jyo},alt={adhaḥ}]{\skm{r }adhaḥ}
		\rdg[wit={Gamma,Delta}]{asau}
		\rdg[wit={N19}]{adhraḥ}
		}
	\app{\lem[wit={J5,Delta,Epsi,Zeta2,Jyo}]{paśyann apaśyann api}
		\rdg[wit={N3}]{paśyann apaśyann ivā}% illeg. G4 ##?
		\rdg[wit={Gamma}]{paśyan na paśyaty api}
		\rdg[wit={J10}]{paśyann api}
		\rdg[wit={P11,V3}]{paśyan na paśyet tataḥ}% tata V3; paśyet tadā F
		\rdg[wit={C6}]{paśyen na paśyet tataḥ}}/}\\+}
\tl{
\pada{\app{\lem[wit={ceteri}]{mudreyaṃ}
		\rdg[wit={V15}]{mudre}
		} khalu % khaluṃ N23
	\app{\lem[wit={N3,J5,J10,P11,V3}]{khecarī} % = source; +K1,P6
		\rdg[wit={Gamma,Delta,Epsi,Zeta2,C6,Jyo}]{śāmbhavī}% +F
		}
	\app{\lem[wit={ceteri}]{bhavati sā}% +M3
		\rdg[wit={G11}]{ti kathitā}
		\rdg[wit={V3}]{bhavati}
		}
	\app{\lem[wit={N3,J5,Delta,Zeta2,V3},alt={yuṣmat}]{yuṣma\skp{t}}% +J10pc
		\rdg[wit={N23}]{puṣpat}
		\rdg[wit={J7}]{<<yu>>ṣmat}
		\rdg[wit={Epsi,C6}]{yasya}
		\rdg[wit={J10}]{yuṣmān}
		\rdg[wit={P11}]{yāsya}
		\rdg[wit={Jyo}]{labdhā}
		}tprasādā%d
	\app{\lem[wit={Gamma,V19,V15,J10,P11,V3},alt={guro}]{\skm{d }guro}% +K3,C7
		\rdg[wit={N3}]{gurau}
		\rdg[wit={J5}]{gure}
		\rdg[wit={E2,Epsi,N19,C6,Jyo}]{guroḥ} % J10pc
		}}\\+}
\tl{
\pada{\app{\lem[wit={ceteri}]{śūnyāśūnya}% °nyaṃ N19
		\rdg[wit={C6}]{śūnyāc chūnya}}%
	\app{\lem[wit={ceteri}]{vivarjitaṃ}
		\rdg[wit={J5}]{vivarjito}
		\rdg[wit={N23}]{vivarjite}
		\rdg[wit={V19}]{vivarjiti}
		\rdg[wit={Jyo}]{vilakṣaṇaṃ}}
	\app{\lem[wit={ceteri}]{sphurati}
		\rdg[wit={V19}]{spharati}}
	\app{\lem[wit={ceteri},alt={yat}]{ya\skp{t}}% pat N23
		\rdg[wit={N3,Jyo}]{tat}% ##?
		\rdg[wit={J5}]{ttat}
		\rdg[wit={V19}]{[pta]t}
		\rdg[wit={V3}]{ya}
		}t
	tattvaṃ \app{\lem[wit={ceteri}]{padaṃ}
		\rdg[wit={G11,N19}]{\om}
		\rdg[wit={G5}]{paraṃ}
		} śāmbhavam//\versenr}
		\label{antarlaksya}
		\orgvnr{7}\\!}
\end{tlg}
%</vs7>
\commcite\newpage


%<*vs8>
\begin{tlg}[hp04_008]
\tl{\app{\lem[nolem]{}
	\rdg[wit={Zeta2,J10}]{\om}}% om. G3
\pada{śrīśāmbha\app{\lem[wit={N3,J7,Delta,G5,Jyo},alt={°vyāś ca khecaryā}]{\skp{°}vyāś ca khecaryā} % khecarayā J7, °caryyāḥ V19
		\rdg[wit={J5}]{vyā khecaryā}
		\rdg[wit={G4}]{vavyā khecaryā}
		\rdg[wit={N23}]{vyāḥ khecaryā\,\_}
		\rdg[wit={G11}]{vāś ca khecaryā}
		\rdg[wit={Pi}]{vyā(ḥ) khecaryāś ca} % vyā V3,P6
		}}
\pada{\app{\lem[wit={P11}]{avasthāyām abhedatā}% +K1 (ana°)
		\rdg[wit={N3,G11}]{avasthāyāṃ na bhedataḥ}% bhedatā M3
		\rdg[wit={J5}]{avasthāyasya bhedataḥ}
		\rdg[wit={G4,G5}]{avasthāyā na bhedataḥ}
		\rdg[wit={Gamma}]{avasthā ca na bhedataḥ\,(°naḥ \getsiglum{N23})}
		\rdg[wit={Delta}]{avasthā balabhedataḥ}
		\rdg[wit={C6}]{hy avasthāyām abhedataḥ}
		\rdg[wit={V3}]{avasthāyāṃ ca bhedatā}
		\rdg[wit={Jyo}]{avasthādhāmabhedataḥ}
		}\marma//\versenr}%
		\myfn{After this verse, \getsiglum{Jyo} has an additional line: \devnote{bhavec cittalayānandaḥ śūnye citsukharūpiṇi/}} % +M1,G8
		\orgvnr{8}\\!}
\end{tlg}
%</vs8>
\commcite\newpage

% J5 śobhavyā khecaryā avaschāyasya bhedataḥ
% G4 śrīśāṃbhavavyā khecaryā avasthāyā {{tavalabdhi}} na bhedataḥ
% K1a! śrīśāṃbhavyāś ca khecarya anasthāyām abhedatā
% K1b! śrīśāṃbhavyāḥ khecaryāś ca avasthā tu na bhedataḥ
% M1 śrīśāṃbhavyāś ca khecar.ā avasthāyām abh[e]dataḥ
% M3 śriśāṃbhavyāś ca khecaryā avastāyān na bhedatā
% G5 śriśāṃbhavyāś ca khecaryā avastāyā na bhedataḥ
% G7 śāṃbaryāś ca khecarya atastā naiva bhedataḥ
% P6 śrīśaṃbhavyā khecaryāś ca avasthā tu na bhedataḥ
% N22 śrīśaṃbhavyā khecaryāś ca avasthā <!> na bhedataḥ
% IFP: śriśāṃbhavyāś ca khecaryā avasthā sthānabhedataḥ


%<*vs9>
\begin{tlg}[hp04_009]
\tl{\app{\lem[nolem]{}
	\rdg{\texteng{=\,3.48*1}}
	\rdg[wit={Alpha,Gamma,Delta2}]{\incl}%
 	\rdg[wit={Gamma,Delta2}]{\textapp{found after X4.42}}}%
\pada{\app{\lem[resp=emend]{pātāle yad viśati}
		\rdg[wit={N3,J5}]{pātāle yadvitaya}% °tayas G11
		\rdg[wit={G4}]{pātāḷe yadvita\,..}
		\rdg[wit={Gamma}]{pātālād yad viśati}
		\rdg[wit={Delta2}]{pātālād vā viyati}
		}
	\app{\lem[wit={J5}]{suṣiraṃ}% sukhiraṃ K1,P6
		\rdg[wit={N3}]{suśiraṃ}
		\rdg[wit={N23}]{śikhiraṃ}
		\rdg[wit={J7}]{śikharaṃ}% +K3
		\rdg[wit={Delta2}]{śikhare}
		} merumūle
	\app{\lem[wit={N3}]{tad asmin}% +G11
		\rdg[wit={J5}]{yadismi}
		\rdg[wit={N23}]{tasti}
		\rdg[wit={J7}]{tad asti}% +K1,P6
		\rdg[wit={V19}]{tadāstā}
		\rdg[wit={C7}]{tad āste}% +K3
		}}\\+}
\tl{
\pada{tattvaṃ caitat pravadati % pravahati C7
	\app{\lem[wit={N3,Gamma}]{sudhīs tan mukhaṃ}
		\rdg[wit={J5}]{sudhī sanmukhaṃ}
		\rdg[wit={V19}]{susaṃmukhaṃ}
		\rdg[wit={C7}]{sudhīḥ saṃmukhaṃ}% +K3
		} 
		nimnagānām/}\\+} % ninmaganāṃ V19
\tl{
\pada{candrā%t % caṃdrā N3
	\app{\lem[wit={Gamma},alt={sāraḥ}]{\skm{t }sāraḥ}% sāra K1, sāraṃ P6
		\rdg[wit={N3,J5}]{sāro}
		\rdg[wit={Delta2}]{srāvaḥ}
		}
	\app{\lem[wit={N23,C7}]{sravati}% +K1,P6,G11,K3
		\rdg[wit={N3}]{grasati}
		\rdg[wit={J5}]{\om}
		\rdg[wit={J7}]{savati}
		\rdg[wit={V19}]{śravati}
		}
	\app{\lem[wit={N3,J5,N23,Delta2},alt={vapuṣas}]{vapuṣa\skp{s}}
		\rdg[wit={J7}]{puruṣas}}s
	tena mṛtyur narāṇāṃ}\\+} % mūtyur N23
\tl{
\pada{\app{\lem[wit={Alpha,J7,Delta2},alt={taṃ badhnīyāt}]{taṃ badhnīyā\skp{t}}
		\rdg[wit={N23}]{tadvahmaṃpāt}}%
	\app{\lem[wit={N3,J5},alt={sukaraṇamṛdā}]{\skm{t }sukaraṇamṛdā}
		\rdg[wit={G4}]{sukaraṇāmudā}
		\rdg[wit={N23}]{svakaraṇaimṛdā}
		\rdg[wit={J7,C7}]{svakaraṇamṛdā}% sakaraṇamṛtā K1,P6
		\rdg[wit={V19}]{svakaraṇamṛjā}
		} % sukaraṇam amṛtaṃ G11 (unm.)
		nānyathā
	\app{\lem[wit={N3,J7,C7}]{kāyasiddhiḥ}% +K1,G11
		\rdg[wit={J5,G4,V19}]{kāryasiddhi(ḥ)}% V19 om. h; +P6
		\rdg[wit={N23}]{kāyaḥ siddhiḥ}
		}//\versenr}
	\orgvnr{9}\\!}
\end{tlg}
%</vs9>
\commcite\newpage
%\ \newpage


%<*vs10>
\begin{tlg}[hp04_010]
\tl{\app{\lem[nolem]{}
	\rdg{\texteng{=\,3.77*1}}
	\rdg[wit={Alpha}]{\only}}%
\pada{yat kiṃcit sravate candrā}%
\pada{d amṛtaṃ divyarūpiṇaḥ/}\\+} % N24 has °rūpiṇaḥ too.
\tl{
\pada{tat sarvaṃ grasate sūrya}%
\pada{s tena piṇḍaṃ jarāyutam//\versenr} % piṇḍo jarāyutaḥ N24
\orgvnr{10}\\!} % almost the same J5; G4 damaged
\end{tlg}
%</vs10>
\commcite%\newpage


%<*vs11>
\begin{tlg}[hp04_011]
\tl{\app{\lem[nolem]{}
	\rdg{\texteng{=\,3.77*2}}
	\rdg[wit={Alpha}]{\only}}%
\pada{tatrāsti karaṇaṃ divyaṃ}
\pada{sūryasya \app{\lem[resp=emend]{mukhabandhanam}
	\rdg[wit={N3,J5}]{paribandhanaṃ}
	\rdg[wit={G4}]{\illeg}}/}\\+} % paripaṃthi ca N24
\tl{
\pada{gurūpadeśato jñeyaṃ}
\pada{na tu śāstrārthakoṭibhiḥ//\versenr}
%\sgwit{Alpha} 
%\anm{\textrightarrow\ \manuref{3.77*2}}
\orgvnr{11}\\!} % almost the same J5; G4 damaged
\end{tlg}
%</vs11>
\commcite\newpage


\iffalse
\startaltrecension
%<*vs11-1>
\begin{alttlg}[hp04_011_1]
\tl{\app{\lem[nolem]{}
	\rdg[wit={Delta1}]{\also}}%
\pada{\app{\lem[nolem]{\skp{pāda a}}
	\rdg[wit={Gamma}]{\also}}%
\app{\lem[wit={J7,E2,V15,P11,Jyo}]{tāre}
		\rdg[wit={N23}]{vāre}
		\rdg[wit={V19,V3}]{tāra}% E4
		%\rdg[wit={K3,C7}]{tāraṃ}
		\rdg[wit={G11}]{kalāṃ}
		\rdg[wit={G5}]{kalā}
		\rdg[wit={N19}]{tāvad}
		\rdg[wit={J10}]{tārā}% +F
		\rdg[wit={C6}]{tārāṃ}
		}
	\app{\lem[wit={Gamma,E2,Epsi,V15,C6,Jyo}]{jyotiṣi}% jotiṣi G5
		\rdg[wit={V19}]{jyotiso}
		\rdg[wit={N19}]{yotiṣi}
		\rdg[wit={J10}]{jyotiṣu}
		\rdg[wit={P11}]{jyotiṣīṃ}
		\rdg[wit={V3}]{jyotīṣa}
		}
	\app{\lem[wit={ceteri}]{saṃyojya}
		\rdg[wit={N23}]{samojyaṃ}
		\rdg[wit={V19}]{jojya}
		\rdg[wit={J10}]{saṃyojyā}
		}}
\pada{\app{\lem[nolem]{\skp{pāda b}}
	\rdg[wit={Gamma}]{\also}}%
kiṃci%d
	\app{\lem[wit={Epsi,V15,Pi,Jyo},alt={unnamayed}]{\skm{d }unnamaye\skp{d}} % °yet* V3; +M3,P6,N22
		\rdg[wit={N23,E2}]{uccālayed}
		\rdg[wit={J7}]{uccalayed}
		\rdg[wit={V19}]{uccācayed}
		\rdg[wit={N19}]{uṣṭānnama}
		\rdg[wit={J10}]{uccārayed}% +K1b
		}%
	\app{\lem[wit={ceteri},alt={bhruvau}]{\skm{d }bhruvau}% bhuvau? P11,J10
		\rdg[wit={N23}]{bhūvo<<ḥ>>}}/}\\+}
\tl{
\pada{\app{\lem[wit={E2,Epsi,Zeta2,P11,V3},alt={pūrvayogasya mārgo 'yam}]{pūrvayogasya mārgo'ya\skp{m}}% +K3,C7
		\rdg[wit={V19}]{pūrvayogasya māhātmyam}
		\rdg[wit={J10}]{sūryayogasya mārge ca}
		\rdg[wit={C6}]{pūrvayogasya mārgeṇa}
		\rdg[wit={Jyo}]{pūrvayogaṃ mano yuñjann}
		}}%
\pada{\app{\lem[wit={Delta,Epsi,Zeta2,P11,V3,Jyo},alt={unmanī}]{\skm{m }unmanī} % yaṃmunmanī N19
		\rdg[wit={G11}]{kiṃcid un}
		\rdg[wit={J10}]{yunmanī}
		\rdg[wit={C6}]{hy unmanī}
		}%
	\app{\lem[wit={G11,P11,Jyo}]{kārakaḥ kṣaṇāt}
		\rdg[wit={Delta,V15}]{karaṇaṃ kṣaṇāt}% +M3
		\rdg[wit={N19}]{kārakaṃ kṣaṇāt}
%		\rdg[wit={G11}]{namayet kṣaṇāt | unmanīkārakaḥ}
		\rdg[wit={G5}]{kāraṇaṃ kṣaṇāt}
		\rdg[wit={J10}]{kāralakṣaṇam}% +K1
		\rdg[wit={C6}]{kārakakṣaṇāt}
		\rdg[wit={V3}]{kāraṇaḥ kṣaṇāt}
		}//\versenr}
	\\!} % om. +P6,N22
\end{alttlg}
%</vs11-1>

%<*vs11-2>
\begin{alttlg}[hp04_011_2]
\tl{\app{\lem[nolem]{}
	\rdg[wit={Delta1}]{\also}}%
\pada{keci%d
	\app{\lem[wit={ceteri},alt={āgama}]{\skm{d }āgama}
		\rdg[wit={Epsi}]{nigama}}%
	\app{\lem[wit={ceteri}]{jālena}
		\rdg[wit={J10}]{yogena}
		}}
\pada{keci%n % cecin N19, keccin V19
	\app{\lem[wit={N19,J10,P11,C6,Jyo},alt={nigama}]{\skm{n }nigama}% =source
		\rdg[wit={Delta,V3}]{niyama}
		\rdg[wit={Epsi}]{āgama}
		\rdg[wit={V15}]{nima}
		}%
	\app{\lem[wit={Epsi,N19,J10,P11,C6,Jyo}]{saṃkulaiḥ}% saṃkulai P11
		\rdg[wit={V19}]{saṃkulā}
		\rdg[wit={E2}]{saṃkulāḥ}% +K3,C7
		\rdg[wit={V15,V3}]{saṃkule}
		}/}\\+}
\tl{
\pada{kecit tarkeṇa muhyanti} % ke<ci>nnarkkeṇa N19, tarkena E2
\pada{naiva jānanti tārakam//\versenr} % kārakam G5
%\myfn{Pādas ab and cd are transposed in \getsiglum{V19} and the correct order is indicated by a small number 1 and 2 above the hemistiches.}
\\!} % C6?
\end{alttlg}
%</vs11-2>


%\startaltnormal
%<*vs11-3>
\begin{alttlg}[hp04_011_3]
\tl{\app{\lem[nolem]{}
	\rdg[wit={Gamma,V19}]{\textapp{included after \ref{antarlaksya}}}}% Not in E2!
\pada{\app{\lem[wit={ceteri}]{ardhodghāṭita}
		\rdg[wit={N23}]{arddhocchā[d]ita}
		\rdg[wit={P11}]{arddhoghāṭita}
		\rdg[wit={Jyo}]{ardhonmīlita}}%
	\app{\lem[wit={V19,V15,Jyo}]{locanaḥ}
		\rdg[wit={Gamma,Epsi,N19,J10,Pi}]{locana}}
	\app{\lem[wit={ceteri}]{sthira}
		\rdg[wit={N23}]{sthila}}manā
	nāsāgradatte% nāśā° N19, nāśādagra° V17
	\app{\lem[wit={ceteri},alt={°kṣaṇaḥ/ś}]{\skp{°}kṣaṇaḥ\skp{/-aś}}% ś P11,C6,V19,J10
		\rdg[wit={N23,V3}]{kṣaṇāś}% °āḥś V3
		\rdg[wit={N19}]{kṣaṇaṃ}}}\\+}
\tl{
\pada{\app{\lem[wit={ceteri},alt={candrārkāv}]{candrārkā\skp{v}}
		\rdg[wit={J10}]{candrārkau}
		\rdg[wit={V3}]{cāndrārkāv}
		}%
	\app{\lem[wit={ceteri},alt={api}]{\skm{v }api}% N23,V19,G11,G5,V15,Pi,Jyo
		\rdg[wit={J7}]{avi}
		\rdg[wit={N19}]{aca}
		\rdg[wit={J10}]{ca vi}} līnatā%m
	\app{\lem[wit={G11,Jyo},alt={upanayan}]{\skm{m }upanaya\skp{n}}% M1,M3,F
		\rdg[wit={Gamma,V19,Zeta2}]{upanayen}% upanaye V19,N19,P6; eva nayet G7; +K1
		\rdg[wit={G5}]{apanayan}
		\rdg[wit={J10}]{gatau}
		\rdg[wit={Pi}]{upagatau}% = source
		}%
	\app{\lem[wit={ceteri},alt={niṣpanda}]{\skm{n }niṣpanda} % nispanda K3,C7,Jyo, niṣyaṃda? N19
		\rdg[wit={G5}]{diṣyanda}
		\rdg[wit={J10}]{nikṣipya}
		\rdg[wit={P11}]{nirvyaṃda}
		}%
	\app{\lem[wit={Epsi}]{bhāvāntare}% +C7
		\rdg[wit={N23,V19}]{bhāvo 'ntare}% bhāvottaraḥ P6,F
		\rdg[wit={J7}]{bhāvotare}
		\rdg[wit={N19}]{vāpyaṃ tataḥ}
		\rdg[wit={V15}]{bāṣpaṃ tataḥ}
		\rdg[wit={J10}]{bhāsoṃtare}
		\rdg[wit={P11}]{rūpaṃ tanu}% ##
		\rdg[wit={C6}]{rūpaṃ tataḥ}
		\rdg[wit={V3}]{rūpatanu}
		\rdg[wit={Jyo}]{bhāvena yaḥ}}/}\\+}
\tl{
\pada{jyotī% jyoti P11, jyoṃtī J7, yotī N19, jyotī{{kha}} J10
	\app{\lem[wit={ceteri},alt={rūpam}]{rūpa\skp{m}}
		\rdg[wit={J7}]{yatsyam}
		\rdg[wit={Zeta2}]{rūpa}
		}%
	\app{\lem[wit={ceteri},alt={aśeṣa}]{\skm{m }aśeṣa}% aśaiṣa P11
		\rdg[wit={Zeta2}]{viśeṣa}}%
	\app{\lem[wit={ceteri}]{bāhyarahitaṃ}
		\rdg[wit={Jyo}]{bījam akhilaṃ}}
	\app{\lem[wit={ceteri},alt={dedīpya°}]{dedīpya\skp{°}}
		\rdg[wit={N23}]{devadīpya}}mānaṃ paraṃ}\\+} % puraṃ G5
\tl{
\pada{tattvaṃ % tatva V3
	\app{\lem[wit={ceteri},alt={tat}]{ta\skp{t}}
		\rdg[wit={J10}]{yac}}%
	\app{\lem[wit={Gamma,V19,Jyo},alt={padam eti}]{\skm{t }padam eti}% +K1,P6
		\rdg[wit={Epsi,Pi}]{param eti}% +F
		\rdg[wit={Zeta2}]{param asti}% +HTK
		\rdg[wit={J10}]{carama}}
	\app{\lem[wit={ceteri}]{vastu}% +M1,M3,G7,K1,P6
		\rdg[wit={N23}]{vasta}
		\rdg[wit={P11,V3}]{yastu}
		\rdg[wit={C6}]{yat tu}} 
		paramaṃ % param avācyaṃ J7ac
	\app{\lem[wit={ceteri}]{vācyaṃ} % vācya V19
		\rdg[wit={N23}]{vāpyaṃ}} ki%m
	\app{\lem[wit={ceteri},alt={atrādhikam}]{\skm{m }atrādhikam}% °dikam G5
		\rdg[wit={N23}]{andrādhikaṃ}
		\rdg[wit={V19}]{atrāsanaṃ}}//\versenr}\label{ardhogha} 
	\\!}
\end{alttlg}
%</vs11-3>
%\endaltnormal



%\startaltrecension
%\teimute{\small}
%<*vs11-4>
\begin{alttlg}[hp04_011_4]
\tl{
\pada{\app{\lem[wit={Epsi,Zeta2,Pi,Jyo}]{divā na}
		\rdg[wit={J10}]{vāsare}} pūjayel liṅgaṃ}
\pada{\app{\lem[wit={N19,P11}]{rātrau naiva ca pūjayet}% +G7
		\rdg[wit={G11,C6,V3}]{rātrau naiva prapūjayet}
		\rdg[wit={G5,J10,Jyo}]{rātrau caiva na pūjayet} % pūyet J10; vāpi na M3
		\rdg[wit={V15}]{rātrau liṃgaṃ na pūjayet}}/}\\+}
\tl{
\pada{\app{\lem[wit={Epsi,Zeta2,J10,Pi}]{satataṃ}
		\rdg[wit={Jyo}]{sarvadā}} pūjayel liṅgaṃ} % pū<<ja>>yel J10
\pada{\app{\lem[wit={Jyo}]{divārātrinirodhataḥ}
		\rdg[wit={Epsi,P11,V3}]{divārātraṃ na pūjayet}% +M1,M3,F
		\rdg[wit={Zeta2,J10}]{divārātrau na pūjayet}
		\rdg[wit={C6}]{divārātrau ca pūjayet}% ca varjayet G3,G7
		}//\versenr}\\!}
\end{alttlg}
%</vs11-4>

%<*vs11-5a>
\begin{altava}[hp04_011_5a]
\app{\lem[nolem]{}
 \rdg[wit={P11,C6,Jyo}]{\only}}%
 atha \app{\lem[wit={C6,Jyo}]{khecarī}
	\rdg[wit={P11}]{khecarīsamādhiḥ}% +F
	}/ %\sgwit{P11,C6,Jyo}
\end{altava}
%</vs11-5a>
%<*vs11-5>
\begin{alttlg}[hp04_011_5]
\tl{\app{\lem[nolem]{}
	\rdg[wit={Jyo}]{\om}}%
\pada{\app{\lem[wit={Epsi,N19}]{suṣiro}% +F
		\rdg[wit={V15}]{dṛṅmukhaṃ}
		\rdg[wit={J10,V3}]{sukhiraṃ}
		\rdg[wit={P11}]{susthiro}
		\rdg[wit={C6}]{sukhiro}
		}
	jñāna\app{\lem[wit={Epsi,N19,P11,C6}]{janakaḥ}% +F
		\rdg[wit={V15}]{jaṃnakaṃ}
		\rdg[wit={J10,V3}]{janakaṃ}
		}}
\pada{pañcasrotaḥ% srota G5
%\app{\lem[wit={Epsi,V15,P11,C6}]{srotaḥ}
%	\rdg[wit={N19,J10,V3}]{śrotaḥ}}%
	\app{\lem[wit={Epsi,N19,P11,C6}]{samanvitaḥ}% +F
		\rdg[wit={V15}]{samanvitaṃ}% J10pc
		\rdg[wit={J10}]{samanvite}
		\rdg[wit={V3}]{samanvita}
		}/}\\+}
\tl{
\pada{tiṣṭhate khecarī mudrā} % khacarī G11
\pada{\app{\lem[wit={J10}]{tasmiñ śūnye}% +F; tasmin ś° J10
		\rdg[wit={Epsi,V15,P11,C6}]{tasmāc chūnye}
		\rdg[wit={N19}]{satyaṃ tatra}
		\rdg[wit={V3}]{\om\ (eye-skip?)}} % cf. V3 resumes with tasmin!
	\app{\lem[wit={Epsi,V15,J10,P11,C6}]{nirañjane}
		\rdg[wit={N19}]{na saṃśayaḥ}
		\rdg[wit={V3}]{\om}}//\versenr}%
	\myfn{\getsiglum{J10} and \getsiglum{V3} have this verse both in chapter three (after 3.48) and chapter four, while \getsiglum{Jyo} has it only in chapter three.}
	\\!}
\end{alttlg}
%</vs11-5>
%\myfn{\getsiglum{L2} omits the 2nd half of this verse and the 1st half of the next verse (prob. by haplography) and adds them after the 1st half of 4.25*3.}

%<*vs11-6>
\begin{alttlg}[hp04_011_6]
\tl{
\pada{\app{\lem[nolem]{\skp{pāda a}}
	\rdg[wit={V3}]{\om}}% % J6 has this omission too!!
savyadakṣiṇa\app{\lem[wit={V15,P11,C6}]{nāḍī}
		\rdg[wit={Epsi,N19,J10,Jyo}]{nāḍi}
		}stho}
\pada{\app{\lem[nolem]{\skp{pāda b}}
	\rdg[wit={V3}]{\om}}% % J6 has this omission too!!
\app{\lem[wit={Epsi,N19,J10,P11,C6,Jyo}]{madhye}
		\rdg[wit={V15}]{madhya}
		\rdg[wit={P11}]{madhyaṃ}
		}
		\app{\lem[wit={G11,N19}]{calati}
		\rdg[wit={G5,P11,C6,Jyo}]{carati}%  <<carati>> C6
		\rdg[wit={V15}]{carita}
		\rdg[wit={J10}]{vahati}
		}
		mārutaḥ/}% ḥ om. N19
		\\+}
\tl{
\pada{\app{\lem[nolem]{\skp{pāda c}}
	\rdg[wit={V3}]{\om}}% % J6 has this omission too!!
	tiṣṭhate khecarī mudrā} % tiṣṭhajña? N19
\pada{\app{\lem[wit={Epsi,V15,Pi,Jyo}]{tasmin sthāne}
		\rdg[wit={N19}]{satyaṃ tatra}
		\rdg[wit={J10}]{tatra satyaṃ}}
		na saṃśayaḥ//\versenr}\\!}
\end{alttlg}
%</vs11-6>

%<*vs11-7>
\begin{alttlg}[hp04_011_7]
\tl{\app{\lem[nolem]{}
	\rdg{\texteng{=\,3.37}}
	\rdg[wit={J10,Jyo}]{\NotIn}}%
\pada{cittaṃ carati khe yasmā}% citraṃ N19; cārati V15; yasmā V15,V3
\pada{j jihvā carati
	\app{\lem[wit={G5,N19,Pi}]{khe gatā}% ga<tā> P11
		\rdg[wit={G11}]{khe yadā}
		\rdg[wit={V15}]{vegataḥ}}/}\\+}
\tl{
\pada{\app{\lem[wit={G11,V15,P11,V3}]{tenaiṣā}% P6
		\rdg[wit={G5}]{iyaṃ ca}
		\rdg[wit={N19}]{tenaiva}
		\rdg[wit={C6}]{teneyaṃ}
		} khecarī  % khecarīṃ N19
	\app{\lem[wit={Epsi,N19,P11,V3}]{nāma}
		\rdg[wit={V15,C6}]{mudrā}}}
\pada{\app{\lem[wit={Epsi,N19,P11,V3}]{mudrā}% mudra P11
		\rdg[wit={V15}]{satyaṃ}
		\rdg[wit={C6}]{sarva}} 
		siddhai%r  % siddhai N19
	\app{\lem[wit={Epsi,N19,Pi},alt={namaskṛtā}]{\skm{r }namaskṛtā}% °tāḥ P11
		\rdg[wit={V15}]{nigadyate}}//\versenr}
	\\!}
\end{alttlg}
%</vs11-7>
%<*vs11-8>
\begin{alttlg}[hp04_011_8]
\tl{\app{\lem[nolem]{}
	\rdg[wit={V15,J10}]{\NotIn}}%
\pada{iḍāpiṅgalayo%r  % idā N19
	\app{\lem[wit={Epsi,N19,Pi},alt={yoge}]{\skm{r }yoge}
		\rdg[wit={Jyo}]{madhye}}}
\pada{\app{\lem[wit={Epsi,C6,Jyo}]{śūnyaṃ}
		\rdg[wit={N19,P11}]{śūnye}% +F,P6
		\rdg[wit={V3}]{śūne}} % +J11ac
	\app{\lem[wit={Epsi,N19,V3,Jyo}]{caivānilaṃ}
		\rdg[wit={P11,C6}]{caiva bilaṃ}}
	\app{\lem[wit={Epsi,N19,P11,V3,Jyo}]{graset}% grasit P11
		\rdg[wit={C6}]{viśet}}/}\\+}
\tl{
\pada{\app{\lem[wit={Epsi,N19,C6,V3,Jyo}]{tiṣṭhate}
		\rdg[wit={P11}]{tiṣṭhati}} khecarī mudrā} % khacarī G11
\pada{\app{\lem[wit={Epsi,P11}]{tatra satyaṃ na saṃśayaḥ}% P6,F
		\rdg[wit={N19}]{satyaṃ tatra na saṃśayaḥ}
		\rdg[wit={C6,V3,Jyo}]{tatra satyaṃ punaḥ punaḥ}
		}//\versenr}
		\\!}
\end{alttlg}
%</vs11-8>
%<*vs11-9>
\begin{alttlg}[hp04_011_9]
\tl{
\pada{\app{\lem[wit={Epsi,N19,J10},alt={somasūryadvayor}]{somasūryadvayo\skp{r}} % +G7,P6; sūryā J10
		\rdg[wit={V15}]{candrasūryadvayor}
		\rdg[wit={Pi,Jyo}]{sūryācandramasor}% sūrya C6, sūryāc P11, °masaur V3
		}r madhye}
\pada{\app{\lem[wit={Zeta2,C6,V3}]{nirālambe tale}% +G3
		\rdg[wit={G11}]{nirālambe kale}
		\rdg[wit={G5,P11}]{nirālambatale}% °laṃbatare M3,F
		\rdg[wit={J10}]{nirālambo 'ntarā}%
		\rdg[wit={Jyo}]{nirālambāntare}}
		punaḥ/}\\+}
\tl{
\pada{saṃsthitā vyomacakre yā} % cakreṇa G11; °ścaikacakre yā G5; 
\pada{sā mudrā nāma khecarī//\versenr}\\!}
\end{alttlg}
%</vs11-9>
%<*vs11-10>
\begin{alttlg}[hp04_011_10]
\tl{\app{\lem[nolem]{}
	\rdg[wit={C6}]{\NotIn}}%
\pada{\app{\lem[wit={P11,V3}]{sā mayodbheditā vāmā}
		\rdg[wit={G11}]{sā māyodbhedikā vāmā}%  sā mayābhedatā M3, sā madhyodbheditā F
		\rdg[wit={G5}]{sā māyābhedito vāmā}
		\rdg[wit={N19}]{sā mayodve\,\_\,tā vāmā}
		\rdg[wit={V15}]{sā mayodve[dh]itā vāmā}
		\rdg[wit={J10}]{somayodbheditā dhāma}% °rbhidādhāmāmā P6
		\rdg[wit={Jyo}]{somād yatroditā dhārā}}}
\pada{\app{\lem[wit={Zeta2,P11,V3}]{sākṣāc ca}% +G7
		\rdg[wit={Epsi}]{sā sākṣāt}% +M3
		\rdg[wit={J10}]{sākṣād vai}
		\rdg[wit={Jyo}]{sākṣāt sā}} śivavallabhā/}\\+}
\tl{
\pada{\app{\lem[wit={Epsi,Zeta2,P11,V3},alt={pūrayen}]{pūraye\skp{n}}
		\rdg[wit={J10}]{pūjayed}
		\rdg[wit={Jyo}]{pūrayed}
		}%
	\app{\lem[wit={Zeta2,P11,V3},alt={mārutaṃ divyaṃ}]{\skm{n }mārutaṃ divyaṃ}% +G7
		\rdg[wit={G11}]{na tu tad divyaṃ}% + na tu tan divya M3
		\rdg[wit={G5}]{satataṃ divyaṃ}
		\rdg[wit={J10,Jyo}]{atulāṃ divyāṃ}}}
\pada{\app{\lem[wit={Epsi,Zeta2,J10,P11,V3}]{suṣumṇā}
		\rdg[wit={Jyo}]{suṣumṇāṃ}}
	\app{\lem[wit={Epsi,Zeta2,J10,P11,Jyo}]{paścime}%
		\rdg[wit={V3}]{paścimā}}
	mukhe//\versenr}
	\\!}
\end{alttlg}
%</vs11-10>

%<*vs11-11>
\begin{alttlg}[hp04_011_11]
\tl{
\pada{purastāc caiva pūryeta} % °stāvaiva P11, °stācaiva V3, t* caiva N19; °yetā V3
\pada{\app{\lem[wit={Epsi,Zeta2,Pi,Jyo}]{niścitā}% °taṃ F, °cittā G5
		\rdg[wit={J10}]{niśritā}
		} khecarī bhavet/}\\+}
\tl{
\pada{\app{\lem[wit={Epsi,N19,P11,C6},alt={abhyaset}]{abhyase\skp{t}}
		\rdg[wit={V15}]{\om\ (eye-skip?)}
		\rdg[wit={J10,Jyo}]{abhyastā}
		\rdg[wit={V3}]{abhyase}
		}%
	\app{\lem[wit={Epsi,N19,C6,V3},alt={khecarīmudrām}]{\skm{t }khecarīmudrā\skp{m}}
		\rdg[wit={V15}]{\om}
		\rdg[wit={J10,Jyo}]{khecarīmudrā}
		\rdg[wit={P11}]{khecarīṃ mudrām}
		}}%
\pada{\app{\lem[wit={Epsi,N19,J10,Pi},alt={unmanī}]{\skm{m }unmanī}
		\rdg[wit={V15}]{\om}
		\rdg[wit={Jyo}]{py unmanī}
		}
	\app{\lem[wit={Epsi,N19,J10,Jyo}]{saṃprajāyate}
		\rdg[wit={V15}]{\om}
		\rdg[wit={P11}]{sāṃdrajāyate}
		\rdg[wit={C6,V3}]{sā prajāyate}
		}//\versenr}\\!}
\end{alttlg}
%</vs11-11>
%<*vs11-12>
\begin{alttlg}[hp04_011_12]
\tl{\app{\lem[nolem]{}
	\rdg[wit={Jyo}]{\textapp{transposed with the next verse}}}%
\pada{\app{\lem[wit={G11,N19,Pi,Jyo},alt={abhyaset}]{abhyase\skp{t}}
		\rdg[wit={G5}]{abhyasya}
		\rdg[wit={V15}]{abhyasat}
		\rdg[wit={J10}]{abhyaste}}%
	\app{\lem[wit={Epsi,Zeta2,J10,Pi},alt={khecarī}]{\skm{t }khecarī}
		\rdg[wit={Jyo}]{khecarīṃ}}% +J11
	\app{\lem[wit={G11,V15,J10}]{mudrāṃ}
		\rdg[wit={G5,N19}]{mudrā}
		\rdg[wit={Pi,Jyo}]{tāvad}
		}}
\pada{\app{\lem[wit={Epsi,Zeta2,J10},alt={tāvat}]{tāva\skp{t}}
		\rdg[wit={Pi,Jyo},alt={yāvat}]{yāvat}
		}t syā%d
	\app{\lem[wit={G11,Zeta2,C6,Jyo},alt={yoganidritaḥ}]{\skm{d }yoganidritaḥ}
		\rdg[wit={G5}]{coramudritā}
		\rdg[wit={J10}]{yoganidratāḥ}
		\rdg[wit={P11}]{yoganidritāḥ}
		\rdg[wit={V3}]{yoganiṃdrataḥ}
		}/}\\+}
\tl{
\pada{saṃprāptayoganidrasya} % niṃdrasya V3
\pada{kālo nāsti kadācana//\versenr}% canaṃ V3, canā G11
\\!}
\end{alttlg}
%</vs11-12>
%<*vs11-13>
\begin{alttlg}[hp04_011_13]
\tl{\app{\lem[nolem]{}
	\rdg[wit={Jyo}]{\textapp{transposed with the previous verse}}}%
\pada{\app{\lem[nolem]{\skp{pāda a}}
	\rdg[wit={Epsi}]{\om}}% M3 omits too.
bhruvor madhye  % °vo P11, bhrūvaur N19
\app{\lem[wit={Zeta2,J10,C6,V3,Jyo}]{śiva}
	\rdg[wit={P11}]{bhavet}}sthānaṃ}
\pada{\app{\lem[nolem]{\skp{pāda b}}
	\rdg[wit={Epsi}]{\om}}% M3 omits too.
manas tatra vilīyate/}  % vilīyati P11
	\\+}
\tl{
\pada{jñātavyaṃ tat padaṃ turyaṃ} % talpadaṃ G11; tūryaṃ P11,V15
\pada{\app{\lem[wit={Epsi,N19,J10,Pi,Jyo}]{tatra}
		\rdg[wit={V15}]{yatra}}
	\app{\lem[wit={Epsi,V15,J10,Pi,Jyo}]{kālo}
		\rdg[wit={N19}]{kopi}} 
		na vidyate//\versenr}\\!}
\end{alttlg}
%</vs11-13>
%<*vs11-14>
\begin{alttlg}[hp04_011_14]
\tl{\app{\lem[nolem]{}
	\rdg[wit={Jyo}]{\NotIn}}%
\pada{candrasūryadvayor madhye}
\pada{\app{\lem[wit={Epsi,V15,J10,Pi}]{mudrāṃ}
		\rdg[wit={N19}]{mudrā}}
	\app{\lem[wit={G5,V15,J10,Pi}]{dadyāc ca} % +M3; dadyā V3
		\rdg[wit={G11}]{dadyāt tu}
		\rdg[wit={N19}]{divyā ca}}
	\app{\lem[wit={Epsi,V15,J10,C6}]{khecarīm}
		\rdg[wit={N19,V3}]{khecarī}
		\rdg[wit={P11}]{khecare}}/}\\+}
\tl{
\pada{\app{\lem[wit={G11,J10,C6}]{nirālambe}
		\rdg[wit={G5}]{nirālamba}
		\rdg[wit={Zeta2,V3}]{nirālambaṃ}% +F
		\rdg[wit={P11}]{nirālambas}
		}
	\app{\lem[wit={J10,C6}]{mahāśūnye}
		\rdg[wit={G11}]{mahacchūnye}
		\rdg[wit={G5,V3}]{mahāśūnya}
		\rdg[wit={Zeta2}]{mahāśūnyaṃ}% +F
		\rdg[wit={P11}]{tadā śūnya}}}
\pada{vyoma\app{\lem[wit={Epsi,N19,J10,Pi}]{cakre}
		\rdg[wit={V15}]{cakraṃ}}
	\app{\lem[wit={G11,J10,C6,V3}]{vyavasthitām}
		\rdg[wit={G5,P11}]{vyavasthitā}% P6; °taḥ M3?
		\rdg[wit={Zeta2}]{vyavasthitaṃ}% +G7,F
		}//\versenr}\\!} % cf. 4.25*6
\end{alttlg}
%</vs11-14>
%<*vs11-15>
\begin{alttlg}[hp04_011_15]
\tl{
\pada{nirālambaṃ manaḥ kṛtvā} % manaṃ P11
\pada{na kiṃcid api cintayet/}\\+}
\tl{
\pada{sabāhyā\app{\lem[wit={G11,Zeta2,Pi,Jyo},alt={°bhyantare}]{\skp{°}bhyantare}
		\rdg[wit={G5,J10}]{bhyantaraṃ}} vyomni} % vyomani G11, vyoma G5
\pada{\app{\lem[wit={V15,J10,Pi,Jyo},alt={ghaṭavat}]{ghaṭava\skp{t}}
		\rdg[wit={G11}]{aṭavat}
		\rdg[wit={G5}]{maghaṭat}
		\rdg[wit={N19}]{paṭavat}
		}%
	\app{\lem[wit={Epsi,Zeta2},alt={tiṣṭhate}]{\skm{t }tiṣṭhate}
		\rdg[wit={J10,Pi,Jyo}]{tiṣṭhati}} 
		dhruvam//\versenr}\\!}
\end{alttlg}
%</vs11-15>
%<*vs11-16>
\begin{alttlg}[hp04_011_16]
\tl{
\pada{bāhyavāyu%r
	\app{\lem[wit={J10,Pi,Jyo},alt={yathā}]{\skm{r }yathā}% P6
		\rdg[wit={Epsi}]{tathā}% +M3
		\rdg[wit={Zeta2}]{yadā}}
	\app{\lem[wit={G11,V15,C6}]{līnaḥ}
		\rdg[wit={G5}]{līnā}
		\rdg[wit={N19,P11}]{līna}
		\rdg[wit={J10,Jyo}]{līnas}
		\rdg[wit={V3}]{līnaṃ}
		}}
\pada{\app{\lem[wit={Epsi,P11,V3}]{khasya madhye}
		\rdg[wit={N19}]{khamadhya\,\_}
		\rdg[wit={V15}]{khamadhye ca}
		\rdg[wit={J10}]{tathā madhye}
		\rdg[wit={C6}]{khamadhye tu}
		\rdg[wit={Jyo}]{tathā madhyo}
		}
	 \app{\lem[wit={Epsi,V15,J10,Pi,Jyo}]{na saṃśayaḥ}
		\rdg[wit={N19}]{\_\,\_\,sayaḥ}}/}\\+}
\tl{
\pada{\app{\lem[wit={G11,Zeta2,J10,Pi}]{svasthānaṃ gacchati prāṇaḥ} % gacchatiṃ G11; prāṇa N19,V3,F
	\rdg[wit={G5}]{saṃsthānaṃ gacchati prāṇaḥ}
	\rdg[wit={Jyo}]{svasthāne sthiratām eti}}}
\pada{\crux\app{\lem[wit={Epsi,Zeta2,Pi}]{sūryāṅge manasā tathā}% sūryāgnir M1; manatā P6, mānasaṃ G7
		\rdg[wit={J10}]{sūryāṅge pavane tathā}% +P17
		\rdg[wit={Jyo}]{pavano manasā saha}}\crux//\versenr}\\!} % sveyāṅge F
\end{alttlg}
%</vs11-16>

%<*vs11-17>
\begin{alttlg}[hp04_011_17]
\tl{
\pada{evam abhyasyamānasya}
		%\rdg[wit={Jyo}]{abhyasyatas tasya}}}
\pada{\app{\lem[wit={Epsi,J10,Pi,Jyo}]{vāyumārge}% vāyurmāge C6
		\rdg[wit={Zeta2}]{vāyor mārge}} % vāyomārgre N19, mārgaṃ J10
	\app{\lem[wit={C6,Jyo}]{divāniśam}% +M1,F
		\rdg[wit={Epsi,J10}]{sadāniśaṃ}
		\rdg[wit={Zeta2}]{sadānilaṃ}
		\rdg[wit={P11}]{divā niśi}
		\rdg[wit={V3}]{divādisam}
		}/}\\+}
\tl{
\pada{\app{\lem[wit={Epsi,N19,J10,Pi,Jyo}]{abhyāsāj jīryate} % abhyāsā V3; jāyate J10ac?
		\rdg[wit={V15}]{abhyāsāl līyate}}
		vāyu}% % vāyu V3
\pada{r mana\app{\lem[wit={Epsi,Zeta2,J10},alt={tatra vilīyate}]{\skm{s }tatra vilīyate}% P6
		\rdg[wit={Pi,Jyo}]{tatraiva līyate}}//\versenr}\\!}% +F
\end{alttlg}
%</vs11-17>

%<*vs11-18>
\begin{alttlg}[hp04_011_18]
\tl{
\pada{\app{\lem[wit={N19,P11,V3},alt={amṛtaṃ plāvayed deham}]{amṛtaṃ plāvayed deha\skp{m}}
		\rdg[wit={Epsi},postwit=\texteng{(amṛtā \getsiglum{G11})}]{amṛtāt plāvayed deham}% amṛtā G11; +F
		\rdg[wit={V15}]{amṛte plāvayed deham}
		\rdg[wit={J10}]{ajaratvaṃ bhaved dehe}
		\rdg[wit={C6}]{amṛtaṃ plavate \_\,\_}
		\rdg[wit={Jyo}]{amṛtaiḥ plāvayed deham}
		}}%m}
\pada{\app{\lem[wit={ceteri},alt={ā pādatala}]{\skm{m }ā pādatala}
		\rdg[wit={J10}]{apādapala}
		\rdg[wit={C6}]{\lacuna}}%
	\app{\lem[wit={Epsi,V15,Pi,Jyo}]{mastakam} % dehaṃāpādattala°  apādapala J10
		\rdg[wit={N19}]{mastakān}
		\rdg[wit={J10}]{mastake}
		\rdg[wit={C6}]{\lacuna}}/}\\+}
\tl{
\pada{\app{\lem[wit={Epsi,V3,Jyo}]{sidhyaty eva}% °ateva G11
		\rdg[wit={N19}]{siddhaty eva}
		\rdg[wit={V15}]{siddhyaty evaṃ}
		\rdg[wit={J10}]{sidhyate ca}
		\rdg[wit={P11}]{siddhideho}
		\rdg[wit={C6}]{siddhadeho}
		}
	\app{\lem[wit={Epsi,N19,V3}]{sadā kāyo}
		\rdg[wit={V15}]{tadā kāyo}% +F
		\rdg[wit={J10}]{mahāyogo}
		\rdg[wit={P11}]{mahākāryo}
		\rdg[wit={C6,Jyo}]{mahākāyo}
		}}
\pada{mahābalaparākramaḥ//\versenr}\\!} % °kramī P11
\end{alttlg}
%</vs11-18>

% \begin{altpostmula}[hp04_011_18p]
% iti khecarī/ \sgwit{Jyo} % only in the edition
% \end{altpostmula}

%<*vs11-19a>
\begin{altava}[hp04_011_19a]
\app{\lem[nolem]{}
 \rdg[wit={N19,P11}]{\only}}%}
\app{\lem[wit={N19}]{atha}\rdg[wit={P11}]{\om}}
\app{\lem[wit={P11}]{śāmbhavī}
	\rdg[wit={N19}]{śāṃbhavī śaktiḥ}}/ 
\end{altava}
%</vs11-19a>
%<*vs11-19>
\begin{alttlg}[hp04_011_19]
\tl{
\pada{śaktimadhye manaḥ kṛtvā} % mana P11
\pada{\app{\lem[wit={Epsi,N19}]{śaktiṃ ca manamadhyagām}% P6
		\rdg[wit={V15}]{śaktiṃ ca svāṃtamadhyagām}
		\rdg[wit={J10}]{śaktiṃ manasi madhyataḥ}
		\rdg[wit={P11}]{sumadhyagaṃ}
		\rdg[wit={C6,V3}]{manaḥ śaktes tu madhyagam}% +F
		\rdg[wit={Jyo}]{śaktiṃ mānasamadhyagām}
		}/}\\+}
\tl{
\pada{manasā
	\app{\lem[wit={Epsi,V15,J10,P11,C6,Jyo}]{mana ālokya}
		\rdg[wit={N19}]{mana ārokya}
		\rdg[wit={V3}]{manam ālokya}
		}}
\pada{\app{\lem[wit={G11,Zeta2,C6},alt={tad dhyāyet}]{tad dhyāye\skp{t}}
		\rdg[wit={G5}]{taṃ dhyāyet}
		\rdg[wit={J10,Jyo}]{dhārayet}% dhāryate P6
		\rdg[wit={P11}]{taṃ dhātaṃ}
		\rdg[wit={V3}]{vaddhyāyait}
		}t paramaṃ padam//\versenr}\\!}
\end{alttlg}
%</vs11-19>
%<*vs11-20>
\begin{alttlg}[hp04_011_20]
\tl{
\pada{\app{\lem[wit={Epsi,Zeta2,J10,C6,V3,Jyo}]{khamadhye}
		\rdg[wit={P11}]{khaṃmadhye}% +F
		} kuru cātmāna}%m}
\pada{\app{\lem[wit={Epsi,V15,V3,Jyo},alt={ātmamadhye}]{\skm{m }ātmamadhye} % J10pc
		\rdg[wit={N19,J10,P11,C6}]{ātmāmadhye}} % J10ac
		ca khaṃ kuru/}\\+}
\tl{
\pada{\app{\lem[wit={G11,C6,V3}]{ātmānaṃ}% +G7,??
		\rdg[wit={G5,Zeta2,J10,Jyo}]{sarvaṃ ca}% +M1,P6,G10,F
		\rdg[wit={P11}]{evaṃ kṛ}}
	\app{\lem[wit={Zeta2,V3,Jyo}]{khamayaṃ kṛtvā}
		\rdg[wit={Epsi,J10,C6}]{khaṃmayaṃ kṛtvā}% +F
		\rdg[wit={P11}]{tvā tayoś cāpi}}}
\pada{na kiṃcid api cintayet//\versenr}\\!}
\end{alttlg}
%</vs11-20>
%<*vs11-21>
\begin{alttlg}[hp04_011_21]
\tl{\app{\lem[nolem]{}
	\rdg[wit={Zeta2,V3}]{\NotIn}}%
\pada{antaḥśūnyo bahiḥśūnyaḥ} % aṃtaśūnye °śūnyo J10; bahiśūnyaṃ P11
\pada{\app{\lem[wit={Epsi,J10,P11,C6}]{śūnya}
	\rdg[wit={Jyo}]{śūnyaḥ}}kumbha ivāmbare/}\\+} % °baro P11
\tl{
\pada{\app{\lem[nolem]{\skp{pāda c}}
		\rdg[wit={G11}]{\om}}%
	antaḥpūrṇo bahiḥpūrṇaḥ} % bahipūrṇa P11, pūrṇa<<ḥ>> J10
\pada{\app{\lem[nolem]{\skp{pāda d}}
		\rdg[wit={G11}]{\om}}%
	\app{\lem[wit={G5,J10,P11,C6}]{pūrṇa}
		\rdg[wit={Jyo}]{pūrṇaḥ}}kumbha
	\app{\lem[wit={G5,J10,Jyo}]{ivārṇave}% ivādhare P6
		\rdg[wit={P11}]{ivāṃbare}% +N22
		\rdg[wit={C6}]{ivāmbudhau}}//\versenr}
	\\!} % G5 vollständig, M3 om.
\end{alttlg}
%</vs11-21>

%<*vs11-22>
\begin{alttlg}[hp04_011_22]
\tl{\app{\lem[nolem]{}
	\rdg[wit={Zeta2}]{\NotIn}}%
\pada{bāhyacintā na kartavyā} % °tavyo J10
\pada{tathaivāntara%
	\app{\lem[wit={Epsi,J10,Jyo}]{cintanam}% P6/N22
		\rdg[wit={P11}]{ciṃtamān}
		\rdg[wit={C6,V3}]{cintanā}% +F
		}/}\\+}
\tl{
\pada{\app{\lem[wit={Epsi,C6,Jyo}]{sarvacintāṃ parityajya}
	\rdg[wit={J10}]{sarvacintā parityājyā}
	\rdg[wit={P11,V3}]{sarvacintā parityajya}
	}}
\pada{na kiṃcid api cintayet//\versenr}
\\!} % M3 omits. too, but G11,G5 have it.
\end{alttlg}
%</vs11-22>

%<*vs11-23>
\begin{alttlg}[hp04_011_23]
\tl{
\app{\lem[nolem]{\skp{pāda a}}
	\rdg[wit={P11,C6}]{\om}}%
\pada{saṃkalpamātra%
	\app{\lem[wit={Epsi,Zeta2,J10,Jyo}]{kalanaiva}% ka<la>naiva G11
		\rdg[wit={V3}]{kalanaṃ ca}} jaga%t
	\app{\lem[wit={Epsi,Zeta2,V3,Jyo},alt={samagraṃ}]{\skm{t }samagraṃ}
		\rdg[wit={J10}]{samastaṃ}}}\\+}
\tl{
\pada{\app{\lem[nolem]{\skp{pāda b}}
	\rdg[wit={P11,C6}]{\om}}%
saṃkalpamātra% kal<p>a G11
	\app{\lem[wit={Epsi,Zeta2,V3}]{kalanā hi}% karahāṇi P6,N22b
		\rdg[wit={J10,Jyo}]{kalanaiva}} % phalanaiva N22a
	mano\app{\lem[wit={G5,J10,Jyo}]{vilāsaḥ}% +P6,M3
		\rdg[wit={G11}]{vivāsaḥ}
		\rdg[wit={N19}]{vilīnā}
		\rdg[wit={V15}]{valīnā}
		\rdg[wit={V3}]{vilāsā}
		}/}\\+}
\tl{
\pada{\app{\lem[nolem]{\skp{pāda c}}
	\rdg[wit={C6}]{\om}}%
saṃkalpa\app{\lem[wit={G11},alt={°m etam ata}]{\skp{°}m etam ata}
		\rdg[wit={G5}]{mātramanam}
		\rdg[wit={N19}]{mātramata}
		\rdg[wit={V15}]{mātramatam}% P6,N22
		\rdg[wit={J10}]{mātrakalanaiva}
		\rdg[wit={P11}]{mātrami[m]}
		\rdg[wit={V3}]{mātram idam}
		\rdg[wit={Jyo}]{mātramatim}
		}
	\app{\lem[wit={Epsi,V15,P11,V3,Jyo}]{utsṛja}
		%\rdg[wit={Jyo}]{utsṛjya}
		\rdg[wit={N19}]{tsṛja}
		\rdg[wit={J10}]{vikṛtis tu}}
	\app{\lem[wit={Epsi,Zeta2,P11,V3,Jyo}]{nirvikalpaṃ}
		\rdg[wit={J10}]{nityaṃ}}}\\+}
\tl{
\pada{\app{\lem[nolem]{\skp{pāda d}}
	\rdg[wit={C6}]{\om}}%
\app{\lem[wit={Epsi,N19,P11,V3,Jyo}]{āśritya} % āśrītya V3
		\rdg[wit={V15}]{āśrita}
		\rdg[wit={J10}]{saṃkalpa}}
	\app{\lem[wit={G11,J10,Jyo},alt={niścayam}]{niścaya\skp{m}} % P6; = South Ind. 4c
		\rdg[wit={G5,P11}]{niścalam}
		\rdg[wit={Zeta2}]{niścitam}
		\rdg[wit={V3}]{niścalayam}
		}%
	\app{\lem[wit={Epsi,Zeta2,V3,Jyo},alt={avāpnuhi}]{\skm{m }avāpnuhi}
		\rdg[wit={J10}]{avāpnudhi}
		\rdg[wit={P11}]{anāpnuhi}
		}
	\app{\lem[wit={G11,J10,P11,V3,Jyo}]{rāma}% śāmya P6, ?? N22
		\rdg[wit={G5}]{kāma}
		\rdg[wit={N19}]{roga}
		\rdg[wit={V15}]{rāga}
		} śāntim//\versenr}\\!} % śānti P11,V3
\end{alttlg}
%</vs11-23>

%<*vs11-24>
\begin{alttlg}[hp04_011_24]
\tl{\app{\lem[nolem]{}
	\rdg[wit={J10}]{\NotIn}}%
\pada{karpūra%m
	\app{\lem[wit={Epsi,Zeta2,P11,V3,Jyo},alt={anale}]{\skm{m }anale}
		\rdg[wit={C6}]{anile}} yadva}%
\pada{t saindhavaṃ salile yathā/}\\+}
\tl{
\pada{\app{\lem[wit={Epsi,V15,Pi,Jyo}]{tathā}
		\rdg[wit={N19}]{yathā}}
	\app{\lem[wit={Epsi,Pi,Jyo}]{saṃdhīyamānaṃ ca}
		\rdg[wit={Zeta2}]{saṃdīpamānaṃ ca}}}
\pada{mana\app{\lem[wit={G11,V15,C6,Jyo},alt={tattve}]{\skm{s }tattve}
		\rdg[wit={G5,N19}]{tatra}
		\rdg[wit={P11}]{tātva}
		\rdg[wit={V3}]{tatva}
		}
	\app{\lem[wit={Epsi,N19,Pi,Jyo}]{vilīyate}
		\rdg[wit={V15}]{valīyate}}//\versenr}\label{karpura}
	\\!}
\end{alttlg}
%</vs11-24>
%<*vs11-25>
\begin{alttlg}[hp04_011_25]
\tl{
\pada{jñeyaṃ % jñoyaṃ N19
	\app{\lem[wit={Epsi,P11,C6,Jyo}]{sarvaṃ pratītaṃ}
		\rdg[wit={Zeta2,V3}]{sarvapratītaṃ}
		\rdg[wit={J10}]{sarvam atītaṃ}} ca}
\pada{\app{\lem[wit={G11,Zeta2}]{tajjñānaṃ}
		\rdg[wit={G5}]{tat jñātaṃ}
		\rdg[wit={J10,Jyo}]{jñānaṃ ca}% jñāna J10; +F
		\rdg[wit={Pi}]{jñānaṃ tu} % P6
		} mana ucyate/}\\+} % ucyata P11
\tl{
\pada{jñānaṃ % + yaṃ! G11; but jñānaṃ G5,M3
	  jñeyaṃ % m. om. V3, jñoyaṃ N19
	\app{\lem[wit={Epsi,Zeta2,Pi,Jyo}]{samaṃ naṣṭaṃ}
		\rdg[wit={J10}]{manaś caiva}}}
\pada{\app{\lem[wit={Epsi,Zeta2,J10,C6,V3,Jyo}]{nānyaḥ}
		\rdg[wit={P11}]{mānyaḥ}}
	\app{\lem[wit={Epsi,N19,J10,C6,Jyo}]{panthā}
		\rdg[wit={V15}]{paṃtha}
		\rdg[wit={P11}]{paṃthyā}
		\rdg[wit={V3}]{pathā}}
	\app{\lem[wit={Epsi,V15,J10,C6,Jyo}]{dvitīyakaḥ}
		\rdg[wit={N19,P11}]{dvitīyakaṃ}
		\rdg[wit={V3}]{dvitiyaka}}//\versenr}\\!}
\end{alttlg}
%</vs11-25>
%<*vs11-26>
\begin{alttlg}[hp04_011_26]
\tl{
\pada{manodṛśyam idaṃ sarvaṃ} % the last 3 akṣaras damaged G11
\pada{yat kiṃcit sacarācaraṃ/}\\+}
\tl{
\pada{\app{\lem[wit={J10,Jyo}]{manaso hy unmanī}
		\rdg[wit={G11}]{manaso hy amanī}
		\rdg[wit={G5,V15,Pi}]{manasopy unmanī}% +M3
		\rdg[wit={N19}]{mano so 'py unmanī}
		}%
	\app{\lem[wit={V15,J10pc,V3},alt={°bhāve}]{\skp{°}bhāve}% +G7,F
		\rdg[wit={Epsi,C6}]{bhāvo}% +M3
		\rdg[wit={N19}]{\om}
		\rdg[wit={J10ac}]{bhāvavo}
		\rdg[wit={P11}]{bhāvai}
		\rdg[wit={Jyo}]{bhāvād}
		}}
\pada{\app{\lem[wit={V15,P11,C6}]{dvaitābhāvaṃ}% +M3,G7; bhāvaḥ G5
		\rdg[wit={G11}]{dvaitābhā\,+}
		\rdg[wit={G5}]{dvaitābhāvaḥ}
		\rdg[wit={N19}]{bhāvaṃ}
		\rdg[wit={J10,Jyo}]{dvaitaṃ naivo}
		\rdg[wit={V3}]{dvaitābhāva}
		}
	\app{\lem[wit={Epsi,V15,C6,V3}]{pracakṣate}% +Loc-Nom
		\rdg[wit={N19,P11}]{pracakṣyate}% +Nom-Nom
		\rdg[wit={J10,Jyo}]{palabhyate}}//\versenr}\\!}
\end{alttlg}
%</vs11-26>

%<*vs11-27>
\begin{alttlg}[hp04_011_27]
\tl{
\pada{jñeyavastuparityāgā}% % jñeyaṃ?  parī° N19; jñeyaṃ yas tu P11
\pada{d vilayaṃ yāti % yāṃti P11, damaged G11
	\app{\lem[wit={Epsi,V15,J10,Pi,Jyo}]{mānasam}
		\rdg[wit={N19}]{mārutaṃ}}/}\\+}
\tl{
\pada{\app{\lem[wit={G11,Zeta2,Pi}]{mānase}
		\rdg[wit={G5,J10,Jyo}]{manaso}}
	\app{\lem[wit={Epsi,Zeta2,J10,P11,V3}]{vilayaṃ}
		\rdg[wit={C6,Jyo}]{vilaye}}
	\app{\lem[wit={G11,Zeta2,P11}]{yāte}
		\rdg[wit={G5}]{yāti}
		\rdg[wit={J10,C6,V3,Jyo}]{jāte}}}
\pada{kaivalya%m
	\app{\lem[wit={Epsi,V15,Pi,Jyo},alt={avaśiṣyate}]{\skm{m }avaśiṣyate}
		\rdg[wit={N19}]{anasīṣyate}
		\rdg[wit={J10}]{api kalpate}}//\versenr}\\!}
\end{alttlg}
%</vs11-27>

%<*vs11-28>
\begin{alttlg}[hp04_011_28]
\tl{%\app{\lem[nolem]{}
	%\rdg[wit={Jyo},alt=\textapp{found between \ref{yatradrsti} and \ref{vedasastra}}]{\skp{found between xxx and xxx}}}%
\pada{layo laya iti prāhuḥ} % first 2 akṣaras damaged G11; layaṃ? C6; prāhur N19,V15
\pada{\app{\lem[wit={Epsi,J10,Pi,Jyo}]{kīdṛśaṃ}
		\rdg[wit={Zeta2}]{īdṛśaṃ}} layalakṣaṇam/}\\+}
\tl{
\pada{apunarvāsano% apu<na>r° N19
	\app{\lem[wit={N19,J10,P11,C6,Jyo},alt={°tthānāl}]{\skp{°}tthānā\skp{l}}% °nāl P11,Jyo, °tthā .. l C6, nāt* N19, °nād J10
		\rdg[wit={G11,V15,V3}]{tthānā}
		\rdg[wit={G5}]{tthāna}}}% % +P7
\pada{\app{\lem[wit={G5,Zeta2,Pi,Jyo},alt={layo viṣaya}]{\skm{l }layo viṣaya}
		\rdg[wit={G11}]{yalo viṣaya}%
		\rdg[wit={J10}]{vṛttyayā viśva}
		}vismṛtiḥ//\versenr}% °smṛti N19,V3; °smati P11, niśrutiḥ G5
	\label{layo}
	\\!}
\end{alttlg}
%</vs11-28>
%<*vs11-29>
\begin{alttlg}[hp04_011_29]
\tl{\app{\lem[nolem]{}
 \rdg[wit={Epsi,Zeta2,J10,Pi,Jyo}]{\incl}}%
\pada{evaṃ nānāvidhopāyāḥ} % eva N19; ḥ om. P11,V3
\pada{samyaksvānu%
	\app{\lem[wit={G11,N19,J10,Pi,Jyo}]{bhavānvitāḥ} % ḥ om. V3; sānubhavanvitā P11
		\rdg[wit={G5}]{bhavānyuta}
		\rdg[wit={V15}]{bhavātmikāḥ}}/}\\+}
\tl{
\pada{samādhi\app{\lem[wit={Epsi,Zeta2,P11,C6,Jyo}]{mārgāḥ}% ḥ om. P11
		\rdg[wit={J10}]{mārge}
		\rdg[wit={V3}]{\illeg}} 
		kathitāḥ} % °tā V15
\pada{pūrvācāryair mahātmabhiḥ//\versenr}\\!}% °ryai P11,V3
%\myfn{After this verse, \getsiglum{P11,C6} have \devnote{iti viśrāntiḥ} and \getsiglum{N19,V15} \devnote{atha viśrāntiḥ}.}
\end{alttlg}
%</vs11-29>

%<*vs11-30a>
\begin{altava}[hp04_011_30a]
\app{\lem[nolem]{}
	\rdg[wit={J10,V3,Jyo}]{\NotIn}}%
	%\rdg[wit={Epsi,Zeta2,P11,C6}]{\incl}}%}
\app{\lem[wit={Epsi,Zeta2}]{atha}
	\rdg[wit={P11,C6}]{iti}% °ti P11; +F
	} viśrāntiḥ/ 
\end{altava}
%</vs11-30a>

%<*vs11-30>
\begin{alttlg}[hp04_011_30]
\tl{\app{\lem[nolem]{}
	\rdg[wit={J10}]{\NotIn}}%
\pada{\app{\lem[wit={Epsi,V15,Pi,Jyo}]{suṣumṇāyai}
		\rdg[wit={N19}]{sukhayaiḥ}} kuṇḍalinyai}
\pada{sudhāyai candra% ceṃdra V3
	\app{\lem[wit={Epsi}]{maṇḍale}% +M3
		\rdg[wit={Zeta2}]{maṇḍalāt}% +G7,G3
		\rdg[wit={Pi,Jyo}]{janmane}% +M1
		}/}\\+}
\tl{
\pada{manonmanyai namas tubhyaṃ} % °nye P11,G5
\pada{mahā\app{\lem[wit={Epsi,Zeta2,P11,C6}]{śakti}
		\rdg[wit={V3}]{śakte}
		\rdg[wit={Jyo}]{śaktyai}% +F
		}%
	\app{\lem[wit={Zeta2,C6,V3,Jyo}]{cidātmane}
		\rdg[wit={G11}]{cidātmike}% °ātmake M3
		\rdg[wit={G5}]{cidātmine}
		\rdg[wit={P11}]{cidātmani}% +G7 
		}//\versenr}
	\\!}
\end{alttlg}
%</vs11-30>

%<*vs11-31>
\begin{alttlg}[hp04_011_31]
\tl{\app{\lem[alt=\mylem{71ab},nosep,nonum]{\skp{pāda ab}}
	\rdg{\texteng{cf. \manuref{X4.122ab (4.35ab)}}}}%
\pada{\app{\lem[wit={G5,Zeta2,P11,Jyo}]{aśakya}% +G5
		\rdg[wit={G11,J10}]{aśakyaṃ}
		\rdg[wit={C6,V3}]{aśakta}}tattvabodhānāṃ} % ṃ om. G11
\pada{\app{\lem[wit={Epsi,Zeta2,J10,C6,V3,Jyo},alt={mūḍhānām}]{mūḍhānā\skp{m}}
		\rdg[wit={P11}]{gūḍhānām}}%
	\app{\lem[wit={Epsi,J10,Pi,Jyo},alt={api saṃmatam}]{\skm{m }api saṃmatam}
		\rdg[wit={N19}]{atisaṃtataṃ}
		\rdg[wit={V15}]{api saṃtataṃ}
		}/} 
	\\+}
\tl{
\pada{proktaṃ 
	\app{\lem[wit={Zeta2,J10,Pi,Jyo}]{gorakṣa}
		\rdg[wit={Epsi}]{śrīśaṃbhu}}nāthena}
\pada{nādopāsana%m % nāptepā° P11
	\app{\lem[wit={Epsi,Zeta2,J10,P11,V3,Jyo},alt={ucyate}]{\skm{m }ucyate}
		\rdg[wit={C6}]{uttamam}}//\versenr}\\!}
\end{alttlg}
%</vs11-31>
\endaltrecension
\fi


%<*vs12>
\begin{tlg}[hp04_012]
\tl{
\pada{\app{\lem[wit={ceteri}]{śrīādināthena}% nāthona N23
	\rdg[wit={Epsi}]{śrīśaṃbhunāthena}} 
	sapādakoṭi}-\\+}% koṭī G11
\tl{
\pada{\app{\lem[wit={ceteri}]{laya}
		\rdg[wit={N3,Gamma,N19}]{layaḥ}
		\rdg[wit={J5}]{laṣa}
		}prakārāḥ % prakār<<āḥ>> J10
		kathitā  % kathaṃ J10
	\app{\lem[wit={N3,J5,Epsi,N19}]{jayante}% +F
		\rdg[wit={Gamma,E2,V15,J10,Pi,Jyo}]{jayanti} % jayati V17
		\rdg[wit={V19}]{yayaṃti}}/}\\+}
\tl{
\pada{nādānusandhānaka%m % nodānasaṃdhānasaṃdhānakam N19
	\app{\lem[wit={N3,Epsi,P11,C6,Jyo},alt={ekam eva}]{\skm{m }ekam eva}% +F
		\rdg[wit={J5,V3}]{eva}
		\rdg[wit={Gamma,Delta}]{eva kāryaṃ}
		\rdg[wit={N19,J10}]{eva nānyaṃ}
		\rdg[wit={V15}]{eva mānyaṃ}
		}}\\+} % °karmeva kārya N23
\tl{
\pada{\app{\lem[wit={ceteri}]{manyāmahe} % manyāṃmahe N19
	\rdg[wit={C6}]{gaṇyāmahe}}
	\app{\lem[wit={N3,Zeta2,P11,V3}]{mānyatamaṃ}% myanya° P11
		\rdg[wit={J5,Gamma,Delta,Epsi}]{nānyatamaṃ}% +F
		\rdg[wit={J10}]{tātarasaṃ}
		\rdg[wit={C6}]{nānyamataṃ}
		\rdg[wit={Jyo}]{mukhyatamaṃ}} 
		layānām//\versenr}\label{sapadakoti}
		\orgvnr{12}\\!} % layāṇāṃ J10
\end{tlg}
%</vs12>
\commcite\newpage


%<*vs13>
\begin{tlg}[hp04_013]
\tl{%\app{\lem[nolem]{}
	%\rdg[wit={Jyo},alt=\textapp{found betw. X4.72 and X4.73}]{\skp{found betw. X4.72 and X4.73}}}%
\pada{\app{\lem[nolem]{\skp{pāda a}}
	\rdg[wit={N23a,J7a}]{\om}}%
\app{\lem[wit={ceteri}]{muktāsanasthito}% +K3,C7
	\rdg[wit={N23b}]{mudrāsanasthite}
	\rdg[wit={V19a,Jyo}]{muktāsane sthito} % or: °sāna? V19
	} yogī}
\pada{\app{\lem[nolem]{\skp{pāda b}}
	\rdg[wit={N23a,J7a}]{\om}}%
mudrāṃ saṃdhāya śāmbhavīm/}% sadhāya? V19
% mudrā P11,N23b,V3,J10; śāṃbhavī Gammab,P11,V3, sāṃbhavī J10
	\\+}
\tl{
\pada{śṛṇuyād dakṣiṇe karṇe} % śṛṇuyā J7a,V3, śuṇu° V19a, śruṇu° V19b,V15
\pada{\app{\lem[wit={ceteri},alt={nādam}]{nāda\skp{m}}
	\rdg[wit={C6}]{\_\,\_}}%
\app{\lem[resp=emend,alt={antaḥstham ekadhīḥ}]{\skm{m }antaḥstham ekadhīḥ}% 
	\rdg[wit={N3,G4,N23a,J7a,P11,Jyo}]{antastham ekadhīḥ}% aṃtta G4, ekadhī P11
	\rdg[wit={J5}]{atastham ekadhā}
	\rdg[wit={N23b,J7b,V19b,E2b,Epsi,Zeta2}]{antargataṃ sadā}
	\rdg[wit={V19a}]{ekāntake sudhīḥ}
	\rdg[wit={E2a}]{ekāntike sudhīḥ}% +K3,C7
	\rdg[wit={J10,V3}]{antargataṃ mahat}
	\rdg[wit={C6}]{nādamataṃ sadā}
	}//\versenr}\myfn{%
	This verse is found twice in \getsiglum{Gamma,Delta}: %
	first (a) after 4.12/X4.72, and second (b) after 4.36/X4.84.}
	\orgvnr{13}\\!}
\end{tlg}
%</vs13>
\commcite%\newpage



%<*vs14>
\begin{tlg}[hp04_014]
\tl{
\pada{\app{\lem[wit={ceteri}]{kāṣṭhe}
		\rdg[wit={N23}]{kaṣṭaiḥ}
		\rdg[wit={J7,Delta,C6}]{kāṣṭhaiḥ}
		}
	\app{\lem[wit={ceteri}]{pravartito}% pravartte J5
		\rdg[wit={V15,J10}]{pravartate}
		} vahniḥ} % vahni J5,P11,V3
\pada{\app{\lem[wit={ceteri}]{kāṣṭhena}
		\rdg[wit={N23}]{kaṣṭena}}
	\app{\lem[wit={ceteri}]{saha}
		\rdg[wit={V15}]{sa}}
	\app{\lem[wit={ceteri}]{śāmyati}
		\rdg[wit={N3,J5,V19,V3}]{sāmyati}
		\rdg[wit={V15}]{līyate}}/}\\+}
\tl{
\pada{\app{\lem[wit={ceteri}]{nāde}% nādena J5
		\rdg[wit={N23}]{nā}}
	\app{\lem[wit={ceteri}]{pravartitaṃ} %°taṃś N19
		\rdg[wit={V15}]{pravartite}
		\rdg[wit={J10}]{pravartate}}
	\app{\lem[wit={ceteri}]{cittaṃ}
		\rdg[wit={N23}]{\om}}}
\pada{nādena saha līyate//\versenr}
\orgvnr{14}\\!}
\end{tlg}
%</vs14>
\commcite\newpage

%<*vs15>
\begin{tlg}[hp04_015]
\tl{\app{\lem[nolem]{}
	\rdg[wit={J10,Jyo}]{\om}}%
\pada{\app{\lem[wit={ceteri}]{vismṛtya}
%		\rdg[wit={C7}]{nismṛtya}
		\rdg[wit={E2}]{niḥsṛtya}
		\rdg[wit={G5}]{visṛjya}
		} sakalaṃ bāhyaṃ} % bāhya N19
\pada{\app{\lem[wit={N3,J5,J7,Delta,G5,V15,Pi}]{nāde}
		\rdg[wit={N23}]{na\,\_}
		\rdg[wit={G11}]{nādo}
		\rdg[wit={N19}]{nāda}
		}
	\app{\lem[wit={ceteri}]{dugdhāmbu}
		\rdg[wit={N23}]{gugyāṃbu}}va%n
	\app{\lem[wit={ceteri},alt={manaḥ}]{\skm{n }manaḥ}
		\rdg[wit={N23,Delta}]{naraḥ}
		\rdg[wit={V3}]{mana}
		}/}\\+}
\tl{
\pada{\app{\lem[wit={G4,Gamma,E2,Epsi,Zeta2,C6}]{ekībhūyātha}% °ācca G5
		\rdg[wit={N3}]{°bhūtvātha}
		\rdg[wit={J5}]{°bhūyotha}
		\rdg[wit={V19}]{°bhūyāya}
		\rdg[wit={P11}]{°bhūyādya}
		\rdg[wit={V3}]{°bhūyā}
		}
	\app{\lem[wit={ceteri}]{sahasā}
		\rdg[wit={J5}]{manasā}
		\rdg[wit={V3}]{sahasā ca}
		}}
	% ekībhūyād atha saha cidā(page break)ekībhūyātha sahasā J7
\pada{\app{\lem[wit={cetwG4}]{cidākāśe}
		\rdg[wit={J5}]{cidāśe}
		\rdg[wit={N23}]{vidāktośe}
		\rdg[wit={J7}]{cidākaro}} 
\app{\lem[wit={ceteri}]{vilīyate}
	\rdg[wit={N3}]{valīyate}
	\rdg[wit={G4}]{na lipyate}}//\versenr}
	\orgvnr{15}\\!}
\end{tlg}
%</vs15>
\commcite%\newpage

%<*vs16>
\begin{tlg}[hp04_016]
\tl{\app{\lem[nolem]{}
	\rdg[wit={Jyo}]{\om}}%
\pada{\app{\lem[wit={Delta,Epsi,J10,P11}]{audāsīnya}
		\rdg[wit={N3}]{udāsonya}
		\rdg[wit={J5}]{udāsinya}
		\rdg[wit={G4}]{audāśinya}
		\rdg[wit={N23}]{odāsīnya}
		\rdg[wit={J7,V3}]{udāsīnya}
		\rdg[wit={N19}]{ṛdāsīnya}
		\rdg[wit={V15}]{audāsinya}
		\rdg[wit={C6}]{audāsīna}
		}paro bhūtvā}
\pada{sadābhyāsena saṃyamī/}\\+} % °bhyosena V3
\tl{
\pada{unmanī\app{\lem[wit={N3,Gamma,Delta,P11,C6}]{karaṇaṃ}
		\rdg[wit={J5}]{karaṇe}
		\rdg[wit={Epsi,Zeta2,J10}]{kārakaṃ}
		\rdg[wit={V3}]{karaṇa}
		} sadyo} % sadyā N3, vastu G5
\pada{\app{\lem[wit={ceteri},alt={nādam}]{nāda\skp{m}}
		\rdg[wit={N19}]{bhāda}}%
	\app{\lem[wit={ceteri},alt={evāvadhārayet}]{\skm{m }evāvadhārayet}
		\rdg[wit={J5}]{evāvadhārayan}
		\rdg[wit={V15}]{eva sadābhyaset}}//\versenr}
	\orgvnr{16}\\!}
\end{tlg}
%</vs16>
\commcite\newpage

%<*vs17a>
\begin{ava}[hp04_017a]
\app{\lem[nolem]{}
	\rdg[wit={G5,Jyo}]{\om}}%
\app{\lem[wit={N3,N23,G11,P11},alt={kīdṛśam},post=\texteng{(ki° \getsiglum{N3})}]{kīdṛśa\skp{m}} % ki° N3
%		\rdg[wit={C7}]{kīdṛṣam}
		\rdg[wit={J5,J7}]{kīdṛśīm}
		\rdg[wit={V19}]{kim}
		\rdg[wit={E2,V15}]{\om}
		\rdg[wit={N19,J10}]{idṛśam}% +F
		\rdg[wit={C6,V3}]{kīdṛśyam}% +K3
		}%
	\app{\lem[wit={Gamma,V19,G11,J10,Pi},alt={audāsīnyam}]{\skm{m }audāsīnyam} % °śīnyaṃ V3
		\rdg[wit={N3}]{audasīnyaṃ}
		\rdg[wit={J5}]{audāsinyā}
		\rdg[wit={E2}]{athaudāsīnyam}
		\rdg[wit={Zeta2}]{audāsinyaṃ}
		}/
\end{ava}
%</vs17a>
%<*vs17>
\begin{tlg}[hp04_017]
\tl{\app{\lem[nolem]{}
	\rdg[wit={Jyo}]{\om}}%
\pada{\app{\lem[wit={ceteri}]{śīte} % sīte N19,V3
		\rdg[wit={J5}]{śīta}
		\rdg[wit={V15}]{śīti}
		\rdg[wit={J10}]{jñāte}
		}
	\app{\lem[wit={ceteri}]{kāle}
		\rdg[wit={N3}]{\om}
		\rdg[wit={J5}]{rakṣa}
		\rdg[wit={J7}]{kāla}
		\rdg[wit={J10}]{kā}
		}
	\app{\lem[wit={G5,J10,V3}]{caupaṭī vā kuṭī vā}% +J6,G7,M3
		\rdg[wit={N3}]{caupaṭī vā paṭī vā}
		\rdg[wit={J5}]{ṇe kathā vā paṭī vā}
		\rdg[wit={N23}]{cāpaṭī cāpaṭī vā}% +C7
		\rdg[wit={J7,E2}]{cāpaṭī vā paṭī vā}% cāpaṭe E2ac
		\rdg[wit={V19}]{cāpaṭī vā paṭīkā}
		\rdg[wit={G11}]{dvaupaṭī vā kuṭī vā}% dvau paṭau vā kuṭī vā F
		\rdg[wit={N19}]{copaṭī vā paṭī vā}
		\rdg[wit={V15}]{paṭī vā}
		\rdg[wit={P11}]{copaṭī vā kuṭī vā}
		\rdg[wit={C6}]{cāpaṭī vā kuṭī vā}
		}}\\+}
\tl{
\pada{\app{\lem[wit={N3,J5,E2,Epsi,N19,P11,V3}]{pathyāhāre}% payo° G5
		\rdg[wit={N23}]{yathāhārā}
		\rdg[wit={J7,V15,J10,C6}]{pathyāhāro}% +K3,C7
		\rdg[wit={V19}]{<<mi>>thyāhāro}
		}
	\app{\lem[wit={ceteri}]{gopayo}% +K3
		\rdg[wit={V19}]{gopatho}
%		\rdg[wit={E2,C7}]{gomayo}% +F
		}
	\app{\lem[wit={ceteri}]{vā}
		\rdg[wit={N23}]{\om}
		\rdg[wit={J10}]{co}
		}
	\app{\lem[wit={ceteri}]{payo vā}
		\rdg[wit={N23}]{<<payo>> vā}
		\rdg[wit={V19}]{patho vā}
		\rdg[wit={C6}]{tha pānaṃ}}/}\\+}
\tl{
\pada{\app{\lem[wit={Alpha,Epsi,P11,V3}]{bhojye}% G11pc
		\rdg[wit={Gamma}]{bhakṣe}
		\rdg[wit={V19,C6}]{bhakṣyaṃ}
		\rdg[wit={E2}]{bhikṣye}
		\rdg[wit={N19}]{bhojya}
		\rdg[wit={V15,J10}]{bhojyaṃ}% +M3
		}
	\app{\lem[wit={ceteri}]{bhikṣā} % bhīkṣā V15, bhikṣyā N3,P11
		\rdg[wit={J10}]{bhuktaṃ}}%
	\app{\lem[wit={ceteri},alt={vṛndam}]{vṛnda\skp{m}}% vṛndapā° E2
		\rdg[wit={Epsi}]{kandam}% +M3
		\rdg[wit={J10}]{cānnam}
		\rdg[wit={P11}]{mṛdam}
		}%
	\app{\lem[wit={Alpha,J7,Delta,V15},alt={āraṇyakandaṃ}]{\skm{m }āraṇyakandaṃ}% +F
		\rdg[wit={N23}]{āramyakaṃdaṃ}
		\rdg[wit={Epsi}]{āraṇyakaṃ vā}% +N22,N12,M3
		\rdg[wit={N19,J10,V3},alt={°kaṃda}]{āraṇyakaṃda}
		\rdg[wit={P11},alt={°kaṃdā}]{āraṇyakaṃdā}
		\rdg[wit={C6}]{āpaṇyakaṃ vā}}}\\+}
\tl{
\pada{\app{\lem[wit={N3,J7,Delta,Epsi,P11}]{pāṇī droṇī}% vāṇī G5
		\rdg[wit={J5,V15,J10}]{pāṇi droṇī}
		\rdg[wit={G4}]{pāṇi droṇi}
		\rdg[wit={N23}]{pāṇīndrāṇī}
		\rdg[wit={N19}]{pāṇī drāṇi}
		\rdg[wit={C6}]{pāṇiṃ droṇe}
		\rdg[wit={V3}]{pāṇi}
		}
	\app{\lem[wit={N3,G4,G11,Zeta2,P11}]{kāpi vā}
		\rdg[wit={J5}]{vā kapī}
		\rdg[wit={N23}]{khapaḍā}
		\rdg[wit={J7}]{kāpaṭo}
		\rdg[wit={V19}]{kharparo}
		\rdg[wit={E2}]{karparā}% +K3,C7
		\rdg[wit={G5}]{cātha vā}
		\rdg[wit={J10}]{kāthivā}
		\rdg[wit={C6}]{karpaṭaṃ}
		\rdg[wit={V3}]{kāpivāṃ}
		}
	\app{\lem[wit={J5,G4,Epsi,N19,P11}]{bhojyapātre}% patre P11
		\rdg[wit={N3,Delta,V15,J10,V3}]{bhojyapātraṃ}% +F
		\rdg[wit={N23}]{bhājapatraṃ}
		\rdg[wit={J7}]{bhūrjapatraṃ}
		\rdg[wit={C6}]{bhojapatraṃ}
		}//\versenr}
	\orgvnr{17}\\!}
\end{tlg}
%</vs17>

\avacite{17a}
\commcite\newpage


%<*vs18>
\begin{tlg}[hp04_018]
\tl{\app{\lem[nolem]{}
	\rdg[wit={Jyo}]{\om}}%
\pada{\app{\lem[wit={J7,Delta,Epsi,N19}]{sarvacintāṃ}
		\rdg[wit={N3,J5,V15,J10,Pi}]{sarvacintā}
		\rdg[wit={N23}]{\om}}
	\app{\lem[wit={J5,Zeta2,J10,P11,V3}]{samutsṛjya}% +M3
		\rdg[wit={N3}]{samutyajya}
		\rdg[wit={G11}]{samṛtsṛjya}
		\rdg[wit={N23}]{\om}
		\rdg[wit={J7,Delta,G5,C6}]{parityajya}
		}}
\pada{sarva\app{\lem[wit={N3,Epsi,V15,Pi}]{ceṣṭāṃ}
		\rdg[wit={J5}]{ceṣṭā}
		\rdg[wit={Gamma,Delta}]{kāle}
		\rdg[wit={N19}]{ceṣṭī}
		\rdg[wit={J10}]{ceṣṭāś}
		} ca sarvadā/}\\+}
\tl{
\pada{nādam % -m- is a hiatus bridge.
	evānu\app{\lem[wit={N3,P11,C6},alt={°saṃdhānān}]{\skp{°}saṃdhānā\skp{n}}
		\rdg[wit={J5,Epsi,Zeta2,J10}]{saṃdadhyān} % °dhyā N19, saṃ<<da>>dhyān J10
		\rdg[wit={Gamma,Delta}]{saṃdhatte}% dhartte N23
		\rdg[wit={V3}]{saṃdhānā}
		}}%n 
\pada{\app{\lem[wit={ceteri},alt={nāde}]{\skm{n }nāde}
		\rdg[wit={C6}]{devi}
		} cittaṃ vilīyate//\versenr} % vittaṃ vileyate P11
	\orgvnr{18}\\!}
\end{tlg}
%</vs18>
\commcite\newpage



%<*vs19>
\begin{tlg}[hp04_019]
\tl{
\pada{ārambha%ś  % araṃbha° P7
	\app{\lem[wit={ceteri},alt={ca}]{\skm{ś }ca}
		\rdg[wit={V19}]{ca\,\_}}
	\app{\lem[wit={ceteri},alt={ghaṭaś}]{ghaṭa\skp{ś}}
		\rdg[wit={N23}]{gha\,\_\,ś}}%
	\app{\lem[wit={ceteri},alt={caiva}]{\skm{ś }caiva}
		\rdg[wit={V19}]{ca}
		\rdg[wit={J10}]{caivas}
		}}
\pada{tathā
	\app{\lem[wit={N3,G4,Epsi,N19,J10,Pi},alt={paricayas}]{paricaya\skp{s}}
		\rdg[wit={J5,N23,Delta,Jyo}]{paricayo}
		\rdg[wit={J7}]{pariyo}
		\rdg[wit={V15}]{paricas}
		}%
	\app{\lem[wit={N3,V15,V3},alt={tathā}]{\skm{s }tathā}% =ŚS; trayaṃ N24; +F
		\rdg[wit={J5,Gamma,E2,Jyo}]{'pi ca}
		\rdg[wit={G4,G11,N19,J10,P11,C6}]{tataḥ}% +M3
		\rdg[wit={G5}]{smṛtaḥ}
		\rdg[wit={V19}]{pi vā}
		}/}\\+}
\tl{
\pada{niṣpattiḥ % niṣpaṃti N23, niṣyandi G5
	\app{\lem[wit={ceteri}]{sarvayogeṣu}% yogyeṣu G5
		\rdg[wit={E2}]{sarvayoge ca}
		\rdg[wit={Pi}]{ceti yogeṣu}
		}}
\pada{\app{\lem[wit={N3,G4}]{yogāvasthā bhavanti tāḥ}
		\rdg[wit={J5}]{yogāvasthā bhavanti te}
		\rdg[wit={Gamma,Delta}]{yogāvasthā prakīrtitā}
		\rdg[wit={Epsi,Zeta2,J10,Pi,Jyo}]{syād avasthācatuṣṭayaṃ}% +F
		}\marma//\versenr}
		\orgvnr{19}\\!}
\end{tlg}
%</vs19>
\commcite\newpage


%<*vs20a>
\begin{ava}[hp04_020a]
\app{\lem[nolem]{}
	\rdg[wit={N3,J5,Pi}]{\om}}%
\app{\lem[resp=emend]{tatrārambhāvasthā}
		\rdg[wit={G4,G5,Zeta2}]{tatra ārambhaḥ}
		\rdg[wit={N23,Jyo}]{athārambhāvasthā}% atha a° Jyo
		\rdg[wit={J7}]{ārambhāvasthātha}
		\rdg[wit={V19}]{athārambharakṣā}% +K3,C7
		\rdg[wit={E2}]{athārambhadīkṣā}
		\rdg[wit={G11}]{tatrārambhaḥ}
		\rdg[wit={J10}]{tatra cārambhaḥ}
		}/
\end{ava}
%</vs20a>
%<*vs20>
\begin{tlg}[hp04_020]
\tl{
\pada{brahma\app{\lem[wit={N3,G5,Jyo},alt={granther}]{granthe\skp{r}}% +J11pc,M3
		\rdg[wit={J5}]{granthid}
		\rdg[wit={N23,C6}]{granthi}
		\rdg[wit={J7,V19,V15,V3}]{granthir}% grathir J7; +K3,C7
		\rdg[wit={E2}]{granthau}
		\rdg[wit={G11}]{gra\,+}
		\rdg[wit={N19}]{raṃdhre}
		\rdg[wit={J10}]{granthiṃ}
		\rdg[wit={P11}]{granthe}
		}r bhave%d % bhavod C6
	\app{\lem[wit={N3,Epsi,C6,V3},alt={bhedād}]{\skm{d }bhedā\skp{d}}
		\rdg[wit={J5,P11}]{bhedā}
		\rdg[wit={Gamma,V19}]{bhinna}
		\rdg[wit={E2}]{bhinne}
		\rdg[wit={N19}]{bhed}
		\rdg[wit={V15}]{bhinnād}
		\rdg[wit={J10}]{bhinnā}
		\rdg[wit={Jyo}]{bhedo hy}
		}}%
\pada{\app{\lem[wit={ceteri},alt={ānandaḥ}]{\skm{d }ānandaḥ}
		\rdg[wit={J5,N23,C6}]{ānaṃda}
		\rdg[wit={J10}]{nādaḥ}
		\rdg[wit={P11}]{nanādaḥ}
		}
	śūnya\app{\lem[wit={ceteri}]{saṃbhavaḥ} % śūnyaṃ J10; °bhava N19, °bhavaṃ P7, °bhavā J5
		\rdg[wit={J10}]{samaṃbhavaḥ}}/}\\+}
\tl{
\pada{vicitra\app{\lem[wit={E2,Epsi}]{kvaṇako}
		\rdg[wit={N3}]{kvana˟ko}
		\rdg[wit={J5}]{kanako}
		\rdg[wit={Gamma}]{s tatkṣaṇād}% ta<<t>> N23
		\rdg[wit={V19}]{kṣike}
		%\rdg[wit={K3,C7}]{kṣaṇike}
		\rdg[wit={N19,V3}]{kaṇako}
		\rdg[wit={V15}]{kvaṇiko}
		\rdg[wit={J10}]{kuṇako}
		\rdg[wit={P11}]{ṣkāṇako}
		\rdg[wit={C6}]{kuṇape}
		\rdg[wit={Jyo}]{ḥ kvaṇako}
		} 
	\app{\lem[wit={ceteri}]{dehe}
		\rdg[wit={J5}]{deho}
		\rdg[wit={C6}]{caivā}
		}}% caiva C6ac
\pada{\app{\lem[wit={N3,J5,Epsi,Zeta2,J10,Pi,Jyo}]{'nāhataḥ śrūyate}% °hata P11, °hate G5
		\rdg[wit={Gamma}]{sarvataḥ śrūyate}
		\rdg[wit={Delta}]{śrūyate (')nāhata}
		}
		dhvaniḥ//\versenr}
		\orgvnr{20}\\!} % dhvani J5,P11,V3,J7, ddhavaniḥ N19
\end{tlg}
%</vs20>

\avacite{20a}
\commcite\newpage

%<*vs21>
\begin{tlg}[hp04_021]
\tl{
\pada{%\app{\lem[nolem]{\skp{pāda a}}
%	\rdg[wit={Delta,V3}]{\om}}%
\app{\lem[wit={N3,J5,Gamma,G5,P11,C6,Jyo}]{divyadehaś ca tejasvī} % tejasvā N3; +M3
		\rdg[wit={Delta,V3}]{\om}
		\rdg[wit={G11}]{divyadehasya tejasvī}
		\rdg[wit={N19}]{ādityatejaś ca tejasvī}
		\rdg[wit={V15}]{tejasvī divyagandhaś ca}
		\rdg[wit={J10}]{divyagandho divyacakṣuś ca}
		}}
\pada{%\app{\lem[nolem]{\skp{pāda b}}
%	\rdg[wit={Delta,V3}]{\om}}%
\app{\lem[wit={N3,G4,Gamma,P11,C6,Jyo}]{divyagandhas tv arogavān} % gaṃdhās N3
		\rdg[wit={J5}]{divyadeham arogavān}
		\rdg[wit={Delta,V3}]{\om}
		\rdg[wit={Epsi,N19}]{divyagandho py arogavān}% paro° N19
		\rdg[wit={V15}]{divyadeho py arogavān}
		\rdg[wit={J10}]{tejasvī ārogavān}
		}/}
		\\+}
\tl{
\pada{\app{\lem[wit={ceteri}]{saṃpūrṇa}
		\rdg[wit={V15}]{saṃpūrṇe}}%
	\app{\lem[wit={Alpha,G5,N19,P11,Jyo}]{hṛdayaḥ}
		\rdg[wit={N23,V19,V15,J10,C6,V3}]{hṛdaye}% +J11pc
		\rdg[wit={J7,G11}]{hṛdaya}
		%\rdg[wit={E2,C7}]{nilaye}
		}
	\app{\lem[wit={Alpha,Zeta2},alt={śūnye tv}]{śūnye\skp{ tv}}
		\rdg[wit={Gamma,Delta,G11,J10,C6}]{śūnye}
		\rdg[wit={P11}]{śūra}
		\rdg[wit={G5,V3,Jyo}]{śūnya}% +M3,F
		}}
\pada{\app{\lem[wit={ceteri},alt={ārambhe}]{\skm{tv }ārambhe}
		\rdg[wit={J10}]{āraṃbho}
		\rdg[wit={V3}]{ārambha}
		}
	\app{\lem[wit={ceteri},alt={yogavān}]{yogavā\skp{n}}
		\rdg[wit={N23}]{bhogavān}
		}n bhavet//\versenr}
		\orgvnr{21}\\!}
\end{tlg}
%</vs21>
\commcite\newpage


%<*vs22a>
\begin{ava}[hp04_022a]
atha
\app{\lem[wit={ceteri}]{ghaṭāvasthā}% ghaṭavaṃsthā N19, ghaṃṭā° V17
		\rdg[wit={J5}]{ghaṭā arthaḥ}
		\rdg[wit={G4}]{khaṭavasthā}
		\rdg[wit={Delta}]{ghaṭarakṣā}
		\rdg[wit={P11}]{ghaṭaḥ}}/
\end{ava}
%</vs22a>
%<*vs22>
\begin{tlg}[hp04_022]
\tl{
\pada{\app{\lem[wit={ceteri}]{dvitīyāyāṃ} % dviti° V3
		\rdg[wit={J5}]{dvitī}
		\rdg[wit={V19,V15ac}]{dvitīyā}
		\rdg[wit={J10}]{dvitīye}
		}
	\app{\lem[wit={N3,Gamma,Delta,Pi,Jyo}]{ghaṭī}
		\rdg[wit={J5}]{ghaṭikā}
		\rdg[wit={Epsi}]{sphuṭī}% +M3
		\rdg[wit={N19}]{ghaṭāṃ}
		\rdg[wit={V15}]{ghaṃṭi}
		\rdg[wit={J10}]{bheda}
		}%
	\app{\lem[wit={ceteri}]{kṛtya}
		\rdg[wit={V15}]{kṛtvā}
		\rdg[wit={J10}]{mukte tu}}}
\pada{vāyur bhavati madhyagaḥ/}\\+} % vāḥsur P7
%	\app{\lem[wit={ceteri}]{madhyagaḥ}
%		\rdg[wit={K3,C7}]{madhyamaḥ}}/}\\+}
\tl{
\pada{\app{\lem[wit={ceteri}]{dṛḍhāsano}% dṛṣaḥsa? G5
		\rdg[wit={J10}]{haṭhāsano}} bhaved yogī}
\pada{jñānī
	\app{\lem[wit={N3,J5,Gamma,V19,Epsi,Zeta2,Jyo}]{deva}
		\rdg[wit={E2,J10,P11,C6}]{deha}
		\rdg[wit={V3}]{devaḥ}
		}sama%s  % samās N3
	\app{\lem[wit={N3,J5,G5,Pi,Jyo},alt={tadā}]{\skm{s }tadā}
		\rdg[wit={Gamma,Delta,G11,Zeta2,J10}]{tathā}}//\versenr}
		\orgvnr{22}\\!}
\end{tlg}
%</vs22>

\avacite{22a}
\commcite\newpage


%<*vs23>
\begin{tlg}[hp04_023]
\tl{
\pada{viṣṇu\app{\lem[wit={N3,P11}]{granthes tadā}% +N24,F
		\rdg[wit={J5,G5,J10}]{granthes tathā}
		\rdg[wit={Gamma,Delta,V15}]{granthir yadā}% <<graṃthi>> C7
		\rdg[wit={G11}]{granthe tathā}
		\rdg[wit={N19}]{granthe sadā}
		\rdg[wit={C6}]{granther yadā}
		\rdg[wit={V3}]{granthis tadā}
		\rdg[wit={Jyo}]{granthes tato}
		}
	\app{\lem[wit={N3,Epsi,N19,J10,Pi,Jyo},alt={bhedāt}]{bhedā\skp{t}}
		\rdg[wit={J5}]{bhidā}
		\rdg[wit={Gamma,Delta}]{bhinnaḥ} % vimugraṃthipadābhinna N23
		\rdg[wit={V15}]{bhinnā}
		}}%
\pada{\app{\lem[wit={ceteri},alt={paramānanda}]{\skm{t }paramānanda}
		\rdg[wit={N19}]{sadānandasya}}%
	\app{\lem[wit={ceteri}]{sūcakaḥ} % śū° N23, °caka V3
		\rdg[wit={V15}]{sūcakā<<ḥ>>}
		\rdg[wit={C6}]{kārakaḥ}}/}\\+}% °kāḥ V15pc
\tl{
\pada{\app{\lem[wit={Alpha,Epsi,P11,V3,Jyo}]{atiśūnye}% śunye N3
		\rdg[wit={Gamma,Delta,V15,J10}]{atiśūnya}
		\rdg[wit={N19}]{api śūnyo}
		\rdg[wit={C6}]{aṃtyaśūnye}
		}
	\app{\lem[wit={N3,G4,Pi,Jyo}]{vimardaś ca}
		\rdg[wit={J5}]{vimardasya}
		\rdg[wit={Gamma,Delta,V15}]{vibhedaś ca}
		\rdg[wit={G11}]{visanmarde}
		\rdg[wit={G5}]{pi sammardo}
		\rdg[wit={N19}]{'saṃmardā}
		\rdg[wit={J10}]{visaṃmardo}
		}}
\pada{bherīśabda%s % sabdas V3
	\app{\lem[wit={N3,V15,Pi,Jyo},alt={tadā}]{\skm{s }tadā}
		\rdg[wit={J5}]{tatho}
		\rdg[wit={G4,Gamma,Delta,Epsi,N19,J10}]{tathā}% +N24
		} bhavet//\versenr}
		\orgvnr{23}\\!}
\end{tlg}
%</vs23>
\commcite\newpage


%<*vs24a>
\begin{ava}[hp04_024a]
\app{\lem[wit={ceteri}]{atha}
		\rdg[wit={E2}]{\om}
		\rdg[wit={C6}]{tathā}
		}
\app{\lem[wit={ceteri}]{paricayāvasthā}% <<yā>> N23
		\rdg[wit={Zeta2,P11}]{paricayaḥ}}/ % °caya P11,N19
\end{ava}
%</vs24a>
%<*vs24>
\begin{tlg}[hp04_024]
\tl{
\pada{\app{\lem[wit={N3,Delta,Epsi,V15,Pi}]{tṛtīyāyāṃ tato bhittvā} % tṛtiyāyāṃ V3
		\rdg[wit={J5}]{tṛtīyāyāṃ tathā bhitvā}
		\rdg[wit={Gamma}]{karṇikāṃ tu tato bhittvā}
		\rdg[wit={N19}]{karttikāyāṃ tato bhittvā}
		\rdg[wit={J10}]{atha granthitrayaṃ bhittvā}
		\rdg[wit={Jyo}]{tṛtīyāyāṃ tu vijñeyo}
		}}
\pada{\app{\lem[nolem]{\skp{pāda b}}
	\rdg[wit={N3}]{\unavbl}}%
\app{\lem[wit={J5,G11,N19,Jyo}]{vihāyo}
		\rdg[wit={Gamma,V15}]{vihāya}
		\rdg[wit={Delta}]{vimalo}
		\rdg[wit={G5}]{vihāye}
		\rdg[wit={J10}]{jāyate}
		\rdg[wit={P11}]{vikāryo}
		\rdg[wit={C6}]{visphāro}
		\rdg[wit={V3}]{vimāyo}
		}%
	\app{\lem[wit={J5,Gamma,N19,J10,Pi,Jyo}]{mardala}
		\rdg[wit={Delta}]{mandala} % maṃdala V19
		\rdg[wit={G11}]{maddala}
		\rdg[wit={G5}]{daṃviya}
		\rdg[wit={V15}]{mṛḍula}}%
	\app{\lem[wit={ceteri}]{dhvaniḥ}
		\rdg[wit={J7}]{dhvaniṃ}
		\rdg[wit={P11,V3}]{dhvani}}/}\marma\\+}
\tl{
\pada{\app{\lem[nolem]{\skp{pāda c}}
	\rdg[wit={N3}]{\unavbl}}%
\app{\lem[wit={ceteri}]{mahāśūnyaṃ}
		\rdg[wit={G11}]{mahāśūnyas}
		\rdg[wit={V15,P11}]{mahāśūnya}
		}
	\app{\lem[wit={J5,G11,Pi,Jyo}]{tadā}% +N24
		\rdg[wit={Gamma,N19}]{tathā}
		\rdg[wit={Delta,G5}]{tato}
		\rdg[wit={V15}]{tayā}
		\rdg[wit={J10}]{samā}
		}
	\app{\lem[wit={ceteri}]{yāti}
		\rdg[wit={J5}]{jāti}
		\rdg[wit={N19}]{jātiḥ}
		}}
\pada{\app{\lem[nolem]{\skp{pāda d}}
	\rdg[wit={N3}]{\unavbl}}%
\app{\lem[wit={ceteri}]{sarvasiddhi} % +P11
		\rdg[wit={G5}]{tatsaṃsiddhis}
		\rdg[wit={N19}]{sarva}
		\rdg[wit={C6}]{siddhisādha}
		\rdg[wit={V3}]{mahāsiddhi}
		}% 
	\app{\lem[wit={ceteri}]{samāśrayam}
		\rdg[wit={J5}]{matāśrayāt}
		\rdg[wit={G5}]{tadāśrayaḥ}
		\rdg[wit={P11}]{samāśriyaṃ}
		\rdg[wit={C6}]{kam āśrayaṃ}}//\versenr}
	\orgvnr{24}\\!}
\end{tlg}
%</vs24>

\avacite{24a}
\commcite\newpage

%<*vs25>
\begin{tlg}[hp04_025]
\tl{\app{\lem[nolem]{}
	\rdg[wit={N3}]{\unavbl}}%
\pada{\app{\lem[wit={G4,Gamma,Delta,G11,C6,Jyo}]{cittānandaṃ}% <<da>> C7; +N24
		\rdg[wit={J5,V15,V3}]{cidānaṃda}% cidānaṃdaṃ HR,F, cidātmānaṃ M3 ##
		\rdg[wit={G5}]{cittamānaṃ}
		\rdg[wit={N19}]{virāmānaṃ}
		\rdg[wit={J10}]{ciṃtāmanas}
		\rdg[wit={P11}]{vivarttānaṃdaṃ}
		}
	\app{\lem[wit={ceteri}]{tato}
		\rdg[wit={Jyo}]{tadā}}
	\app{\lem[wit={ceteri}]{jitvā}% J5,G4,G11,N19,V15,J10,Pi,Jyo
		\rdg[wit={Gamma,Delta}]{bhittvā}
		\rdg[wit={G5}]{hitvā}}}
\pada{sahajānandasaṃbhavaḥ/}% ḥ om. N19,P11
	% \app{\lem[wit={ceteri}]{saṃbhavaḥ}
		% \rdg[wit={N19,P11}]{saṃbhava}}
		\\+}
\tl{
\pada{\app{\lem[wit={ceteri}]{doṣaduḥkha}
		\rdg[wit={N23}]{dokhaduḥkhe}
		\rdg[wit={P11}]{doṣaduḥkhaṃ}
		}%
	\app{\lem[wit={G4,Epsi,V15,J10,Pi}]{jarāmṛtyu}
		\rdg[wit={J5,N19}]{jarāmṛtyuḥ}
		\rdg[wit={Gamma,Delta}]{kṣudhānidrā}
		\rdg[wit={Jyo}]{jarāvyādhi}
		}}%
\pada{\app{\lem[wit={J5,G4,Epsi,Zeta2,J10,P11,C6,Jyo}]{kṣudhānidrā}
		\rdg[wit={Gamma,Delta}]{jarāmṛtyu}% mṛrtyu N23
		\rdg[wit={V3}]{kṣudhātṛṣā}
		}% 
	\app{\lem[wit={ceteri}]{vivarjitaḥ}
		\rdg[wit={G5}]{vivarjakaḥ}
		\rdg[wit={J10}]{tṛṣā tathā}
		\rdg[wit={C6}]{vivarjitāḥ}
		\rdg[wit={V3}]{vivarjitaṃ}
		}//\versenr}
	\orgvnr{25}\\!}
\end{tlg}
%</vs25>
\commcite\newpage


%<*vs26a>
\begin{ava}[hp04_026a]
\app{\lem[nolem]{}
	\rdg[wit={N3}]{\unavbl}
	\rdg[wit={J5,J7,Delta,Pi}]{\textapp{found after \manuref{4.26b/X4.81b}}}
	\rdg[wit={Jyo}]{\om}
	}%	
	atha 
	\app{\lem[wit={Gamma,C6,V3}]{niṣpattyavasthā} % niḥpatya° N23, niṣpatya° J7
		\rdg[wit={J5}]{niḥpatti-avasthā}
		\rdg[wit={Delta}]{niṣṭhāvasthā}
		\rdg[wit={Epsi,Zeta2,J10,P11}]{niṣpattiḥ}% niṣpatti P11, niṣpati N19, °tiḥ J10; +F ##
		}/
\end{ava}
%</vs26a>
%<*vs26>
\begin{tlg}[hp04_026]
\tl{\app{\lem[nolem]{}
	\rdg[wit={N3}]{\unavbl}}%
\pada{rudragranthiṃ % urddha J5; rudraṃ graṃ{{ga}}thi N23; grathiṃ J7, graṃthi C7,N19,P11,V3, grandhiṃ G11
	\app{\lem[wit={ceteri}]{tato}
		\rdg[wit={Jyo}]{yadā}}
	\app{\lem[wit={ceteri}]{bhittvā} % bhītvā J5,P7
		\rdg[wit={N19}]{bhūtvā}}}
\pada{\app{\lem[wit={ceteri}]{sarva}
		\rdg[wit={P11}]{satva}
		\rdg[wit={Jyo}]{śarva}
		}pīṭha% pīca P11, pīḍa G11
	\app{\lem[wit={ceteri},alt={gato 'nilaḥ}]{gato'nilaḥ}% +M3
		\rdg[wit={J5,V3}]{gatānila}
		\rdg[wit={J7}]{gatonalaḥ}
		\rdg[wit={Epsi}]{gatānilaḥ/ṃ}% °laṃ G5
		}/}\\+}
\tl{
\pada{\app{\lem[wit={J5,J7,Pi,Jyo}]{niṣpattau}
		\rdg[wit={N23}]{niṣpatto}
		\rdg[wit={Delta}]{niṣṭhāto}
		\rdg[wit={G11,J10}]{niṣpanno} % niḥpanno J10
		\rdg[wit={G5}]{niṣpanne}
		\rdg[wit={Zeta2}]{niṣpannau} % niḥṣpannau? V15
		}
	\app{\lem[wit={ceteri}]{vaiṇavaḥ śabdaḥ}% veṇavaḥ V17
		\rdg[wit={J5}]{vauṇāvat sado}
		\rdg[wit={N23}]{veṇacaśabdaṃ}
		\rdg[wit={J7}]{vaiṇavaśabdaḥ}
		}}
\pada{\app{\lem[wit={V15,Jyo}]{kvaṇadvīṇākvaṇo}% viṇā V15; +F
		\rdg[wit={J5}]{kṛṇanityakṛṇo}
		\rdg[wit={N23}]{karṇavīṇādgato}
		\rdg[wit={J7}]{kvaṇadvīṇotvaṇo}% < °vīṇolbaṇo?
		\rdg[wit={Delta}]{kvaṇadvīṇāsamo}% vīnā V19; kvaṇadvīvakvaṇo K1
		\rdg[wit={G11}]{kvaṇan vīṇakvaṇo}
		\rdg[wit={G5}]{kvaṇan gītakvaṇo}
		\rdg[wit={N19}]{kaṇatvītakvaṇo}% sic!
		\rdg[wit={J10}]{kvaṇantenākvuṇo}
		\rdg[wit={P11}]{kvaṇan vītaḥ kvaṇo} 
		\rdg[wit={C6}]{kvacid vīṇākvaṇo}
		\rdg[wit={V3}]{kvaṇatuvītakvaṇo}
		}\marmas
	\app{\lem[wit={ceteri}]{bhavet}
		\rdg[wit={C6}]{dayaḥ}}//\versenr}\label{rudra}
	\orgvnr{26}\\!}
\end{tlg}
%</vs26>

\avacite{26a}
\commcite\newpage

% MD: marking für die andere Recensions hier abgebrochen.

%<*vs27>
\begin{tlg}[hp04_027]
\tl{\app{\lem[nolem]{}
	\rdg[wit={N3}]{\unavbl}
	\rdg[wit={Zeta2}]{\om}}%
\pada{ekībhūtaṃ % ekāṃ P11, eka° C6
	\app{\lem[wit={J5,Epsi,Pi,Jyo}]{tadā}
		\rdg[wit={G4,Gamma,Delta,J10}]{tathā}} cittaṃ}
\pada{\app{\lem[nolem]{\skp{pāda b}}
	\rdg[wit={P11}]{\om}}%
\app{\lem[wit={ceteri},alt={rājayogā°}]{rājayogā\skp{°}}
		\rdg[wit={J10}]{rājayoga}
		\rdg[wit={V3}]{rājayogo}}%
	\app{\lem[wit={J7,G11,V3},alt={°bhidhāyakam}]{\skp{°}bhidhāyakam}
		\rdg[wit={J5}]{vidhāyakaḥ}
		\rdg[wit={G4,Delta,G5,J10,C6,Jyo}]{bhidhānakaṃ}% +F
		\rdg[wit={N23}]{bhidhāyanaṃ}
		}\marma/}\\+} % +N24
\tl{
\pada{\app{\lem[nolem]{\skp{pāda c}}
	\rdg[wit={P11}]{\om}}%
	sṛṣṭisaṃhārakartāsau} % karttasau N23, karttāso V3
\pada{\app{\lem[nolem]{\skp{pāda d}}
	\rdg[wit={P11}]{\om}}%
	yogīśvarasamo bhavet//\versenr}%
	%\myfnx{After this verse, \getsiglum{Jyo} has X4.117, 119, and 121ab.}
	\orgvnr{27}\\!}
\end{tlg}
%</vs27>
\commcite%\newpage


\iffalse
%<*vs27-1a>
\begin{altava}[hp04_027_1a]
\app{\lem[nolem]{}
	\rdg[wit={Epsi}]{\only}}%
	atha nādānusandhānam/
\end{altava}
%</vs27-1a>

\startaltrecension % B1 has this verse too.
%<*vs27-1>
\begin{alttlg}[hp04_027_1]
\tl{%\app{\lem[nolem]{}
	%\rdg[wit={Pi},alt=\textapp{found after \ref{astuva}}]{\skp{found after xxx}}}%
\pada{rājayoga\app{\lem[wit={G5,P11,C6}]{padaprāptau}% +C8,P7,M3,F
		\rdg[wit={G11}]{padaprāptā}
		\rdg[wit={N19}]{padaprāptaḥ}
		\rdg[wit={V15}]{padaṃ prāpti}
		\rdg[wit={J10,Jyo}]{padaṃ prāptuṃ}
		\rdg[wit={V3}]{padaṃ prāptaṃ}
		}}
\pada{\app{\lem[wit={Epsi,Zeta2,J10,P11,C6,Jyo},alt={sukhopāyo 'lpa}]{sukhopāyo'lpa}
		\rdg[wit={V3}]{sukhopāyogya}}cetasām/}\\+} % cetasam P11
\tl{\app{\lem[alt=\mylem{83cd},nosep,nonum]{\skp{pāda cd}}%
	\rdg{\texteng{cf. \manuref{X4.122cd (4.35cd)}}}}%
\pada{sadyaḥpratyaya%
	\app{\lem[wit={N19,J10,C6,V3,Jyo}]{saṃdhāyī}
		\rdg[wit={Epsi}]{saṃdāyī}
		\rdg[wit={V15,P11}]{saṃdhāyi}
		}}
\pada{\app{\lem[wit={G5,Zeta2,Pi,Jyo}]{jāyate}% +M3
		\rdg[wit={G11}]{līyate}
		\rdg[wit={J10}]{sevyate}
		}
	\app{\lem[wit={N19,C6,V3,Jyo}]{nādajo layaḥ} % laya V3
		\rdg[wit={G5}]{nādamūlayā}
		\rdg[wit={V15}]{nātra saṃśayaḥ}
		\rdg[wit={J10,P11}]{nādayo layaḥ}
		}//\versenr}
	%\sgwit{G11,G5,N19,V15,J10,Pi,Jyo}
	%\myfn{cf. \ref{saukhya}cd}
	\\!}
\end{alttlg}
%</vs27-1>
\endaltrecension
\fi



%<*vs28>
\begin{tlg}[hp04_028]
\tl{\app{\lem[nolem]{}
	\rdg[wit={N3}]{\unavbl}
	\rdg[wit={P11}]{\om}}%
\pada{astu vā
	\app{\lem[wit={J5,Gamma,Delta,Epsi,N19,J10,C6,Jyo}]{māstu} % māsta N19
		\rdg[wit={V15,V3}]{mastu}
		} vā
	\app{\lem[wit={J5,E2,Epsi,N19,J10,C6,Jyo},alt={muktir}]{mukti\skp{r}}% +K3,C7
		\rdg[wit={Gamma}]{śaktir}
		\rdg[wit={V19}]{kiṃcid}
		\rdg[wit={V15}]{muktis}
		\rdg[wit={V3}]{muktiṃ}
		}}%
\pada{\app{\lem[wit={Delta,Epsi,N19,C6,Jyo},alt={atraivākhaṇḍitaṃ}]{\skm{r }atraivākhaṇḍitaṃ}% °ṣaṃḍitaṃ N19
		\rdg[wit={J5,J7}]{atraiva khaṇḍitaṃ}% ṣaṃḍitaṃ J5
		\rdg[wit={N23}]{ātrevikhaṇḍitaṃ}
		\rdg[wit={V15,V3}]{tatraivākhaṇḍitaṃ}% +F
		\rdg[wit={J10}]{atra vākhaṇḍitaṃ}
		}
	\app{\lem[wit={J5,J7,E2,Epsi,Zeta2,J10,V3}]{mahat}
		\rdg[wit={N23}]{marut}
		\rdg[wit={V19}]{bhavet}
		\rdg[wit={C6}]{manaḥ}
		\rdg[wit={Jyo}]{sukham}}/}\\+}
\tl{
\pada{\app{\lem[wit={J5,Epsi,Zeta2,C6}]{layāmṛtamayaṃ}% leyā° J5, staye° N24
		\rdg[wit={N23}]{layāmṛdaṃmitaṃ}
		\rdg[wit={J7,Delta}]{layāmṛtam idaṃ}% +F
		\rdg[wit={J10}]{layāmṛtakaraṃ}
		\rdg[wit={V3}]{layāmṛtalayaṃ}
		\rdg[wit={Jyo}]{layodbhavam idaṃ}}
	% \app{\lem[wit={Delta1,G11,V15,C6,V3,Jyo}]{saukhyaṃ}
		% \rdg[wit={J5,J7,J10}]{sauṣyaṃ}
		% \rdg[wit={N23}]{sokhyaṃ}
		% \rdg[wit={N19}]{saukṣaṃ}}
		saukhyaṃ}
\pada{\app{\lem[wit={ceteri}]{rājayogād avāpyate}% avāpnute G5
		\rdg[wit={V19}]{\om}
		\rdg[wit={J10}]{rājayogam avāpyate}
		}//\versenr}%
		\label{astuva}
	\orgvnr{28}\\!}
\end{tlg}
%</vs28>
\commcite\newpage


%<*vs29>
\begin{tlg}[hp04_029]
\tl{\app{\lem[nolem]{}
	\rdg{\texteng{=\,2.77}}
	\rdg[wit={N3}]{\unavbl}
	\rdg[wit={Epsi,J10,V3,Jyo}]{\om}% +M3
	\rdg[wit={Zeta2}]{\textapp{abbreviated to} haṭhaṃ vinā rājayoga iti}
	}%
\pada{\app{\lem[nolem]{\skp{pāda a}}
	\rdg[wit={V19}]{\om}}%
	haṭhaṃ vinā rāja%
	\app{\lem[wit={J5,N23,C6},alt={rājayogo}]{yogo}
		\rdg[wit={G4,J7,E2},alt={°yogaṃ}]{\skp{°}yogaṃ}
		\rdg[wit={P11},alt={°yogaḥ}]{\skp{°}yogaḥ}
		}} 
\pada{rājayogaṃ vinā % yoga J5,J11
	\app{\lem[wit={J5,G4,N23,C6}]{haṭhaḥ}% N23ac
		\rdg[wit={J7,Delta,P11}]{haṭhaṃ}% +N23pc
		}/}\\+}
\tl{
\pada{\app{\lem[nolem]{\skp{pāda c}}
	\rdg[wit={Gamma,Delta}]{\om}}%
	na sidhyati tato yugma}%
\pada{m ā\app{\lem[nolem]{\skp{pāda d}}
	\rdg[wit={Gamma,Delta}]{\om}}
	niṣpatteḥ samabhyaset//\versenr}% +M1,J11; °ttiḥ P11, °tte C6
	\label{hathamvina}
	\orgvnr{29}\\!}
\end{tlg}
%</vs29>
\commcite%\newpage


%<*vs30>
\begin{tlg}[hp04_030]
\tl{\app{\lem[nolem]{}
	\rdg[wit={N3}]{\unavbl}}%
\pada{\app{\lem[nolem]{\skp{pāda a}}
	\rdg[wit={Gamma,Delta}]{\om}}%
	rājayogam ajānantaḥ} % °nataḥ V3, °naṃteḥ N19
\pada{\app{\lem[nolem]{\skp{pāda b}}
	\rdg[wit={Gamma,Delta}]{\om}}%
	kevalaṃ
	haṭha\app{\lem[wit={G11,V15,P11}]{karmaṭhāḥ}
		\rdg[wit={J5}]{karmaṭhaḥ}
		\rdg[wit={G5}]{karmataḥ}
		\rdg[wit={N19}]{karmacā}
		\rdg[wit={J10}]{karmaṇaḥ}
		\rdg[wit={C6,V3}]{karmaṇā}
		\rdg[wit={Jyo}]{karmiṇaḥ}}/}\\+}
\tl{
\pada{\app{\lem[wit={Epsi,P11,C6}]{ye tu tān karṣakān manye}
		\rdg[wit={J5}]{ye ca te kāmukān manne}
		\rdg[wit={Gamma}]{ye vai varanti tān manye}
		\rdg[wit={Delta}]{ye vai caranti tān manye}% madhye J7
		\rdg[wit={Zeta2}]{ye tu tān karkaśān manye}
		\rdg[wit={J10}]{ye tuṃgān karmavasān manye}
		\rdg[wit={V3}]{\lacuna}
		\rdg[wit={Jyo}]{etān abhyāsino manye}
		}} % V3 gap
\pada{\app{\lem[wit={ceteri}]{prayāsaphalavarjitān}%ptalevi N23 
		\rdg[wit={J5}]{prayāsakalavarjitaḥ}
		\rdg[wit={J10}]{prayāsaphalavarjitāḥ}
		\rdg[wit={C6}]{prāyaśaḥ phalavarjitāḥ}
		\rdg[wit={V3}]{\lacuna}
		}//\versenr}%
	\xinclude{\myfnx{\getsiglum{N19} ends with this verse.}}
	\label{ajananta}
	\orgvnr{30}\\!}
\end{tlg}
%</vs30>
\commcite\newpage


%\startaltrecension
%%<*vs30-1>
%\begin{alttlg}[hp04_030_1]
%\tl{\pada{\app{\lem[wit={Gamma,E2}]{haṭhaṃ vinā}
%		 \rdg[wit={V19}]{\om}}
%	 \app{\lem[wit={J7,E2}]{rājayogaṃ}
%		 \rdg[wit={N23}]{rājayogo}
%		 \rdg[wit={V19}]{\om}}}
%\pada{rājayogaṃ vinā
%	 \app{\lem[wit={J7,Delta}]{haṭham}
%		 \rdg[wit={N23}]{haṭhaḥ}}/}\\+} % hathaḥ N23ac
%\tl{
%\pada{ye \app{\lem[wit={N23,Delta}]{vai}
%		 \rdg[wit={J7}]{cai}}
%	 \app{\lem[wit={Delta}]{caranti}
%		 \rdg[wit={Gamma}]{varaṃti}} tā%n
%	 \app{\lem[wit={N23,Delta},alt={manye}]{\skm{n }manye}
%		 \rdg[wit={J7}]{madhye}}}
%\pada{prayāsa\app{\lem[wit={J7,Delta}]{phala} % varjitāḥ J10
%		 \rdg[wit={N23}]{ptalevi}}varjitān//\versenr}
%		 \sgwit{Gamma,Delta}
%	\myfn{This verse is found in \getsiglum{Gamma,Delta} as a substitute for \ref{hathamvina}--\ref{ajananta}.}
%	\\!}
%\end{alttlg}
%%</vs30-1>
%\endaltrecension
%
%\altcommcite
%\newpage



%<*vs31>
\begin{tlg}[hp04_031]
\tl{\app{\lem[nolem]{}
	\rdg[wit={N3}]{\unavbl}
	\rdg[wit={E2}]{\om}}%
\pada{\app{\lem[wit={ceteri}]{tattvaṃ}% +M3
		\rdg[wit={N23,G11,V3}]{tattva}}\marmas 
		bījaṃ
	\app{\lem[wit={V19,Jyo}]{haṭhaḥ}% +M3
		\rdg[wit={J5,Gamma,Epsi,V15,P11}]{haṭha}% +J11
		\rdg[wit={G4,J10,C6,V3}]{haṭhaṃ}% +F,K3,C7
		} kṣetra}%m 
\pada{\app{\lem[wit={Gamma,J10,C6,V3,Jyo},alt={audāsīnyaṃ}]{\skm{m }audāsīnyaṃ}% +K3,C7
		\rdg[wit={J5,V15,P11}]{audāsinyaṃ}
		\rdg[wit={G4}]{audāśinyaṃ}
		\rdg[wit={Epsi}]{audāsīnya}
		\rdg[wit={V19}]{<<sau>>dāmanyaṃ}}
	\app{\lem[wit={J5,G5,V15,J10,P11,V3,Jyo}]{jalaṃ tribhiḥ}% tribhi V3
		\rdg[wit={G4,Gamma,V19,C6}]{jalaṃ smṛtam}
		\rdg[wit={G11}]{layaṃ tribhiḥ}
		}/}\\+}
\tl{
\pada{unmanīkalpalatikā}
\pada{sadya
	\app{\lem[wit={J5,V19,Epsi,V15,J10,C6,V3}]{evodbhaviṣyati}% yevo C6; K3,C7
		\rdg[wit={G4,Gamma}]{eva bhaviṣyati}% +F
		\rdg[wit={P11}]{evādbhaviṣyati}
		\rdg[wit={Jyo}]{eva pravartate}}//\versenr}
	\orgvnr{31}\\!}
\end{tlg}
%</vs31>
\commcite\newpage


%<*vs32>
\begin{tlg}[hp04_032]
\tl{\app{\lem[nolem]{}
	\rdg[wit={N3}]{\unavbl}}%
\pada{\app{\lem[wit={J7,Delta,G11,V15,V3,Jyo}]{rājayogaḥ}
		\rdg[wit={J5,N23,G5,N19,J10,P11,C6a,C6b}]{rājayoga}
		}
	\app{\lem[wit={ceteri}]{samādhiś ca}
		\rdg[wit={Zeta2}]{samādhiḥ syād}% samādhi V15; +F,N22,P6
		\rdg[wit={P11}]{samādhīś cā}
		}}
\pada{\app{\lem[wit={ceteri}]{unmanī}
		\rdg[wit={J5,G11}]{py unmanī}
		\rdg[wit={P11}]{nmatī}
		\rdg[wit={C6b}]{hy unmanī}
		} ca manonmanī/}\\+} % va P11
\tl{
\pada{\app{\lem[wit={V15,J10}]{amaraugho}
		\rdg[wit={J5,P11}]{amarodyo}
		\rdg[wit={N23}]{araughau}
		\rdg[wit={J7}]{amaraudhyai}
		\rdg[wit={Delta}]{amaroly a}
		\rdg[wit={G11,C6b}]{amaraughā}
		\rdg[wit={G5}]{aromaro}
		\rdg[wit={N19}]{avaraubhū}
		\rdg[wit={C6a}]{amaraughi}
		\rdg[wit={V3}]{amarogho}
		\rdg[wit={Jyo}]{amaratvaṃ}
		}%
	\app{\lem[wit={J5,P11,C6b,V3}]{'pi cādvaitaṃ}% vādvaitaṃ P11
		\rdg[wit={N23}]{ghatvīṃdrī ca}
		\rdg[wit={J7}]{ghacāṃdrī ca}
		\rdg[wit={Delta}]{bhicāndrī ca}
		\rdg[wit={G11,Zeta2,J10,C6a,Jyo}]{layas tattvaṃ}% layes V15
		\rdg[wit={G5}]{yas tat tulyaḥ}
		\rdg[wit={J10}]{layas tatra}
		}}
\pada{\app{\lem[wit={J5,Gamma,P11,C6b,V3}]{nirālambaṃ}
		\rdg[wit={Delta}]{nirālambo}
		\rdg[wit={G11,Zeta2,J10,Jyo}]{śūnyāśūnyaṃ} % °śūnya N19
		\rdg[wit={G5}]{śūnyāt śūnya}
		\rdg[wit={C6a}]{śūnyācūnyaṃ}
		} 
	\app{\lem[wit={ceteri}]{nirañjanam}
		\rdg[wit={J5}]{nirāmayaṃ}
		\rdg[wit={Epsi,Zeta2,J10,C6a,Jyo}]{paraṃ padam}
		}//\versenr}\label{A1}%
	\myfn{\getsiglum{C6} has this pair of verses twice: first (\getsiglum{C6a}) as X4.3–4 of the expanded version, and second (\getsiglum{C6b}) as 4.32–33 of the older version.}
	\orgvnr{32}\\!}
\end{tlg}
%</vs32>
%\commcite%\newpage


%<*vs33>
\begin{tlg}[hp04_033]
\tl{\app{\lem[nolem]{}
	\rdg[wit={N3}]{\unavbl}
	}% C6 twice
\pada{\app{\lem[wit={J7,V19,P11,C6b,V3}]{amanasko} % °skoṃ? V19
		\rdg[wit={J5}]{amarasko}
		\rdg[wit={N23}]{amanaskau}
		\rdg[wit={E2,Epsi,Zeta2,J10,C6a,Jyo}]{amanaskaṃ} % +K3,C7
		}
	\app{\lem[wit={P11,C6b,V3}]{layas tattvaṃ}
		\rdg[wit={J5}]{layas tatra}
		\rdg[wit={N23}]{lyayāś caiva}
		\rdg[wit={J7,Delta}]{layaś caiva}
		\rdg[wit={Epsi,Zeta2,J10,C6a,Jyo}]{tathādvaitaṃ}
		}}
\pada{\app{\lem[wit={J5,J7,Delta,P11}]{śūnyāśūnyaṃ}% °sūnyaṃ J5
		\rdg[wit={N23,V3}]{śūnyāśūnya}
		\rdg[wit={Epsi,Zeta2,J10,C6a,Jyo}]{nirālambaṃ}
		\rdg[wit={C6}]{śūnyāc chūnyaṃ}
		}
	\app{\lem[wit={J5,G4,P11,C6b,V3}]{paraṃ padam}
		\rdg[wit={N23,Delta}]{parāparaṃ}
		\rdg[wit={J7}]{parāvaraṃ}
		\rdg[wit={Epsi,Zeta2,J10,C6a,Jyo}]{nirañjanam}
		}/}\\+}
\tl{
\pada{\app{\lem[wit={ceteri}]{jīvanmuktiś ca}
		\rdg[wit={G4,G5}]{jīvanmuktaś ca}
		\rdg[wit={N23}]{jīvanmuktiḥ}} 
	\app{\lem[wit={ceteri}]{sahajaṃ} % ṃ om. N23
%		\rdg[wit={Jyo}]{sahajā}
		\rdg[wit={V15}]{\om}
		}}
\pada{\app{\lem[wit={J5,G4,Gamma,E2,V15,P11,C6a}]{turyaṃ} % turya J5
		\rdg[wit={V19}]{turjaṃ}
		\rdg[wit={G11}]{turīyaṃ}
		\rdg[wit={G5}]{tulyaṃ}
		\rdg[wit={N19}]{turyai}
		\rdg[wit={J10}]{muktiś}
		\rdg[wit={C6}]{turyāṃ}
		\rdg[wit={V3}]{tuṣkaṃ}
		\rdg[wit={Jyo}]{turyā}
		}
	\app{\lem[wit={J5,J7,Delta,G5,J10,P11,C6,Jyo}]{cety eka}
		\rdg[wit={G4}]{..\,ty eka}
		\rdg[wit={N23}]{vatyaka}
%		\rdg[wit={C7}]{cety eva}
		\rdg[wit={G11}]{caika}
		\rdg[wit={N19}]{ciṃtaika}
		\rdg[wit={V15}]{cittaika}% +N22,P6
		\rdg[wit={C6a}]{caityeka}
		\rdg[wit={V3}]{caiyeka}
		}%
	\app{\lem[wit={N23,Jyo}]{vācakāḥ}% +J10pc
		\rdg[wit={J5,J10}]{vācakaḥ}
		\rdg[wit={G4,Delta,Epsi,Zeta2,P11,C6a,C6b,V3}]{vācakaṃ}% +N22,P6
		\rdg[wit={J7}]{vācakīṃ}
		}//\versenr}\label{A2}
	\orgvnr{33}\\!}
\end{tlg}
%</vs33>

\commciterange{33}{32–33}\newpage


%<*vs34>
\begin{tlg}[hp04_034]
\tl{\app{\lem[nolem]{}
	\rdg[wit={N3}]{\unavbl}
	\rdg[wit={E2,Zeta2,J10}]{\om}
	}%
\pada{\app{\lem[nolem]{\skp{pāda a}}
	\rdg[wit={Gamma}]{\om}}%
\app{\lem[wit={J5,Epsi,P11,V3,Jyo}]{unmanyavāptaye}% °staye G5
		\rdg[wit={G4}]{unmanyaye}
		\rdg[wit={V19}]{unmanyavāsayet}
		%\rdg[wit={K3,C7}]{unmanyā vāsayec}
		\rdg[wit={C6}]{unmanyā\,\_\,\_\,ye}
		} śīghraṃ} % chīghraṃ V19,K3,C7
\pada{\app{\lem[nolem]{\skp{pāda b}}
	\rdg[wit={Gamma}]{\om}}%
\app{\lem[wit={J5,Epsi,P11,C6}]{mārgau dvau}
		\rdg[wit={G4}]{mārgā\,..}
		\rdg[wit={V19}]{dvau mārgau}% +K3,C7
		\rdg[wit={V3}]{mārgo dvau}
		\rdg[wit={Jyo}]{bhrūdhyānaṃ}
		}
	\app{\lem[wit={J5,Epsi,V3}]{mama saṃmatau}
		\rdg[wit={G4}]{myama saṃ[m].\,+}
		\rdg[wit={V19,P11}]{samasaṃmatau}% +K3,C7
		\rdg[wit={C6}]{mamatau}
		\rdg[wit={Jyo}]{mama saṃmatam}
		}/}
	\\+}
\tl{
\pada{\app{\lem[nolem]{\skp{pāda c}}
	\rdg[wit={V19,Jyo}]{\om}}%
tattvaṃ parama% padaṃ J5
	\app{\lem[wit={Gamma,G5,C6}]{saukhyaṃ} % 2 x ṃ om. N23
		\rdg[wit={J5}]{sākhyaṃ}
		\rdg[wit={G11,V3}]{sāṃkhyaṃ}% saṃkhyaṃ F
		\rdg[wit={P11}]{vāgraṃ}
		} vā} % ca J5
\pada{\app{\lem[nolem]{\skp{pāda d}}
	\rdg[wit={V19,Jyo}]{\om}}%
nādopāsanam eva  % nadipā°? J7ac
	\app{\lem[wit={J5,Gamma,V3}]{ca}
		\rdg[wit={Epsi,P11,C6}]{vā}% +F
		}//\versenr} % cā N24
	\orgvnr{34}\\!}
\end{tlg}
%</vs34>
%\commcite%\newpage

%<*vs35>
\begin{tlg}[hp04_035]
\tl{\app{\lem[nolem]{}
	\rdg[wit={N3}]{\unavbl}
	\rdg[wit={E2,Zeta2,J10,Jyo}]{\om}
	}%
\pada{\app{\lem[nolem]{\skp{pāda a}}
	\rdg[wit={V19}]{\om}}%
\app{\lem[wit={N23,G5,C6}]{saukhya}
		\rdg[wit={J5}]{sākṣaṃ}
		\rdg[wit={J7}]{saukhyā}
		\rdg[wit={G11}]{sāṃkhyaṃ}
		\rdg[wit={P11,V3}]{sāṃkhya}
		}%
	\app{\lem[wit={J7,Epsi,C6,V3}]{praviṣṭa}
		\rdg[wit={J5}]{pravṛṣṭa}
		\rdg[wit={N23,P11}]{pratiṣṭha}
		}cittānāṃ} % cittānyu N23
\pada{\app{\lem[nolem]{\skp{pāda b}}
	\rdg[wit={V19}]{\om}}%
	mūḍhānām api saṃmatam/} % dṛḍhānām! G11; avivekinām G5, saṃmataḥ P11
	\\+}
\tl{
\pada{\app{\lem[wit={J5,Delta2,Epsi,Pi}]{sadya}% +K3; sadyaivānanda G5
		\rdg[wit={Gamma}]{satyam}}%
	\app{\lem[wit={J5,Gamma,C7,Epsi,Pi}]{ānanda}% +K3
		\rdg[wit={V19}]{ādāya}}%
	\app{\lem[wit={G4,J7,Delta2}]{saṃdhāyī}
		\rdg[wit={J5}]{saṃdāï}
		\rdg[wit={N23}]{saṃdhyāyī}
		\rdg[wit={Epsi,P11}]{saṃdāyī}% +K3 % saṃ√dā ist nicht gebräuchlich.
		\rdg[wit={C6}]{saṃdāyi}
		\rdg[wit={V3}]{sadāyī}
		}}
\pada{\app{\lem[wit={ceteri}]{jāyate}
		\rdg[wit={V19}]{jāvate}}
	\app{\lem[wit={G4,Gamma,Delta2,Epsi,C6,V3}]{nādajo}% °jau N23; +K3
		\rdg[wit={J5}]{natato}
		\rdg[wit={P11}]{nādato}
		} layaḥ//\versenr} % layā G5, laya V3
	\label{saukhya}
	\orgvnr{35}\\!}
\end{tlg}
%</vs35>
\commciterange{35}{34–35}\newpage


\iffalse
\startaltrecension
%<*vs35-1>
\begin{alttlg}[hp04_035_1]
\tl{\app{\lem[nolem]{}
	\rdg{\texteng{=\,3.48}}
	\rdg[wit={Epsi,V15,J10}]{\only}}%
\pada{ekaṃ sṛṣṭimayaṃ bīja}% svasti° G11
\pada{m ekā mudrā 
\app{\lem[wit={Epsi,J10}]{ca}
	\rdg[wit={V15}]{tu}% +J11
	} khecarī/}\\+}
\tl{
\pada{eko devo \app{\lem[wit={V15,J10}]{nirālamba}% +J11
	\rdg[wit={Epsi}]{nirālambo hy}}}
\pada{ekāvasthā manonmanī//\versenr} % °mani V15
\\!} % Not in V3,G7, but in G11
\end{alttlg}
%</vs35-1>

%<*vs35-2>
\begin{alttlg}[hp04_035_2]
\tl{\app{\lem[nolem]{}
	\rdg[wit={N19}]{\NotIn}}%
\pada{śaṅkhadundubhi%
	\app{\lem[wit={G5,V15,J10,P11,Jyo}]{nādaṃ ca}
		\rdg[wit={G11,V3}]{nādaś ca}
		\rdg[wit={C6}]{nādāṃś ca}}}
\pada{na śṛṇoti kadācana/}\\+} % sṛṇoti V3,J10; canaḥ J11, canā G11
\tl{
\pada{\app{\lem[wit={G5,V15,J10,Jyo}]{kāṣṭhavaj jāyate}
		\rdg[wit={G11}]{kāṣṭhavaj jñāyate}
		\rdg[wit={P11}]{sthāṇuvarddhattayed}% °yej jogī P11
		\rdg[wit={C6}]{sthāṇuvad vartate}% =M1,F,G3
		\rdg[wit={V3}]{sthāṇu vardhate}}
	\app{\lem[wit={J10,Jyo}]{deha}% G7 deha-m?
		\rdg[wit={Epsi}]{nādam}
		\rdg[wit={V15}]{dehe}
		\rdg[wit={Pi}]{yogī hy}}} % =M1,G3
\pada{unmanyā\app{\lem[wit={Epsi,V15,Pi,Jyo},alt={°vasthayā}]{\skp{°}vasthayā} % unmanyava° would be unmetrical. V15 has avagraha: unmanyā'va°
		\rdg[wit={J10}]{vasthāyāṃ}} dhruvam//\versenr} 
		\\!}
\end{alttlg}
%</vs35-2>

%<*vs35-3>
\begin{alttlg}[hp04_035_3]
\tl{\app{\lem[nolem]{}
	\rdg[wit={N19}]{\NotIn}}%
\pada{sarvāvasthāvinirmuktaḥ} % mukto G5
\pada{sarvacintā%
	\app{\lem[wit={Epsi,V15,J10,P11,C6,Jyo}]{vivarjitaḥ} % jita P11
		\rdg[wit={V3}]{vivarjitaṃ}}/}\\+}
\tl{
\pada{\app{\lem[wit={Epsi,V15,J10,Jyo},alt={mṛtavat}]{mṛtava\skp{t}}
		\rdg[wit={Pi}]{kāṣṭhavat}}% =M1,F
	\app{\lem[wit={G5,V15,J10,P11,C6,Jyo},alt={tiṣṭhate}]{\skm{t }tiṣṭhate}
		\rdg[wit={G11}]{vartate}
		\rdg[wit={V3}]{tiṣṭhayed}
		} yogī}
\pada{sa mukto nātra saṃśayaḥ//\versenr}
\\!}
\end{alttlg}
%</vs35-3>


%<*vs35-4>
\begin{alttlg}[hp04_035_4]
\tl{\app{\lem[nolem]{}
	\rdg[wit={N19,J10}]{\NotIn}}%
\pada{na \app{\lem[wit={Epsi,P11}]{hi jānāti}% +F
		\rdg[wit={V15,Jyo}]{vijānāti}
		\rdg[wit={V3}]{hi jānaṃti}} śītoṣṇaṃ}
\pada{\app{\lem[wit={G11,V15,P11,Jyo}]{na duḥkhaṃ na sukhaṃ}
		\rdg[wit={G5}]{na duḥkhaṃ sukham eva vā}
		\rdg[wit={V3}]{na ca duḥkhaṃ sukhaṃ}} tathā/}\\+}
\tl{
\pada{\app{\lem[wit={V15,Jyo}]{na mānaṃ nāpamānaṃ}
		\rdg[wit={Epsi}]{na mānaṃ nāvamānaṃ}% +F
		\rdg[wit={P11}]{na mānaṃ cāpamānaṃ}
		\rdg[wit={V3}]{na ca mānāpamānaṃ}} ca}
\pada{yogī \app{\lem[wit={G11,P11,C6,Jyo}]{yuktaḥ}
		\rdg[wit={G5,V15}]{muktaḥ}
		\rdg[wit={V3}]{yukti}} samādhinā//\versenr}%
\myfn{In \getsiglum{Jyo}, this verse appears after \manuref{X4.131}. Instead, the following verse has been inserted here, which has no counterpart in the other manuscripts of the \HP:
\devnote{khādyate na ca kālena bādhyate na ca karmaṇā/
sādhyate na sa kenāpi yogī yuktaḥ samādhinā//} (4.108)}
\\!}
\end{alttlg}
%</vs35-4>

%<*vs35-5>
\begin{alttlg}[hp04_035_5]
\tl{\app{\lem[nolem]{}
	\rdg[wit={Epsi,V3,Jyo}]{\only}}%
\pada{na gandhaṃ na rasaṃ rūpaṃ} % nāgandhaṃ G11, raso G5
\pada{\app{\lem[resp=emend]{na sparśanaṃ na ca śrutam}
	\rdg[wit={Epsi,V3}]{na\,(\om\ \getsiglum{V3}) sparśaṃ na ca na śrutaṃ}% +G3; niśrutaṃ F
	\rdg[wit={Jyo}]{na ca sparśaṃ na niḥsvanam}% = VM
	}/}\\+}
\tl{
\pada{nātmānaṃ
	\app{\lem[wit={Epsi,Jyo}]{na paraṃ}
	\rdg[wit={V3}]{paramaṃ}}
	vetti}
\pada{yogī
	\app{\lem[wit={G11,Jyo}]{yuktaḥ}
	\rdg[wit={G5}]{muktaḥ}
	\rdg[wit={V3}]{yukti}
	} samādhinā//\versenr} 
	\\!}
\end{alttlg}
%</vs35-5>
%<*vs35-6>
\begin{alttlg}[hp04_035_6]
\tl{\app{\lem[nolem]{}
	\rdg[wit={Epsi,V15,J10,Jyo}]{\only}}%
\pada{\app{\lem[resp=emend]{avedhyaḥ}
	\rdg[wit={G11}]{adhyāpyāḥ}% avāpyaḥ M3, avyāpya G7; G3 does not have this verse?
	\rdg[wit={G5}]{adhyāpaḥ}
	\rdg[wit={V15,J10,Jyo}]{avadhyaḥ}% abodhya P6
	} 
	sarva\app{\lem[wit={V15,J10,Jyo},alt={śastrāṇām}]{śastrāṇā\skp{m}}
	\rdg[wit={Epsi}]{śāstrāṇām}
	}}%m
\pada{\app{\lem[wit={Epsi,V15,J10},alt={avadhyaḥ}]{\skm{m }avadhyaḥ}% ā(!)bādhyaḥ P6
	\rdg[wit={Jyo}]{aśakyaḥ}} 
	sarva\app{\lem[wit={ceteri}]{dehinām}
	\rdg[wit={G11}]{\om}}/}\\+}
\tl{
\pada{\app{\lem[wit={G5,V15,J10,Jyo}]{agrāhyo}
	\rdg[wit={G11}]{\om}} 
\app{\lem[wit={V15,J10}]{mantratantrāṇāṃ} % tatrāṇāṃ V15
	\rdg[wit={G11}]{\om}
	\rdg[wit={G5,Jyo}]{mantrayantrāṇāṃ}% +M3; yantramantrāṇāṃ G7
	}}
\pada{yogī \app{\lem[wit={J10,Jyo}]{yuktaḥ}% M3
	\rdg[wit={G11}]{\om}
	\rdg[wit={G5,V15}]{muktaḥ}
	} 
	\app{\lem[wit={G5,V15,J10,Jyo}]{samādhinā}
	\rdg[wit={G11}]{mādhinā}}//\versenr}
	\\!} % NotIn{V3}
\end{alttlg}
%</vs35-6>


%<*vs35-7>
\begin{alttlg}[hp04_035_7]
\tl{\app{\lem[nolem]{}
	\rdg[wit={N19}]{\NotIn}}%
\pada{\app{\lem[nolem]{\skp{pāda a}}
	\rdg[wit={G5}]{\om}}%
cittaṃ na suptaṃ no jāgrat} % nna P11; jāgrati J10ac
\pada{\app{\lem[nolem]{\skp{pāda b}}
	\rdg[wit={G5}]{\om}}%
\app{\lem[wit={G11}]{smṛtiman na}% +F; smṛtaṃ na ca nānyathā M3
		\rdg[wit={V15}]{smṛtivarṇaṃ}
		\rdg[wit={J10}]{spṛśati vastu}
		\rdg[wit={P11}]{na smṛtir na}
		\rdg[wit={C6}]{smṛtyaman}
		\rdg[wit={V3}]{sṛtinannaṃ}
		\rdg[wit={Jyo}]{smṛtivismṛ}
		}
	\app{\lem[resp=emend]{na cānyathā}
		\rdg[wit={G11,V15,J10,P11,V3}]{ca nānyathā}
		\rdg[wit={C6}]{na nānyathā}
		\rdg[wit={Jyo}]{tivarjitam}}/}
		\\+}
\tl{
\pada{\app{\lem[wit={Epsi,V15,Pi}]{nāstam eti}
		\rdg[wit={J10}]{na vāstum eti}
		\rdg[wit={Jyo}]{na cāstam eti}}
	\app{\lem[wit={Epsi,V15,J10,P11,C6}]{na codeti}
		\rdg[wit={V3}]{na cādeti}
		\rdg[wit={Jyo}]{nodeti}}}
\pada{\app{\lem[wit={Epsi,V15,P11,C6,Jyo}]{yasyāsau}
		\rdg[wit={J10}]{yathāsau} % mukti J10ac
		\rdg[wit={V3}]{\illeg}}
	\app{\lem[wit={Epsi,V15,J10,P11,C6,Jyo}]{mukta eva saḥ} % mukti J10ac
		\rdg[wit={V3}]{\illeg}}//\versenr} 
		\label{nasuptam}\\!} % yavaprasa in J8
\end{alttlg}
%</vs35-7>

%<*vs35-8>
\begin{alttlg}[hp04_035_8]
\tl{\app{\lem[nolem]{}
	\rdg[wit={N19,J10}]{\NotIn}}%
\pada{\app{\lem[wit={G11,V3,Jyo}]{svastho}% +M3?
		\rdg[wit={G5}]{svapne}
		\rdg[wit={V15}]{svecchā}
		\rdg[wit={P11}]{svapno}
		\rdg[wit={C6}]{supto}
		} jāgradavasthāyāṃ} % °sthīyāṃ P11; ṃ om. V3
\pada{\app{\lem[wit={G11,P11,V3,Jyo}]{suptavad yo}
		\rdg[wit={G5}]{pūrvavad yo}
		\rdg[wit={V15}]{suptaḥ sadyo}
		\rdg[wit={C6}]{suptavadhyo}
		}%
	\app{\lem[wit={Epsi,V15,V3,Jyo}]{'vatiṣṭhate}
		\rdg[wit={P11,C6}]{vatiṣṭhati}}/}\\+} % °tiṣṭhati V15pc?
\tl{
\pada{\app{\lem[wit={V15,Jyo}]{niḥśvāsocchvāsa}
		\rdg[wit={Epsi}]{niśvāsocchvāsa}% +F
		\rdg[wit={P11}]{nisvāsośvaḥsa}
		\rdg[wit={C6}]{niḥśvāsaśvāsa}
		\rdg[wit={V3}]{niśvāsośvāsa}
		}%
	\app{\lem[wit={V15,V3,Jyo}]{hīnaś ca}% source & testimonia
		\rdg[wit={G11,P11,C6}]{hīnas tu}% +F
		\rdg[wit={G5}]{hīnasya}
		}} 
\pada{\app{\lem[wit={Epsi,V15,Jyo}]{niścitaṃ}
		\rdg[wit={P11}]{niścitto}
		\rdg[wit={C6}]{niśceṣṭo}
		\rdg[wit={V3}]{niścito}
		} mukta eva saḥ//\versenr}% sa V3
		\\!}
\end{alttlg}
%</vs35-8>
\endaltrecension
\fi


%==================================
%<*vs36>
\begin{tlg}[hp04_036]
\tl{
\pada{\app{\lem[nolem]{\skp{pāda a}}
	\rdg[wit={N3}]{\unavbl}}%
nādānusandhānasamādhibhājāṃ}\\+} % dānā° V15, saṃdhāta P11, samādhinā J5
\tl{
\pada{\app{\lem[nolem]{\skp{pāda b}}
	\rdg[wit={N3}]{\unavbl}}%
yogīśvarāṇāṃ % <<yo>>gośvarāṇāṃ J7, °svarā° N23
	\app{\lem[wit={J5,J7,Delta,Epsi,V15,C6,V3}]{hṛdaye prarūḍham}% prarūḍha V3
		\rdg[wit={N23,P11}]{hṛdayaprarūḍhaṃ}% [rū]ḍhaṃ N23, rūṇāṃ P11; +F
		\rdg[wit={N19,J10,Jyo}]{hṛdi vardhamānaṃ}}/}\\+}  % vaddha° N19, vaddhra°? N26
\tl{
\pada{\app{\lem[nolem]{\skp{pāda c}}
	\rdg[wit={N3}]{\unavbl}
	\rdg[wit={J5}]{\om}}%
ānandam ekaṃ vacasā%m
	\app{\lem[wit={ceteri},alt={avācyaṃ}]{\skm{m }avācyaṃ} % yacasām V3
%		\rdg[wit={G11}]{avācaṃ}
		\rdg[wit={N19}]{avākyaṃ}
		\rdg[wit={C6,Jyo}]{agamyaṃ}% +G3
		}}\\+}
\tl{
\pada{\app{\lem[nolem]{\skp{pāda d}}
	\rdg[wit={J5}]{\om}}%
\app{\lem[wit={ceteri}]{jānāti}
		\rdg[wit={N3}]{\lost}
		\rdg[wit={N19}]{jānaṃti}
		\rdg[wit={P11,C6}]{jānāty a}
		}
	\app{\lem[wit={J7,Epsi,Zeta2,J10,V3,Jyo}]{taṃ śrī}
		\rdg[wit={N3}]{\lost}
		\rdg[wit={N23}]{tatvaṃ śrī}
		\rdg[wit={Delta}]{tattvaṃ}
		\rdg[wit={P11}]{tītaṃ}
		\rdg[wit={C6}]{taḥ śrī}
		}%
	\app{\lem[wit={ceteri}]{gurunātha}
		\rdg[wit={N3}]{+\,+\,nātha}
		\rdg[wit={Delta}]{guṇanātha}
		}
	\app{\lem[wit={N3,J7,Delta,G11,V15,Pi}]{eva}
		\rdg[wit={N23,G5}]{evaṃ}
		\rdg[wit={N19,Jyo}]{ekaḥ}% +F
		\rdg[wit={J10}]{ekaṃ}}//\versenr}
	\orgvnr{36}\\!} % N3 resumes with nātha eva
\end{tlg}
%</vs36>
\commcite%\newpage


%<*vs37>
\begin{tlg}[hp04_037]
\tl{%\app{\lem[nolem]{}
	%\rdg[wit={G11,Zeta2,J10},alt=\textapp{found after \ref{sarvacinta}}]{\skp{found after \manuref{4.18}}}%
	%\rdg[wit={Jyo},alt=\textapp{found \expnr{???}}]{\skp{found ???}}
	%\rdg[wit={Pi},alt=\textapp{transposed with the next one}]{\skp{transposed with the next one}}}%
\pada{sarvacintāṃ parityajya} % ciya J5, ciṃtā C6,V3, ciṃtāḥ J10; samutsṛjya G5
\pada{\app{\lem[wit={ceteri}]{sāvadhānena}
	\rdg[wit={N19,J10}]{sarvadānena}}
	cetasā/}\\+}
\tl{
\pada{\app{\lem[wit={ceteri}]{nāda evānusandheyo}% nāda yecānusaṃdhyeyo J5
	\rdg[wit={V19,J10}]{nādam evānusaṃdhatte}% +K3
	\rdg[wit={G5}]{nādam evānubandheyo}
	\rdg[wit={N19}]{nādam evānusaṃdhe}% (yo omitted)
	}}
\pada{yoga\app{\lem[wit={ceteri},alt={sāmrājyam}]{sāmrājya\skp{m}}
	\rdg[wit={V19}]{samrājyam}
	\rdg[wit={C6}]{samrājam}
	}%
	\app{\lem[wit={N3,J7,Delta,G11,V15,Pi,Jyo},alt={icchatā}]{\skm{m }icchatā}
		\rdg[wit={J5}]{iṣṭatā}
		\rdg[wit={G4,G5,N19}]{icchatāṃ}
		\rdg[wit={N23,J10}]{icchati}
		}//\versenr}%
		\label{sarvacinta2}
		\orgvnr{37}\\!}
\end{tlg}
%</vs37>
\commcite\newpage


%<*vs38>
\begin{tlg}[hp04_038]
\tl{\app{\lem[nolem]{}
	\rdg[wit={J10}]{\om}}%
\pada{\app{\lem[wit={ceteri}]{karṇau}
		\rdg[wit={N3,N23}]{karṇo}
		\rdg[wit={G4}]{karṇā}
		\rdg[wit={P11}]{karṇa}
		}
	\app{\lem[wit={ceteri}]{pidhāya}% N3,Gamma,E2,G11,N19,V15,Pi,Jyo
		\rdg[wit={J5}]{nidhāya}
		\rdg[wit={G4}]{pidhāna}
		\rdg[wit={V19}]{pi}
		}
	\app{\lem[wit={G4,G5,N19}]{tūlena}% 
		\rdg[wit={N3,J5,G11,V3}]{mūlena}% +F,G3
		\rdg[wit={Gamma}]{hastena}
		\rdg[wit={V19}]{hastābhya[ṃ]}
		\rdg[wit={E2,C6,Jyo}]{hastābhyāṃ}
		\rdg[wit={V15}]{śū\,\_\,na}
		\rdg[wit={P11}]{tulyena}
		}}
\pada{\app{\lem[wit={N3,J5,G11,Zeta2,Jyo}]{yaṃ}
		\rdg[wit={G4,Gamma,Delta,C6}]{yaḥ}% +F
		\rdg[wit={G5,P11}]{saṃ}
		\rdg[wit={V3}]{sa}} śṛṇoti % °tiṃ P11, śruṇoti V3
	\app{\lem[wit={N3,J5,Delta,Epsi,Zeta2,Pi,Jyo}]{dhvaniṃ muniḥ} % dhvani V3, ddhvaniṃ N19; muni P11,V3
		\rdg[wit={N23}]{dhvaniṃ muniṃ}
		\rdg[wit={J7}]{munir dhvanim}
%		\rdg[wit={C7}]{dhvaniṃ dhvaniḥ}
		}/}\\+}
\tl{
\pada{\app{\lem[wit={ceteri}]{tatra cittaṃ} % citta N23
		\rdg[wit={J5,P11}]{tatra ciṃtāṃ} % ṃ oṃ J5
		\rdg[wit={G5}]{cittavṛtti}
		}
	\app{\lem[wit={N3,J5,G5,Pi,Jyo}]{sthirī}
		\rdg[wit={Gamma,Delta,Zeta2}]{sthiraṃ}
		\rdg[wit={G11}]{sthitaṃ}
		}kuryā}%
\pada{d yāva%t  % yāva Gamma,J5,V3
	\app{\lem[wit={ceteri},alt={sthirapadaṃ}]{\skm{t }sthirapadaṃ}% sira° P11
		\rdg[wit={V3}]{sthiparamaṃ}}
	\app{\lem[wit={ceteri}]{vrajet}% vraje N23
		\rdg[wit={Zeta2}]{bhavet}}//\versenr}
		\orgvnr{38}\\!}
\end{tlg}
%</vs38>
\commcite\newpage


%========= common verses ==================
%<*vs39>
\begin{tlg}[hp04_039]
\tl{
\pada{abhyasyamāno %  abhyāsya° N23, °mano N3
	\app{\lem[wit={ceteri}]{nādo}% nādau J5
		\rdg[wit={N23}]{nātho}}%
	\app{\lem[wit={ceteri}]{'yaṃ}% +G5, ya G11
		\rdg[wit={C6}]{yo}}}
\pada{\app{\lem[wit={J7,Epsi,C6,Jyo}]{bāhyam āvṛṇute}
		\rdg[wit={N3}]{bāhyam āśṛṇu}
		\rdg[wit={J5}]{bāhyaṃ ca śṛṇute}
		\rdg[wit={N23}]{bāhyanā\,\_\,ṇute}
		\rdg[wit={V19,V15}]{bāhyam āvartaye}% °yed +K3,C7
		\rdg[wit={E2}]{bāhyād āvartayed}
		\rdg[wit={N19}]{bāhyamānaśṛṇvate}
		\rdg[wit={J10}]{cānyam āśṛṇute}
		\rdg[wit={P11}]{bāhyanāvṛṇute}
		\rdg[wit={V3}]{bāhyam āsṛṇate}
		}\marmas 
	\app{\lem[wit={N3,J7,Delta,V15,J10,Jyo}]{dhvanim}
		\rdg[wit={J5}]{dhvaniṃḥ}
		\rdg[wit={N23}]{dhvani}
		\rdg[wit={Epsi,N19,Pi}]{dhvaniḥ}% ddhaniḥ N19; +F
		}/}\\+}
\tl{
\pada{\app{\lem[wit={ceteri},alt={pakṣād}]{pakṣā\skp{d}}
		\rdg[wit={G4,V19,Epsi,J10}]{paścād}}%
	\app{\lem[wit={N3,J5,J7,E2,J10,V3,Jyo},alt={vikṣepam akhilaṃ}]{\skm{d }vikṣepam akhilaṃ}% +C7
		\rdg[wit={G4}]{vikṣiptam a[nila]ṃ}
		\rdg[wit={N23}]{vikṣeyam akhilaṃ}
		\rdg[wit={V19}]{vikṣepam atulaṃ}
		\rdg[wit={Epsi}]{vikṣiptam akhilaṃ}
		\rdg[wit={Zeta2}]{vipakṣam akhilaṃ}
		\rdg[wit={P11}]{vikṣyemanilaṃ}
		\rdg[wit={C6}]{vipakṣayed enaṃ}}}
\pada{\app{\lem[wit={ceteri}]{jitvā} % jītvā P7
		\rdg[wit={G5}]{hitvā}
		\rdg[wit={J10}]{jīvo}
		} yogī sukhī bhavet//\versenr}
		\orgvnr{39}\\!}
\end{tlg}
%</vs39>
\commcite%\newpage


%<*vs40>
\begin{tlg}[hp04_040]
\tl{
\pada{\app{\lem[wit={ceteri}]{śrūyate}% śūyate N3
		\rdg[wit={E2}]{jāyate}% +C7
		}
	\app{\lem[wit={ceteri}]{prathamābhyāse}
		\rdg[wit={N3}]{prathamābhyāso}
		\rdg[wit={V19}]{prathame bhyāse}
		}}
\pada{nādo nānāvidho mahān/}\\+} % nādau J5, nādā N23; vidhā N19; mahāt N19 
\tl{
\pada{\app{\lem[wit={ceteri},alt={vardhamāne tato 'bhyāse}]{vardhamāne tato'bhyāse}
		\rdg[wit={V15,Jyo}]{tato 'bhyāse vardhamāne}
		}}
\pada{śrūyate % srū˟yate N3, śrūyete P11
	\app{\lem[wit={N3,J5,Delta,G11,J10,C6,V3}]{sūkṣmasūkṣmataḥ}
		\rdg[wit={N23}]{sūjyasūjyakaḥ}
		\rdg[wit={J7,G5,V15,Jyo}]{sūkṣmasūkṣmakaḥ} % 
		\rdg[wit={N19,P11}]{sūkṣmataḥ}
		}//\versenr}
		\orgvnr{40}\\!} % śū° P11; ḥ om. N19; haplogr.
\end{tlg}
%</vs40>
\commcite\newpage


%<*vs41>
\begin{tlg}[hp04_041]
\tl{
\pada{ādau jaladhi% adau P7
	\app{\lem[wit={ceteri}]{jīmūta}% jāmūta J5
		\rdg[wit={N23,P11,V3}]{jīmūte}}}% °te J10pc
\pada{bherī% bhīrī K3
	\app{\lem[wit={Epsi,Zeta2,J10,P11}]{nirjhara}
		\rdg[wit={N3}]{rsara} % unm.
		\rdg[wit={J5}]{nisara}
		\rdg[wit={N23}]{sarāva}
		\rdg[wit={J7}]{śabdatu}
		\rdg[wit={V19}]{nirjara}
		\rdg[wit={E2}]{bhūrbhūra}% +C7
		\rdg[wit={C6}]{nigama}
		\rdg[wit={V3}]{nirbhara} % +P7
		\rdg[wit={Jyo}]{jharjhara}
		}%
	\app{\lem[wit={N19,C6,Jyo}]{saṃbhavāḥ}
		\rdg[wit={N3,J5,P11}]{saṃbhavā}
		\rdg[wit={Gamma,Delta,Epsi,V15}]{saṃbhavaḥ}
		\rdg[wit={J10,V3}]{nisvanaḥ}}/}\\+} % niśvanaḥ
\tl{
\pada{madhye
	\app{\lem[wit={ceteri}]{mardala}% mardvala N23
		\rdg[wit={E2}]{mandala}% +K3,C7
		\rdg[wit={Epsi}]{maddala}
		}%
	\app{\lem[wit={N3,J5,G11,Zeta2,Jyo}]{śaṃkhotthā}% śaṃdho J5; °otchā N3
		\rdg[wit={G4}]{śaṃkhoddhāḥ}
		\rdg[wit={Gamma}]{śaṃkhotha}
		\rdg[wit={Delta,G5,J10,P11,C6pc,V3}]{śaṃkhottha}% saṃkho V3, śaṃ<<kho>>tha C6; °occha P11, °okta G5
		\rdg[wit={C6ac}]{śaṅkhottho}% +K3
		}\marma}
\pada{ghaṇṭā\app{\lem[wit={J5,G4,J7,Epsi,Zeta2,C6,V3,Jyo}]{kāhala}
		\rdg[wit={N3,P11}]{kāhāla}
		\rdg[wit={N23}]{kāhla}
		\rdg[wit={Delta}]{kalaha}
		\rdg[wit={J10}]{kolāha}
		}%
	\app{\lem[wit={N3,J5,Pi,Jyo},alt={°jās}]{\skp{°}jā\skp{s}}
		\rdg[wit={G4,Zeta2}]{kās}
		\rdg[wit={Gamma,Delta,Epsi}]{jas}
		\rdg[wit={J10}]{las}
		}%
	\app{\lem[wit={ceteri},alt={tathā}]{\skm{s }tathā} % tatha J10
		\rdg[wit={C6}]{tataḥ}}//\versenr}
		\orgvnr{41}\\!}
\end{tlg}
%</vs41>
\commcite%\newpage


%<*vs42>
\begin{tlg}[hp04_042]
\tl{
\pada{\app{\lem[wit={ceteri}]{ante}% Alpha,J7,Delta,G11,V15,Pi,Jyo
		\rdg[wit={N23}]{avai}
		\rdg[wit={N19,J10}]{anye}
		} tu kiṅkiṇī% kikiṇī N3, °nī V19, kiṃkaṇī J5
	\app{\lem[wit={N3,Epsi,Zeta2,J10,Jyo}]{vaṃśa}% vaṃśaṃ
		\rdg[wit={J5}]{śabda}
		\rdg[wit={G4}]{bṛṃdā}
		\rdg[wit={Gamma,Delta,C6,V3}]{vṛnda}% vṛṃda -> <<śa>>bdaṃva?! J7
		\rdg[wit={P11}]{vaṃda}
		}}%
\pada{\app{\lem[wit={ceteri}]{vīṇā}% Alpha,Gamma,Delta,G11,J10,Pi,Jyo
		\rdg[wit={Zeta2}]{nādā}
		}bhramara% bhrumara N23, bhrasara N19
	\app{\lem[wit={N3,G4,Epsi,N19,C6}]{nisvanāḥ}
		\rdg[wit={J5}]{niḥśvanā}
		\rdg[wit={N23,E2,P11}]{niḥsvanaḥ}% nissvanaḥ K3
		\rdg[wit={J7,V19}]{nisvanaḥ} % niśvanaḥ V19
		\rdg[wit={V15,Jyo}]{niḥsvanāḥ}
		\rdg[wit={J10,V3}]{nisvanā} % nisvānā J10
		}/}\\+}
\tl{
\pada{iti \app{\lem[wit={N3,J5,Epsi,Zeta2,J10,P11,C6,Jyo}]{nānāvidhā}
		\rdg[wit={Gamma,Delta,V3}]{nānāvidho}}
	\app{\lem[wit={N3,J10,C6,Jyo}]{nādāḥ}
		\rdg[wit={J5,Epsi,V15,P11,V3}]{nādā}
		\rdg[wit={N23}]{nādaṃ}
		\rdg[wit={J7,Delta}]{nādaḥ}
		\rdg[wit={N19}]{vādāḥ}}}
\pada{\app{\lem[wit={J5,Epsi,V15,J10,P11,C6,Jyo}]{śrūyante}
		\rdg[wit={N3,Gamma,Delta,N19,V3}]{śrūyate}} % śṛyate N3
	\app{\lem[wit={ceteri}]{deha}% N3,J5,Gamma,Delta,G11,Pi,Jyo
		\rdg[wit={N19,J10}]{yatra}
		\rdg[wit={V15}]{tatra}}%
	\app{\lem[wit={N3,J5,Epsi,Zeta2,J10,P11,V3}]{madhyataḥ}
		\rdg[wit={Gamma,Delta}]{madhyagaḥ}
		\rdg[wit={C6,Jyo}]{madhyagāḥ}
		}//\versenr}
		\orgvnr{42}\\!}
\end{tlg}
%</vs42>
\commcite\newpage


%<*vs43>
\begin{tlg}[hp04_043]
\tl{
\pada{\app{\lem[wit={ceteri}]{mahati}
		\rdg[wit={J5}]{mahatiḥ}
		\rdg[wit={G5,V15}]{mahatī}
		\rdg[wit={C6}]{\om}}
	\app{\lem[wit={ceteri},alt={śrūyamāṇe/-māne}]{śrūyamāṇe}% °ṇo J5, °ne J10
		\rdg[wit={N23}]{{[ṇya]}yatamāne}}%
	\app{\lem[wit={ceteri}]{'pi}
		\rdg[wit={Gamma}]{ti}
		\rdg[wit={C6}]{pi nāde vai}}} % +P7
\pada{meghabhe%ry
	\app{\lem[wit={J5,Gamma,G5,N19,J10},alt={ādikadhvanau}]{\skm{ry}ādikadhvanau}% ddhūnau N19
		\rdg[wit={N3}]{ādidaṃ dhvanau}
		\rdg[wit={Delta,C6,V3,Jyo}]{ādike dhvanau}
		\rdg[wit={G11}]{ākadhvanau}
		\rdg[wit={V15}]{ādike svane}
		\rdg[wit={P11}]{ādike dhṛti}
		}/}\\+}
\tl{
\pada{\app{\lem[wit={ceteri}]{tatra}% N3,J5,G11,N19,V15,J10,Pi,Jyo
		\rdg[wit={Gamma,Delta}]{tataḥ}}
	\app{\lem[wit={ceteri},alt={sūkṣmāt}]{sūkṣmā\skp{t}}% sūkṣmā<<t>> J7
		\rdg[wit={J5,N19}]{sūkṣmā}
		\rdg[wit={P11}]{sūkṣmāṃ}
		\rdg[wit={J10}]{sūkṣmaṃ}
		}%
	\app{\lem[wit={ceteri},alt={sūkṣmataraṃ}]{\skm{t }sūkṣmataraṃ}% °tara N23, paraṃ G5
%		\rdg[wit={C7}]{sūkṣmatamaṃ}
		\rdg[wit={J10}]{nādam eva}
		\rdg[wit={P11}]{taraṃ nādaṃ}
		}}
\pada{\app{\lem[wit={ceteri}]{nādam eva}
		\rdg[wit={J7}]{nādam evaṃ}
		\rdg[wit={J10}]{paritopi}}
	\app{\lem[wit={ceteri}]{parāmṛśet} % pasamṛśet N23
		\rdg[wit={J5}]{parāmṛśaṃ}
		\rdg[wit={J7}]{samabhyaset}
		\rdg[wit={V19}]{parāmṛṣet} % pasamṛṣet V19ac
		}//\versenr}
		\orgvnr{43}\\!}
\end{tlg}
%</vs43>
\commcite%\newpage


%<*vs44>
\begin{tlg}[hp04_044]
\tl{\app{\lem[nolem]{}
	\rdg[wit={E2}]{\om}}%
\pada{\app{\lem[wit={ceteri},alt={ghanam}]{ghana\skp{m}}
		\rdg[wit={J10}]{dhvanam}
		}m utsṛjya vā % yā G5
	\app{\lem[wit={N3,Epsi,Zeta2,J10,Pi,Jyo}]{sūkṣme}
		\rdg[wit={J5,G4,Gamma,V19}]{sūkṣmaṃ}
		}} % sūkṣmo P7
\pada{sūkṣmam utsṛjya vā % sūkṣmasūtsṛjya N23; śū° P11
	\app{\lem[wit={Alpha,Epsi,Zeta2,P11,C6,Jyo}]{ghane}
		\rdg[wit={Gamma,V19}]{ghanam}% +K3,C7
		\rdg[wit={J10}]{dhune}
		\rdg[wit={V3}]{ghanen}
		}\marma/}\\+}
\tl{
%\pada{\app{\lem[wit={N3,J5,G11,Pi},alt={tau tyaktvā ... syād vā}]{%	dau F; tyaktrā J5, tyaktā P11, tyakvā V3; tu tyaktvā madh. + + + G4
%			tau tyaktvā
%			\app{\lem[wit={J5},alt={madhyame}]{madhyame}% +F
%				\rdg[wit={N3,G11,P11,V3},alt={madhyama}]{madhyama}
%				\rdg[wit={C6},alt={madhyama<<ḥ>>}]{madhyama<<ḥ>>}} 
%			\app{\lem[wit={N3,Pi},alt={syād vā}]{syād vā}% madhyame vāpi F
%				\rdg[wit={Epsi},alt={syādau}]{syādau}
%				\rdg[wit={J5},alt={syātaṃstā}]{syātaṃstā}
%				}}% nested!
%		\rdg[wit={Zeta2}]{ramamāṇam api kṣipraṃ}
%		\rdg[wit={J10,Jyo}]{ramamāṇam api kṣiptaṃ}
%		\rdg[wit={Gamma,V19}]{paraṃ tatraiva nikṣipya}% +K3,C7
%		}}
\pada{% tu tyaktvā madh. + + + G4
	\app{\lem[wit={J5}]{tau tyaktvā madhyame}% dau F, tyaktrā J5
		\rdg[wit={N3,Epsi,P11,V3}]{tau tyaktvā madhyama}% tyaktā P11, tyakvā V3
		\rdg[wit={Gamma,V19}]{paraṃ tatraiva}
		\rdg[wit={Zeta2,J10,Jyo}]{ramamāṇam api}
		\rdg[wit={C6}]{tau tyaktvā madhyama<<ḥ>>} 
			} 
	\app{\lem[wit={N3,Pi}]{syād vā}% madhyame vāpi F
		\rdg[wit={J5}]{syātaṃstā}
		\rdg[wit={Gamma,V19}]{nikṣipya}% +K3,C7
		\rdg[wit={Epsi}]{syādau}
		\rdg[wit={Zeta2}]{kṣipraṃ}
		\rdg[wit={J10,Jyo}]{kṣiptaṃ}
		}}
\pada{mano % manā? N23, manau V19
	\app{\lem[wit={ceteri}]{nānyatra}
		\rdg[wit={Zeta2,J10}]{nātra pra}}
	\app{\lem[wit={ceteri}]{cālayet} % vālayet N23, c<<ā>>layet J7
		\rdg[wit={J10}]{cālet}
		\rdg[wit={V3}]{cālayan}}//\versenr}
		\orgvnr{44}\\!}
\end{tlg}
%</vs44>
\commcite\newpage


%<*vs45>
\begin{tlg}[hp04_045]
\tl{
\pada{yatra kutrāpi vā nāde}
\pada{\app{\lem[wit={ceteri}]{lagati}
		\rdg[wit={N23}]{lagavi}
		\rdg[wit={J10}]{galati}
		\rdg[wit={P11}]{lagnaṃti}
		}
	\app{\lem[wit={ceteri}]{prathamaṃ}% °ma J5
		\rdg[wit={V19}]{prathame}}
	\app{\lem[wit={ceteri}]{manaḥ}
		\rdg[wit={N23}]{mataḥ}}/}\\+} %
\tl{
\pada{tatraiva
	\app{\lem[wit={N3,Epsi,V15,P11,C6},alt={tat}]{ta\skp{t}}% ta(l.br.)tat V15
		\rdg[wit={J5}]{tā}
		\rdg[wit={N23}]{stu}
		\rdg[wit={J7,Delta,Jyo}]{su}% +F
		\rdg[wit={N19,V3}]{ta}
		\rdg[wit={J10}]{niś}
		}%
	\app{\lem[wit={ceteri},alt={sthirī}]{\skm{t }sthirī}% sthirā J5; +G5
		\rdg[wit={G11}]{sthiro}
		\rdg[wit={N19}]{śarī}
		\rdg[wit={J10}]{calo}
		}%
	\app{\lem[wit={Alpha,Epsi,Zeta2,J10,Pi}]{bhūtvā}
		\rdg[wit={Gamma,Delta}]{kuryāt}
		\rdg[wit={Jyo}]{bhūya}
		}} % kuryā J7
\pada{tena sārdhaṃ vilīyate//\versenr}%
%\myfnx{After this verse, \getsiglum{Jyo} has X4.100, 103, and 102.}
\label{yatrakutrapi}
\orgvnr{45}\\!} % vinīyate N3
\end{tlg}
%</vs45>
\commcite%\newpage



%======== passage 4 (makaranda) ======

%<*vs46>
\begin{tlg}[hp04_046]
\tl{
\pada{makarandaṃ % ṃ om. J5,V15,J10
	\app{\lem[wit={ceteri},alt={piban}]{piba\skp{n}}
		\rdg[wit={J5}]{pived}
		\rdg[wit={N19}]{piven}
		}%
	\app{\lem[wit={Alpha,Epsi,V15,J10,Pi,Jyo},alt={bhṛṅgo}]{\skm{n }bhṛṅgo}% bhraṃgo P11
		\rdg[wit={Gamma,Delta}]{bhṛṅgī}
		\rdg[wit={N19}]{śṛṃgo}
		}}
\pada{\app{\lem[wit={N3,G4,Delta,Epsi,V3},alt={gandhān}]{gandhā\skp{n}}
		\rdg[wit={J5,N23,C6}]{gandha}
		\rdg[wit={J7,Zeta2,J10,Jyo}]{gandhaṃ}
		%\rdg[wit={K3,C7}]{gandhā}
		\rdg[wit={P11}]{gandho}
		}%
	\app{\lem[wit={ceteri},alt={nāpekṣate}]{\skm{n }nāpekṣate}% nā[p]i° G4; °kṣata J5
		\rdg[wit={N23}]{napekṣate}
		\rdg[wit={N19,J10}]{nopekṣate}
		}
	\app{\lem[wit={ceteri}]{yathā}
		\rdg[wit={E2}]{yadā}
		\rdg[wit={N19}]{'nyathā}
		}/}\\+}
\tl{
\pada{\app{\lem[wit={ceteri}]{nādāsaktaṃ}
		\rdg[wit={Gamma,G5}]{nādasaktaṃ}} % śaktaṃ N19
	\app{\lem[wit={ceteri}]{tathā} % +P7
		\rdg[wit={G5}]{tadā}
		\rdg[wit={C6}]{yathā}} cittaṃ}
\pada{viṣayā%n % viṣayā J5,J7, viṣayā<<n>> N23, °yāṃ V3,N19
	\app{\lem[wit={ceteri},alt={na hi}]{\skm{n }na hi}% nāhi P11
		\rdg[wit={V15}]{naiva}
%		\rdg[wit={C7}]{api}
		}
	\app{\lem[wit={N3,G11,N19,Pi,Jyo}]{kāṅkṣate}
		\rdg[wit={J5,Gamma,Delta,G5,V15,J10}]{kāṅkṣati}}//\versenr}
	\label{makaranda1}\orgvnr{46}\\!}
\end{tlg}
%</vs46>
\commcite\newpage



%<*vs47>
\begin{tlg}[hp04_047]
\tl{\pada{\app{\lem[nolem]{\skp{pāda a}}
	\rdg[wit={Gamma,Delta}]{\om}}%
\app{\lem[wit={J5,Epsi,Zeta2,Pi,Jyo}]{baddhaṃ}% bandhaṃ G5
		\rdg[wit={N3}]{baṃdhaṃ}
		\rdg[wit={J10}]{buddhaṃ}
		}
	\app{\lem[wit={N3,J5,Epsi,P11,C6,Jyo}]{vimukta}
		\rdg[wit={N19}]{vimuktaṃ}
		\rdg[wit={V15,J10}]{viyuktaṃ}
		\rdg[wit={V3}]{timukta}
		}cāñcalyaṃ}
\pada{\app{\lem[nolem]{\skp{pāda b}}
	\rdg[wit={Gamma,Delta}]{\om}}%
nāda\app{\lem[wit={N3,J5,Epsi,Zeta2,J10,V3,Jyo}]{gandhaka}
		\rdg[wit={P11}]{gandhāya}
		\rdg[wit={C6}]{gandhena}
		}%
	\app{\lem[wit={N3,J5,Epsi,V15,C6,V3,Jyo}]{jāraṇāt}
		\rdg[wit={N19,J10,P11}]{jīraṇāt}}/} 
		\\+}
\tl{
\pada{\app{\lem[nolem]{\skp{pāda c}}
	\rdg[wit={E2}]{\om}}%
\app{\lem[wit={ceteri}]{manaḥ}% +K3,C7
		\rdg[wit={N23}]{vona}
		\rdg[wit={P11,V3}]{mana}
		}%
	\app{\lem[wit={J5,Epsi,N19,J10,P11,C6,Jyo}]{pāradam āpnoti}
		\rdg[wit={N3}]{pārajam āpnoti}
		\rdg[wit={N23}]{cāvam avāpnoti}
		\rdg[wit={J7,V19}]{pākam avāpnoti}% +K3,C7
		\rdg[wit={V15}]{pārada āpnoti}
		\rdg[wit={V3}]{pāradham āpnoti}
		}}
\pada{\app{\lem[nolem]{\skp{pāda d}}
	\rdg[wit={E2}]{\om}}%
\app{\lem[wit={ceteri}]{nirālambākhya}
		\rdg[wit={J5}]{nirālaṃbaratha}
%		\rdg[wit={C7}]{nirālambākṣa}
		\rdg[wit={P11}]{nirālambāsthya}
		}%
	\app{\lem[wit={P11,V3}]{khoṭatām}
		\rdg[wit={N3,J5,Epsi,C6}]{ghoṭatāṃ}% ṃ om. J5 ##
		\rdg[wit={G4}]{gopitāṃ}
		\rdg[wit={Gamma}]{ghoṭanam}
		\rdg[wit={V19}]{codanaṃ}
		\rdg[wit={N19}]{khoṭatī}
		\rdg[wit={V15}]{khoṭakaṃ}
		\rdg[wit={J10}]{khegataṃ}
		\rdg[wit={Jyo}]{khe 'ṭanam} % +J11pc
		}//\versenr}
		\orgvnr{47}\\!}
\end{tlg}
%</vs47>
\commcite\newpage


\iffalse
\startaltrecension
%<*vs47-1>
\begin{alttlg}[hp04_047_1]
\tl{\app{\lem[nolem]{}
	\rdg[wit={Epsi,Zeta2,J10,Pi,Jyo}]{\incl}}%
\pada{\app{\lem[wit={G11,Zeta2,V3}]{baddhaḥ}
		\rdg[wit={G5,Jyo}]{baddhaṃ}% bandhaṃ G5
		\rdg[wit={J10}]{baddha}
		\rdg[wit={P11}]{baṃdhaḥ}
		\rdg[wit={C6}]{baddhas}
		}
	\app{\lem[wit={Epsi,V3}]{sunādagandhena}
		\rdg[wit={N19}]{sunāde gandhena}
		\rdg[wit={V15}]{suṃdhanādena}
		\rdg[wit={J10}]{svennādagandhena}
		\rdg[wit={P11}]{sunādavānpana}
		\rdg[wit={C6}]{tu nādagandhena}
		\rdg[wit={Jyo}]{tu nādabandhena}
		}}
\pada{\app{\lem[wit={Epsi,Zeta2,J10,Pi}]{sadyaḥ}
		\rdg[wit={Jyo}]{manaḥ}}%
	\app{\lem[wit={Epsi,Zeta2,J10,P11,C6,Jyo}]{saṃtyakta}
		\rdg[wit={V3}]{sa tyakta}}%
	\app{\lem[wit={Epsi,Zeta2,J10,Pi}]{cāpalaḥ}
		\rdg[wit={Jyo}]{cāpalam}}/}\\+}
\tl{
\pada{prayāti
	\app{\lem[wit={G11}]{cetaḥsūtendraḥ}
		\rdg[wit={G5}]{cetaḥśailendra}
		\rdg[wit={N19}]{sūtaś citteṃdra}
		\rdg[wit={V15}]{sūtacittendraḥ}
		\rdg[wit={P11}]{cet sthūlendraḥ}
		\rdg[wit={C6}]{cetaḥsūtrendre}
		\rdg[wit={V3}]{cetaḥsuteṃdra}
		\rdg[wit={J10}]{svataś caikyaṃ iṃdra}
		\rdg[wit={Jyo}]{sutarāṃ sthairyaṃ}
		}}
\pada{\app{\lem[wit={Epsi,Zeta2,P11,C6}]{pakṣachinna}
		\rdg[wit={J10}]{pacchacchinna}
		\rdg[wit={V3}]{\lacuna}
		\rdg[wit={Jyo}]{chinnapakṣaḥ}
		}
	\app{\lem[resp=emend]{iti prathām}% +F,M1
		\rdg[wit={G11}]{iva prathāṃ}
		\rdg[wit={G5}]{iti prathā}% +M3?
		\rdg[wit={N19}]{iva prabhāṃ}
		\rdg[wit={V15}]{ivāprabhuḥ}
		\rdg[wit={J10}]{iva parvataḥ drumāḥ}
		\rdg[wit={P11}]{dṛti pṛthāṃ}
		\rdg[wit={C6}]{\_\,va patham}
		\rdg[wit={V3}]{\lacuna}
		\rdg[wit={Jyo}]{khago yathā}
		}//\versenr}\\!}
\end{alttlg}
%</vs47-1>
% ithi pradhā M3, iti prathāṃ M1, iva prathāṃ G11, iva dyāmaḥ G7
% cf. Rasendracūḍāmaṇi 16.52-54
%{pañcamo grāsaḥ}
%evaṃ ca pañcamo grāsaḥ pradātavyo'ṣṭamāṃśataḥ /
%sa pātrastho'gnisaṃtapto na gacchati kathañcana // Rcūm_16.52 //
%sa pakṣacchinna ity uktaḥ sa mukto'khiladurguṇaiḥ /
%so'yaṃ niṣevitaḥ sūtastrimāsaṃ rājikāmitaḥ // Rcūm_16.53 //
%viḍaṅgatriphalākṣaudraiḥ khe devaiḥ saha saṅgamam /
%ghrāṇamātreṇa sūtendraḥ sarvaroganikṛntanaḥ // Rcūm_16.54 //
%guṇā ete vinirdiṣṭā rasasya rasavādibhiḥ /
%sakalāste guṇāḥ satyā bhairaveṇa prakīrtitāḥ // Rcūm_16.55 //
% cf. also NWS pakṣaccheda
\endaltrecension
\fi


%<*vs48>
\begin{tlg}[hp04_048]
\tl{\app{\lem[nolem]{}
	\rdg[wit={G4,G5}]{\om}
	}% om. G5,M3
\pada{\app{\lem[wit={N3,J7,Delta,V15,P11,C6},alt={nādaśravaṇataś cittam}]{nādaśravaṇataś citta\skp{m}} % cittaṃm V19
		\rdg[wit={J5}]{nādaḥ śravaṇañ [c]ittaṃm}
		\rdg[wit={N23}]{nādaśravaṇaś cittaṃ matam}
		\rdg[wit={G11}]{nadaśravaṇakṛc cittaṃ}
		\rdg[wit={N19}]{nādaḥ śravaṇataś cittam}
		\rdg[wit={J10}]{nādena praṇataṃ cittam}
		\rdg[wit={V3}]{nādaḥ śravaṇataḥś citam}
		\rdg[wit={Jyo}]{nādaśravaṇataḥ kṣipram}
		}}%
\pada{\app{\lem[wit={N3,Gamma,E2,G11,Pi,Jyo},alt={antaraṅga}]{\skm{m }antaraṅga}
		\rdg[wit={J5}]{anataraṃga}
		\rdg[wit={V19}]{aṃtaraṃ sa}
		\rdg[wit={Zeta2}]{aṃtaraṃgaṃ}
		\rdg[wit={J10}]{aṃtaraṃgā}
		}%
	\app{\lem[wit={ceteri}]{bhujaṅgamaḥ} % ṃ om. V3
		\rdg[wit={N23}]{turaṃgavaḥ}
		\rdg[wit={J7,E2}]{turaṅgamaḥ}% +C7
		}/}\\+}
\tl{
\pada{\app{\lem[wit={Gamma,Zeta2,J10,P11,V3,Jyo}]{vismṛtya}
		\rdg[wit={N3,J5,G11,C6}]{saṃsmṛtya}% +G3
		\rdg[wit={Delta}]{viśūnyaṃ}}
	\app{\lem[wit={ceteri},alt={sarvam}]{sarva\skp{m}}
		\rdg[wit={Zeta2,J10}]{viśvam}}%
	\app{\lem[wit={N3,Jyo},alt={ekāgraḥ}]{\skm{m }ekāgraḥ}
		\rdg[wit={J5}]{(e)kāgra}% me om. J5
		\rdg[wit={N23,Delta,G11,J10,Pi}]{ekāgraṃ}% aikā° G11
		\rdg[wit={J7}]{ekāgryaṃ}
		\rdg[wit={N19}]{evāgra}
		\rdg[wit={V15}]{evāgraḥ}
		}}
\pada{kutracin na hi dhāvati//\versenr} % naṃ hi P11
	\orgvnr{48}\\!}
\end{tlg}
%</vs48>
\commcite\newpage


%<*vs49>
\begin{tlg}[hp04_049]
\tl{
\pada{\app{\lem[wit={ceteri}]{manomatta}
		\rdg[wit={N23}]{manomantra}
		\rdg[wit={J10,V3}]{manonmatta}
		}gajendrasya} % maje° N19; °āsya V15
\pada{\app{\lem[wit={ceteri}]{viṣayodyāna}% °naṃ P11
		\rdg[wit={J5}]{viṣayodhanu}
		\rdg[wit={G4}]{viṣayeṣudra}
		\rdg[wit={G5}]{viṣayāraṇya}
		\rdg[wit={C6}]{viṣayoḍyā}
		\rdg[wit={V3}]{viṣayodhāma}
		}%
	\app{\lem[wit={ceteri}]{cāriṇaḥ}
		\rdg[wit={J5}]{vāriṇaḥ}
		\rdg[wit={G4}]{cāraṇā[ḥ]}
		\rdg[wit={N23}]{vāriṇaṃ}
		\rdg[wit={P11}]{cāriṇaṃ}
		}/}\\+}
\tl{
\pada{\app{\lem[wit={N3,G4,Delta,G5,V3}]{niyāmana}
		\rdg[wit={J5,P11,C6}]{niyamena}
		\rdg[wit={N23}]{niyamitra}
		\rdg[wit={J7}]{niryāmana}
		\rdg[wit={G11,V15}]{niyāmane}
		\rdg[wit={N19}]{niryāsane}
		\rdg[wit={J10}]{nīyamānaḥ}
		\rdg[wit={Jyo}]{niyamane}
		}%
	\app{\lem[wit={ceteri},alt={samartho 'yaṃ}]{samartho'yaṃ}% samatho C6f
		\rdg[wit={G11}]{samartheyaṃ}
		}}
\pada{\app{\lem[wit={ceteri}]{ninādo}% N3,J5,Gamma,Delta,G11,Pi
		\rdg[wit={Zeta2,J10,Jyo}]{nināda}
		}
	\app{\lem[wit={ceteri}]{niśitāṅkuśaḥ} % niśinā° N23; °kuśa V3, °kuśaṃ J10, ṅ om. G11
		\rdg[wit={N3}]{niyatāṃkuśaḥ}
		\rdg[wit={Delta}]{niścayāṅkuśaḥ}
		\rdg[wit={N19}]{niśatāṅkuḥ}
		}//\versenr}\label{manomatta}
		\orgvnr{49}\\!}
\end{tlg}
%</vs49>
\commcite\newpage


%<*vs50>
\begin{tlg}[hp04_050]
\tl{
\pada{\app{\lem[nolem]{\skp{pāda a}}
	\rdg[wit={G5}]{\om}}%
	\app{\lem[wit={ceteri}]{antaraṅga}
		\rdg[wit={V19}]{aṃtaraṃgaṃ}
		\rdg[wit={J10}]{nādoṃtaraṃ}}%
	\app{\lem[wit={G11,C6,V3},alt={°sya javino}]{\skp{°}sya javino}
		\rdg[wit={N3,J5}]{sya javinaḥ}
		\rdg[wit={Gamma,Delta,Zeta2}]{turaṅgasya} % truraṃga N19
		\rdg[wit={J10}]{tu saṃgamya}
		\rdg[wit={P11}]{sya ca mano}
		\rdg[wit={Jyo}]{sya yamino}
		}}
\pada{\app{\lem[nolem]{\skp{pāda b}}
	\rdg[wit={G5}]{\om}}%
	\app{\lem[wit={Zeta2,J10,Pi,Jyo}]{vājinaḥ}
		\rdg[wit={N3,J5}]{kariṇaḥ}
		\rdg[wit={Gamma,Delta}]{vijñānaṃ}% °na N23
		\rdg[wit={G11}]{<<ga>>jasya}% G5 omits this hemistich.
		} 
	\app{\lem[wit={N3,G11,Jyo}]{parighāyate}% +C8,P7
		\rdg[wit={J5,Gamma,N19,J10,V3}]{paridhāyate}
		\rdg[wit={V19}]{parimeyate}
		\rdg[wit={E2}]{parameyate}
		%\rdg[wit={K3,C7}]{parimīyate}
		\rdg[wit={V15}]{paridhāvataḥ}
		\rdg[wit={P11}]{parighātayaḥ}
		\rdg[wit={C6}]{pariṣāyate}
		}/}\\+}
\tl{
\pada{\app{\lem[wit={ceteri}]{nādopāstir ato}% N3,J5,Delta,G11,G5,N19,V15,Pi,Jyo
		\rdg[wit={Gamma}]{nādopāstivato}
%		\rdg[wit={C7}]{nādopāstimato}
%		\rdg[wit={V19}]{nādopāstiratir}% 1st occurrence
		\rdg[wit={J10}]{\om}
		} 
		nitya}%m % nityaṃm N3
\pada{\app{\lem[wit={N3,J5,V19a,P11,V3},alt={avadhāryāpi}]{\skm{m }avadhāryāpi}% +V19(1),C7(1)
		\rdg[wit={N23a}]{anadhāyāpi}% 1st occurrence
		\rdg[wit={J7a}]{avadhāyāpi}
		\rdg[wit={N23b,J7b,E2b,G11}]{avagamyā hi}
		\rdg[wit={V19b}]{avagamya hi}
		\rdg[wit={G5}]{avagamyo hi}
		\rdg[wit={N19}]{avagamyaṃ hi}
		\rdg[wit={V15,Jyo}]{avadhāryā hi}
		\rdg[wit={J10}]{\om}
		\rdg[wit={C6}]{avadhāryo pi}
		}
	\app{\lem[wit={J5,Pi,Jyo}]{yoginā}
		\rdg[wit={N3,G11,Zeta2}]{yogināṃ}
		\rdg[wit={N23a,J7a,V19a,G5}]{yoginaḥ}% +K3,C7   1st
		\rdg[wit={N23b,J7b,V19b,E2b}]{yogibhiḥ}
		\rdg[wit={J10}]{\om}
		}//\versenr}%
\myfn{\xexclude{\getsiglum{Gamma,Delta} have a different verse order: 4.50cd (except \getsiglum{E2}) \textrightarrow\ 4.51 \textrightarrow\ \ref{anahata} \textrightarrow\ 4.50. }\getsiglum{Gamma,V19} have 4.50cd\,=\,X4.105cd twice. The first time (a), their reading of the last \emph{pāda} is closer to the \textalpha\ reading \textit{avadhāryāpi yoginaḥ}, while the second time (b) it is \textit{avagamyā hi yogibhiḥ}, which is closer to the reading of the expanded version.}
		\label{IV95Vu}
		\orgvnr{50}\\!}
\end{tlg}
%</vs50>
\commcite%\newpage


\iffalse
\startaltrecension
%<*vs50-1>
\begin{alttlg}[hp04_050_1]
\tl{\app{\lem[nolem]{}
	\rdg[wit={Gamma,Delta1}]{\also}}%
\pada{\app{\lem[wit={Gamma,E2,Epsi,V15,P11,Jyo},alt={nādo 'ntaraṅga}]{nādo'ntaraṅga}
		\rdg[wit={V19}]{nādaturaṃga}
		\rdg[wit={N19}]{nādāṃtaraṅga}
		\rdg[wit={J10}]{\om}
		\rdg[wit={C6,V3}]{nādotaraṅga}
		}%
	\app{\lem[wit={ceteri}]{sāraṅga}% sātaṅga? E2
%		\rdg[wit={C7}]{mātaṃga}
		\rdg[wit={J10}]{\om}}}% sāraṅgaṃ V15
\pada{\app{\lem[wit={ceteri}]{bandhane}
		\rdg[wit={N23}]{baṃdhāna}
		\rdg[wit={J10}]{\om}
		\rdg[wit={V3}]{baṃdhana}
		}
	\app{\lem[wit={ceteri}]{vāgurāyate}
		\rdg[wit={N23}]{yāgurāyate}
		\rdg[wit={J10}]{\om}}/}\\+}
\tl{
\pada{\app{\lem[nolem]{\skp{pāda c}}
	\rdg[wit={Epsi}]{\om}}% +M3
\app{\lem[wit={ceteri}]{antaraṅga}
		\rdg[wit={Zeta2}]{antaraṅgaṃ}}%
	\app{\lem[wit={V15,Jyo}]{kuraṅgasya}
		\rdg[wit={Gamma,V19,N19,J10,Pi}]{turaṅgasya}
		\rdg[wit={E2}]{turaṅgasyā}% +K3,C7
		}}
\pada{\app{\lem[nolem]{\skp{pāda d}}
	\rdg[wit={Epsi}]{\om}}% +M3
\app{\lem[wit={N19,Pi}]{rodhe}
		\rdg[wit={N23}]{bāhye}
		\rdg[wit={J7}]{bodho}
		\rdg[wit={V19}]{\lacuna}
		\rdg[wit={E2}]{vabo}% +C7; °varodhe K3
		\rdg[wit={V15}]{nādo}
		\rdg[wit={J10}]{rogo}
		\rdg[wit={Jyo}]{vadhe}
		}\marmas
	\app{\lem[wit={V15,Jyo}]{vyādhāyate}
		\rdg[wit={Gamma}]{pi līyate}
		\rdg[wit={V19}]{\lacuna}
		\rdg[wit={E2}]{dhe līyate}% +K3,C7
		\rdg[wit={N19}]{vā gāyate}
		\rdg[wit={J10}]{vā gīyate}
		\rdg[wit={P11}]{vādyāyate}
		\rdg[wit={C6}]{pi pariṣā}%  parighāyate P7
		\rdg[wit={V3}]{vādhāyate}
		}%
	\app{\lem[wit={ceteri}]{'pi ca}
		\rdg[wit={V19}]{\lacuna}
		\rdg[wit={P11}]{ti ca}
		\rdg[wit={C6}]{yate}
		}//\versenr}%
		\myfn{In \getsiglum{G11,G5}, the first hemistich is found between \ref{IV95Vu}ab and cd, and the second hemistich is omitted;
		in \getsiglum{Pi,Jyo}, the whole verse is found before \ref{IV95Vu};
		\getsiglum{J10} merges the two verses into one:
		\devnote{nādo'ntaraṃ tu saṃgamya vājinaḥ paridhāyate/
		antaraṅgaturaṃgasya rogo vā gīyate'pi ca//}}
		\\!}
\end{alttlg}
%</vs50-1>
\endaltrecension
\fi


%
%% 4.064_1 (cf. 4.65cd)
%\pada{nādopāsti\app{\lem[wit={K3}]{r ato}
%		\rdg[wit={V19}]{ratir}
%		\rdg[wit={Gamma}]{vato}
%		\rdg[wit={C7}]{mato}} nityam}
%\pada{\app{\lem[wit={Delta2}]{avadhāryāpi}
%		\rdg[wit={J7}]{avadhāyāpi}
%		\rdg[wit={N23}]{anadhāyāpi}
%		\rdg[wit={K3}]{avidhāryaṃ hi}} yoginaḥ/}
%		\sgwit{Gamma,Delta} \anm{≈ \ref{IV95Vu}cd}\\!}


%<*vs51>
\begin{tlg}[hp04_051]
\tl{
\pada{\app{\lem[nolem]{\skp{pāda a}}
	\rdg[wit={E2,G5,Zeta2,J10}]{\om}
	}%
\app{\lem[wit={N3,J5,P11,V3,Jyo},post=\texteng{(ādī \getsiglum{N3})}]{ghaṇṭādināda}
%	\rdg[wit={N3}]{ghaṃṭādīnāda}
	\rdg[wit={Gamma,V19,G11,C6}]{ghaṇṭānināda}}% ghaṃṭāki? V19
\app{\lem[wit={V3,Jyo}]{sakta}
	\rdg[wit={N3}]{śaktaś ca}
	\rdg[wit={J5}]{śakti}
	\rdg[wit={Gamma,V19,Epsi}]{saktasya}% sukta N23, °syaṃ? V19
	\rdg[wit={P11}]{sadaṃkatā}
	\rdg[wit={C6}]{kuliśa}
	}}%
\pada{\app{\lem[wit={G5,Jyo}]{stabdhāntaḥ}
	\rdg[wit={N3}]{stavyāṃtaḥ}
	\rdg[wit={J5}]{stadhvāṃta}
	\rdg[wit={N23}]{sabdāntaḥ}
	\rdg[wit={J7}]{śabdataḥ}
	\rdg[wit={V19}]{śuddhāntaḥ}%  °syaṃ? V19; +K3,C7
	\rdg[wit={G11}]{stabdhasyāntaḥ}
	\rdg[wit={P11}]{stabdhyaṃtaḥ}
	\rdg[wit={C6}]{pradhvānta}
	\rdg[wit={V3}]{statravadhātaḥ}
	}%
\app{\lem[wit={G11,P11,V3,Jyo}]{karaṇahariṇasya}
	\rdg[wit={N3}]{karaṇaṃ hariṇasya}
	\rdg[wit={J5}]{karaṇaṃ mṛgasya}
	\rdg[wit={N23}]{karaṇasya na}
	\rdg[wit={J7,V19,G5}]{karaṇasya ca}
	\rdg[wit={C6}]{hariṇasya ca}
	}/} 
	\\+}
\tl{
\pada{\app{\lem[nolem]{\skp{pāda c}}
	\rdg[wit={Gamma,Delta,Zeta2,J10}]{\om}}%
	praharaṇa%m % karṇam! J5
\app{\lem[wit={Epsi,Pi},alt={atisukaraṃ}]{\skm{m }atisukaraṃ}
	\rdg[wit={N3}]{atisukasteraṃ}
	\rdg[wit={J5}]{avisukaraṇaṃ}
	\rdg[wit={Jyo}]{api sukaraṃ}}
\app{\lem[wit={N3,Epsi,P11,C6,Jyo},alt={syāc/t}]{syāc\skp{/t}}% so P7
	\rdg[wit={J5}]{\om}
	\rdg[wit={V3}]{syā}
	}}%
\pada{\app{\lem[nolem]{\skp{pāda d}}
	\rdg[wit={Gamma,Delta,Zeta2,J10}]{\om}}%
\app{\lem[wit={N3,J5,G11,P11,Jyo}]{chara}% so P7
	\rdg[wit={G5}]{kara}
	\rdg[wit={C6}]{sadṛ}
	\rdg[wit={V3}]{ra}
	}%
\app{\lem[wit={N3,G11,P11,V3}]{saṃdhātā}
	\rdg[wit={J5}]{saṃdhā}
	\rdg[wit={C6}]{śaṃ dhātā}
	\rdg[wit={G5,Jyo}]{saṃdhāna}} 
	pravīṇaś cet//\versenr} % vra° N3; °ṇāś G11, ca G5
	\orgvnr{51}\\!}
\end{tlg}
%</vs51>
\commcite\newpage
% P11
% ghaṃṭādinādasadaṃkatāstabdhāṃtaḥkaraṇahariṇasya/
% praharaṇam atisukaraṃ syāc charasaṃdhātā pravīṇaś cet//
% M1:
% ghaṃṭāninādasaktastabdhāṃtaḥkaraṇahariṇasya/
% praharaṇe sukaraṃ syā(c) [ch](ara)[sa]ṃdhānapravīṇaś cet// 4.103//
% P7:
% ghaṃṭādinādakuliśapradhvāṃtaharisasya ca/
% praharaṇam atisūkaraṃ syā<c> charasaṃghātā praviṇaś cet// 106/107 //
% C6:
% ghaṇṭāninādakuliśapradhvāṃtahariṇasya ca/
% praharaṇam atisukaraṃ syāt sadṛśaṃdhātā pravīnaś cet//
% N12:
% ghaṇṭādinākalādhvāṃtaḥkaraṇahariṇyaḥ syāt/
% praharaṇam atisukaraṃ syāc charasaṃdhānā pravīṇaś ca/


%<*vs52>
\begin{tlg}[hp04_052]
\tl{\app{\lem[nolem]{}
	\rdg[wit={G5,Zeta2,J10}]{\om}}%
\pada{\app{\lem[wit={Alpha,Gamma,Delta,G11,P11,V3,Jyo}]{anāhatasya śabdasya}% sabda° V3,N23
%		\rdg[wit={N23,V3}]{anāhatasya sabdasya}
		\rdg[wit={C6}]{anāhatas tu yaḥ śabdas}}}
\pada{\app{\lem[wit={J5,Gamma,Delta,C6}]{tasya śabdasya yo dhvaniḥ}
		\rdg[wit={N3}]{tasya śabdasya ca dhvaniḥ}% +F
		\rdg[wit={G4}]{tasya yo dhvaniḥ}
		\rdg[wit={G11}]{tasya śabdasya yā dhvaniḥ}
		\rdg[wit={P11}]{śabdasyāṃganabho dhvaniḥ}
		\rdg[wit={V3}]{śabdasyāṃtargato dhvaniḥ}
		\rdg[wit={Jyo}]{dhvanir ya upalabhyate}}
		/}\\+}
\tl{
\pada{\app{\lem[wit={N3,Delta,G11,P11,C6,Jyo},postwit=\texteng{\getsiglum{N23}\textsubscript{pc}},alt={dhvaner}]{dhvane\skp{r}}
		\rdg[wit={J5,G4,Gamma,V3}]{dhvanir}}r
		antargataṃ % °gata N23
	\app{\lem[wit={G4,N23,E2,G11},alt={jyotir}]{jyoti\skp{r}}% +G3
		\rdg[wit={N3,Jyo}]{jñeyaṃ}% +F ##
		\rdg[wit={J5,C6}]{\om}
		\rdg[wit={J7,V19}]{jyoti} % jyotī P7
		\rdg[wit={P11,V3}]{geyaṃ}% geya P11
		}}%
\pada{\app{\lem[wit={Gamma,G11},alt={jyotirantar}]{\skm{r }jyotiranta\skp{r}}% +K3
		\rdg[wit={N3}]{yasyāṃtvaṃtar}
		\rdg[wit={J5}]{yotiraṃtar}
		\rdg[wit={G4}]{jyoti\,..\,.. }
		\rdg[wit={Delta,C6}]{jyoterantar}
		\rdg[wit={P11,V3}]{geyasyāntar}
		\rdg[wit={Jyo}]{jñeyasyāntar}% +F
		}rgataṃ manaḥ/}\\+} % mana V3
\tl{
\pada{\app{\lem[wit={N3,J7,P11,V3}]{tan mano vilayaṃ}% lost G4
		\rdg[wit={J5}]{tan maṃnaṃ vilayaṃ}
		\rdg[wit={N23,Delta,C6}]{yan mano vilayaṃ}
		\rdg[wit={G11}]{tan mano nilayaṃ}
		\rdg[wit={Jyo}]{manas tatra layaṃ}% +F
		}
	\app{\lem[wit={J5,N23,Delta,G11,C6,V3,Jyo}]{yāti}
		\rdg[wit={N3,J7,P11}]{yāṃti}}}
\pada{tad viṣṇoḥ paramaṃ padam//\versenr}% viṣṇo N3,J5,C6
	%\sgwit{Alpha,Gamma,Delta,G11,Pi,Jyo}%
	\myfn{cf. \manuref{X4.107}}
	\label{anahata}
	\orgvnr{52}\\!}
\end{tlg}
%</vs52>
\commcite\newpage


\iffalse
\startaltnormal
%<*vs52-1>
\begin{alttlg}[hp04_052_1]
\tl{\app{\lem[nolem]{}
	\rdg[wit={Epsi,Zeta2,J10}]{\incl}}% not in Omega,Jyo
\pada{anāhatadhvaner anta}%r anāhataḥ V17; anta<<r>> J10
\pada{\app{\lem[wit={Zeta2,J10},alt={°r jñeyaṃ yat}]{\skp{°}r jñeyaṃ ya\skp{t}}% gataṃ M3; +G7
	\rdg[wit={G11}]{r geyaṃ yat}
	\rdg[wit={G5}]{m āpnuyāt}}t
	sūkṣma\app{\lem[wit={Zeta2,J10}]{sūkṣmakam} % +M3; sūkṣmya° N19
	\rdg[wit={Epsi}]{sūkṣmataḥ}}/}\\+}
\tl{
\pada{manas tatra layaṃ yāti}
\pada{tad viṣṇoḥ paramaṃ padam//\versenr}% viṣṇo
\myfn{cf. \ref{anahata}}
\\!}
	%\myfn{\getsiglum{G11,G5,N19,V15,J10} have this verse as a substitute for \ref{anahata}; \getsiglum{G11} also contains 4.52 (see ??).}
%	-- Alt\,2 here and Alt\,1 after \ref{sadanada} (preceded by three additional lines: 
%  while \getsiglum{G5,M3} have the latter only.
	% \devnote{vindur bhidyati nādena sa nādaḥ khena bhidyate/
	% oṃkāradhvaninādena vāyus saṃharaṇāntikaṃ/
	% nirālaṃbaṃ samuddiśya yatra nādo layaṃ gataḥ//}) --, 
\end{alttlg}
%</vs52-1>
\endaltnormal
\fi


%<*vs53>
\begin{tlg}[hp04_053]
\tl{\app{\lem[nolem]{}
	\rdg[wit={E2}]{\om}}%
\pada{\app{\lem[wit={ceteri},alt={tāvad ā°}]{tāvad ā}% tavad N23
		\rdg[wit={J10}]{bhāvanā}}kāśasaṃkalpo} % kāsa J10; saḥkalpo J7
\pada{\app{\lem[wit={ceteri}]{yāvac chabdaḥ}% ḥ om. J5,P11,P7; N3,J5,Gamma,G11,G5,V15,J10,Pi,Jyo
		\rdg[wit={V19}]{yāvad bandhaḥ}% +C7
		\rdg[wit={N19}]{yāvad vādhaḥ}} pravartate/}\\+}
\tl{
\pada{niḥśabdaṃ  % niśabdaṃ V19,N3,J5; °bdāṃ V17
	\app{\lem[wit={ceteri}]{tat paraṃ}
		\rdg[wit={N23}]{paramaṃ}} brahma}
\pada{\app{\lem[wit={ceteri}]{paramātmā} % paramatmā V19
		\rdg[wit={Epsi,Jyo}]{paramātme}}
	\app{\lem[wit={N3,J7,V3}]{samīryate}
		\rdg[wit={J5,N23,V19,P11}]{samīyate}% +G3
		\rdg[wit={G4}]{samīkṣate}
		\rdg[wit={Epsi,Jyo}]{ti gīyate}
		\rdg[wit={Zeta2,J10}]{numīyate}
		\rdg[wit={C6}]{yam īryate}
		}//\versenr}
	\orgvnr{53}\\!}
\end{tlg}
%</vs53>
\commcite%\newpage


%<*vs54>
\begin{tlg}[hp04_054]
\tl{\app{\lem[nolem]{}
	\rdg[wit={E2,Zeta2,J10}]{\om}}%
\pada{\app{\lem[wit={Alpha,Gamma,V19,Epsi,P11,C6,Jyo},alt={yat}]{ya\skp{t}}
		\rdg[wit={V3}]{\om}}t kiṃci%n
	\app{\lem[wit={Alpha,Epsi,Pi,Jyo},alt={nāda}]{\skm{n }nāda}
		\rdg[wit={Gamma,V19}]{nāma}}rūpeṇa}
\pada{śrūyate śaktir eva sā/}\\+} % śū° N3
\tl{
\pada{\app{\lem[wit={N3,Gamma,Epsi,P11}]{yas tacchrotā}% +K3,C7
		\rdg[wit={J5}]{yasmin śrato}
		\rdg[wit={V19}]{yat ta[cch]roto}
		\rdg[wit={C6}]{yas tatsrotā}
		\rdg[wit={V3}]{yac chrotā ca}
		\rdg[wit={Jyo}]{yas tattvānto}
		} 
		nirākāraḥ} % °kāra P11,V3; °kārā J5, °dhāraḥ G5
\pada{sa eva parameśvaraḥ//\versenr}
\orgvnr{54}\\!} % ḥ om. J5,V3; e<<va>> N23
\end{tlg}
%</vs54>
\commcite\newpage


%<*vs55>
\begin{tlg}[hp04_055]
\tl{
\pada{śravaṇa\app{\lem[wit={N3,J5,G11,Zeta2,Pi}]{mukha}
		\rdg[wit={Gamma,Delta,J10,Jyo}]{puṭa}
		\rdg[wit={G5}]{\om}}%
	\app{\lem[wit={ceteri}]{nayana} % cayana N23, naya<<na>> N3
		\rdg[wit={J10,Jyo}]{nayanayugala}
		}%
	\app{\lem[wit={ceteri}]{nāsā} % nāśā V19,N19,J5,V3,J10
		\rdg[wit={Jyo}]{ghrāṇa}
		}}%
\pada{\app{\lem[wit={J5,G11,Zeta2,P11,C6}]{nirodhanaṃ naiva kartavyam}% nirodhamaṃ N19
		\rdg[wit={N3,G5}]{nirodhaṃ naiva kartavyaṃ}
		\rdg[wit={Gamma,E2}]{mukhapuṭasaṃrodhanaṃ kāryam}
		\rdg[wit={V19}]{mukhapuṭarodhane kāryaṃ}
		\rdg[wit={J10}]{mukharodhanam eva kartavyaṃ}
		\rdg[wit={V3}]{nirodhanenaiva kartavyaṃ}
		\rdg[wit={Jyo}]{mukhānāṃ nirodhanaṃ kāryam}}/}\\+}
\tl{
\pada{\app{\lem[wit={ceteri}]{śuddha}
		\rdg[wit={Gamma}]{śrīśuddha} % suddha N23
		\rdg[wit={V3}]{\om}}%
	\app{\lem[wit={ceteri}]{suṣumṇā}% śu° P11,V15
		\rdg[wit={N23}]{suṣumū}}% 
	\app{\lem[wit={J7,Delta,G11,Jyo}]{saraṇau}% śaraṇau V19,J10pc
		\rdg[wit={N3}]{tsaraṇaḥ}
		\rdg[wit={J5}]{śarada}
		\rdg[wit={G4}]{saraṇaiḥ}
		\rdg[wit={N23}]{ṇau}
		\rdg[wit={G5,Zeta2,J10}]{śaraṇe}
		\rdg[wit={P11}]{tmakārausaṃ}
		\rdg[wit={C6}]{tmaśaraṇaiḥ}
		\rdg[wit={V3}]{maraṇai}
		}}
\pada{\app{\lem[wit={ceteri}]{sphuṭam amalaḥ}% amala J5,G4, amalaṃ G5,V3
		\rdg[wit={N19}]{vimalaḥ}
		\rdg[wit={V15}]{vimalaḥ saṃ}
		\rdg[wit={C6}]{saṃsphurad amalaḥ}
		}
	śrūyate nādaḥ//\versenr} % śṛyate N3, śraya .. + + G11; nāda J5,P11,V3
	\label{sravanaputa}
	\orgvnr{55}\\!}
\end{tlg}
%</vs55>
\commcite%\newpage


\iffalse
\startaltrecension
%<*vs55-1>
\begin{alttlg}[hp04_055_1]
\tl{\app{\lem[nolem]{}
	\rdg[wit={Jyo}]{\NotIn}}%
\pada{\app{\lem[wit={V15,J10,C6,V3}]{nādaḥ}% +M3
		\rdg[wit={Epsi,N19,P11}]{nāda}} % +F
		śaktir iti
	\app{\lem[wit={V15,J10}]{khyāto} % khyato J10ac
		\rdg[wit={Epsi}]{khyātā}% +M3
		\rdg[wit={N19}]{kṣāto}
		\rdg[wit={P11}]{jñeyaṃ}
		\rdg[wit={C6}]{jñeyā}% +F
		\rdg[wit={V3}]{jñeya}
		}}
\pada{\app{\lem[wit={Epsi,Zeta2,P11,V3}]{nādajñānaṃ}
		\rdg[wit={J10,C6}]{nādo jñānaṃ}} sadāśivaḥ/}\\+} % śiva P11,V3
\tl{
\pada{\app{\lem[wit={G11}]{jñeye jñāne ca naṣṭe tu}
		\rdg[wit={G5}]{jñeyajñāne ca naṣṭe ca}
		\rdg[wit={N19}]{nādajñāne ca neṣṭe tad}% +J11ac
		\rdg[wit={V15}]{nādajñāne vinaṣṭe ca tad}
		\rdg[wit={J10}]{nādajñānena naṣṭena}% +J11pc
		\rdg[wit={P11}]{jñeye jñāne vilīnīṃta}% jñeyajñāne hi līyete
		\rdg[wit={C6}]{jñeyo jñāne vilīne tu}% vilīyeta P7
		\rdg[wit={V3}]{jñeye jñāne vilineṃta}
		}}
\pada{\app{\lem[wit={Epsi,V15},alt={unmany}]{unma\skp{ny}}% +J11
		\rdg[wit={N19}]{unmadhy}
		\rdg[wit={J10}]{hy unmany}
		\rdg[wit={Pi}]{sonmany}% +F
		}%
	\app{\lem[wit={Epsi,J10,C6},alt={evāvaśiṣyate}]{\skm{ny }evāvaśiṣyate}
		\rdg[wit={N19}]{edhāvaśiṣyate}
		\rdg[wit={V15}]{eva śiṣyate}
		\rdg[wit={P11}]{enāvaśiṣyati}
		\rdg[wit={V3}]{avāvaśiṣyate}
		}//\versenr}
		\\!}
\end{alttlg}
%</vs55-1>

%<*vs55-2>
\begin{alttlg}[hp04_055_2]
\tl{\app{\lem[nolem]{}
	\rdg[wit={Jyo}]{\NotIn}}%
\pada{nādo yāvan manas tāva}%n
\pada{\app{\lem[wit={Epsi,N19,J10,P11,V3},alt={nādānte tu}]{\skm{n }nādānte tu}
	\rdg[wit={V15}]{nādānte ca}
	\rdg[wit={C6}]{nādātīte}} manonmanī/}\\+}
\tl{
\pada{saśabdaṃ kathitaṃ vyoma} % na śabdaṃ G5
\pada{niḥśabdaṃ brahma 
	\app{\lem[wit={Zeta2,J10,Pi}]{kathyate}
		\rdg[wit={Epsi}]{ucyate}% +F
	}//\versenr}
\\!}
\end{alttlg}
%</vs55-2>

%<*vs55-3>
\begin{alttlg}[hp04_055_3]
\tl{
\pada{sadā nādānusandhānā}%t °bandhānāt G5
\pada{\app{\lem[wit={Epsi,Zeta2,J10,Pi},alt={saṃkṣīṇe}]{\skm{t }saṃkṣīṇe}% sa G5,P11
		\rdg[wit={Jyo}]{kṣīyante}}
	\app{\lem[wit={Epsi,P11,C6}]{vāsanācaye}% +P7
		\rdg[wit={N19}]{vāsanākṣaye}
		\rdg[wit={V15}]{vāsanākṣaṇe}
		\rdg[wit={J10}]{vāsanodaye}% °dvaye F
		\rdg[wit={V3}]{vāsanāvayo}
		\rdg[wit={Jyo}]{pāpasaṃcayāḥ}
		}/}\\+}
\tl{
\pada{nirañjane
	\app{\lem[wit={G11,Jyo}]{vilīyete}% +F
		\rdg[wit={G5,N19}]{ca līyeta}
		\rdg[wit={V15,J10}]{ca līyete}
		\rdg[wit={P11,V3}]{vilīyaṃte}
		\rdg[wit={C6}]{vilīyeta}
		}}
\pada{\app{\lem[wit={G11}]{niścitaṃ manamārutau}% ##
		\rdg[wit={G5}]{niścitaṃ manamārute}
		\rdg[wit={N19}]{niścitta manamārutau}
		\rdg[wit={V15,Jyo}]{niścitaṃ cittamārutau}
		\rdg[wit={J10}]{niścitau manamārutau}% +F
		\rdg[wit={P11}]{niścitaṃ māruto manaḥ}
		\rdg[wit={C6}]{marutā niścitaṃ manaḥ}
		\rdg[wit={V3}]{niścita māruto mana}
		}//\versenr}
	\label{sadanada}\\!}
\end{alttlg}
%</vs55-3>

%<*vs55-4>
\begin{alttlg}[hp04_055_4]
\tl{\app{\lem[nolem]{}
	\rdg[wit={Gamma,Delta1}]{\textapp{included after \ref{makaranda1}}}
	\rdg[wit={Jyo}]{\NotIn}}%
\pada{nādakoṭisahasrāṇi} % sahasrāṇī J7, sahasrāni N23, °śrāṇi V3,V19
\pada{\app{\lem[wit={ceteri}]{bindu}
	\rdg[wit={C6}]{veda}
	}koṭiśatāni ca/}\\+} % saṭāni N23ac, satāni N23pc
\tl{
\pada{\app{\lem[wit={ceteri}]{sarve}
		\rdg[wit={N23}]{sarvaṃ}
		} tatra layaṃ
	\app{\lem[wit={ceteri}]{yānti}
		\rdg[wit={V19,C6}]{yāti}
		}}
\pada{yatra % ya ca P11
	\app{\lem[wit={ceteri}]{devo}% +G5,M3
		\rdg[wit={G11}]{deve}
		\rdg[wit={N19,V3}]{deva}}
	\app{\lem[wit={ceteri}]{nirañjanaḥ}% +G5,M3; °jana P11, janaṃ V3
		\rdg[wit={G11}]{nirañjane}
		}//\versenr}
	\label{nadakoti}
	\\!}
\end{alttlg}
%</vs55-4>

%<*vs55-4p>
\begin{altpostmula}[hp04_055_4p]
\app{\lem[nolem]{}%
	\rdg[wit={C6}]{\textapp{found between \emph{pāda}s ab and cd of the next verse}}}%
\app{\lem[wit={G11,Zeta2,J10,P11,Jyo}]{iti nādānusandhānam}
	\rdg[wit={G5,C6,V3}]{iti nādānusandhānavidhiḥ} % vidhi V3; +F
}//
%\sgwit{G11,G5,N19,V15,J10,Pi,Jyo}
\end{altpostmula}
%</vs55-4p>

%%%%%%%%%%%%%%%%%%%%%%%


%<*vs55-5a>
\begin{altava}[hp04_055_5a]
\app{\lem[nolem]{}
	\rdg[wit={Epsi}]{\only}}%
atha rājayogaḥ/
\end{altava}
%</vs55-5a>

%<*vs55-5>
\begin{alttlg}[hp04_055_5]
\tl{\app{\lem[nolem]{}
	\rdg{\texteng{=\,1.64}}
	\rdg[wit={Epsi}]{\incl}
	\rdg[wit={N19}]{\textapp{abbreviated to} yathā vṛddho veti}
	\rdg[wit={V15}]{yathā vṛddhaiḥ prabhāṣitaṃ}
	}%
\pada{yuvā 
	\app{\lem[wit={G11},alt={vṛddho 'tivṛddho vā}]{vṛddho'tivṛddho vā}
		\rdg[wit={G5}]{bhavatu vṛddho vā}}} 
\pada{vyādhito durbalo'pi 
	\app{\lem[wit={G11}]{vā}
		\rdg[wit={G5}]{ca}}/}\\+}
\tl{
    \pada{abhyāsāt siddhim āpnoti} 
    \pada{sarvayogeṣv atandritaḥ//\versenr}
    \\!}
\end{alttlg}
%</vs55-5>

\endaltrecension
\fi


%<*vs56>
\begin{tlg}[hp04_056]
\tl{\app{\lem[nolem]{}
	\rdg[wit={G5,Zeta2,J10,Jyo}]{\om}
	}%
\pada{\app{\lem[nolem]{\skp{pāda a}}
	\rdg[wit={J5}]{\om}}%
\app{\lem[wit={N3,G4,Gamma,G11,Pi}]{kāṣṭha}
		\rdg[wit={Delta}]{koṣṭha}}% ab lost J5; none N24
	\app{\lem[wit={Delta,G11}]{goṣṭhī}
		\rdg[wit={N3,G4,J7}]{goṣṭhi}% +F
		\rdg[wit={N23,V3}]{goṣṭha}
		\rdg[wit={P11}]{mathnī}
		\rdg[wit={C6}]{mathnā}}%
	\app{\lem[wit={G11,V3}]{prapañcena}
		\rdg[wit={N3}]{prapaṃce}
		\rdg[wit={G4,Gamma,Delta}]{prasaṅgena}% none J5,N24
		\rdg[wit={P11}]{pravacane}
		\rdg[wit={C6}]{pravartaṃ}}\marma}
\pada{\app{\lem[nolem]{\skp{pāda b}}
	\rdg[wit={J5}]{\om}}%
\app{\lem[wit={N3,G4,G11,Pi}]{kiṃ sakhe śrūyatām idam}% śṛ° N3; °tam P11
		\rdg[wit={N23}]{nāgadaṃtaṃmatargataṃ sṛṇu}
		\rdg[wit={J7,Delta}]{nādam antargataṃ śṛṇu} % śruṇu V19
		}/}
	\\+}
\tl{
\pada{purā matsyendra\app{\lem[wit={N3,J5,G11,Pi},alt={bodhārtham}]{bodhārtha\skp{m}}
		\rdg[wit={Gamma,Delta}]{bodhāya}}}%
\pada{\app{\lem[wit={N3,J5,J7,Delta,G11,P11,C6},alt={ādināthoditaṃ}]{\skm{m }ādināthoditaṃ}
		\rdg[wit={N23}]{ādināthotigaditaṃ}
		\rdg[wit={V3}]{ānināthodinaṃ}}
		vacaḥ//\versenr}\label{kastha}
	\orgvnr{56}\\!} % vaca N3,V3
\end{tlg}
%</vs56>
\commcite\newpage

% C6 4.110: kaṣṭamathnā pravartaṃ (1 syll. ami.) kiṃ sakhe śrūyatām idam/
% purā matsyendrabodhārtham ādināthoditaṃ vacaḥ//
% P7: kāṣṭamathnvāṃ pravartante kiṃ sakhe śrūyatām idam/
% purā matsyendrabodhārtham ādināthoditaṃ vacaḥ//
% C8: kāṣṭamanāpravartaṃte ... bodhāya
% P11: kāṣṭhamathnī pravacane ... bodhārtham


%<*vs57>
\begin{tlg}[hp04_057]
\tl{\app{\lem[nolem]{}
	\rdg[wit={Jyo}]{\textapp{found at the end of the chapter}}}%
\pada{yāvan naiva % yāvān eva G5
	\app{\lem[wit={ceteri}]{praviśati}
		\rdg[wit={N23}]{\_viśati}}
	\app{\lem[wit={ceteri},alt={caran}]{cara\skp{n}}
		\rdg[wit={N3}]{care}
		\rdg[wit={N23}]{palan}
		\rdg[wit={J7}]{calan}
		\rdg[wit={G5}]{caren}
		\rdg[wit={V3}]{\om}}%n 
	\app{\lem[wit={ceteri},alt={māruto}]{\skm{n }māruto}
		\rdg[wit={N3}]{mārutaṃ}} % mādgato N23
	madhya% mādhya V15
	\app{\lem[wit={N3,J5,J7,V19,Epsi,N19,J10,C6,Jyo}]{mārge}
		\rdg[wit={N23,P11}]{mārgo}
		\rdg[wit={E2,V15}]{mārgaṃ}% +C7
		\rdg[wit={V3}]{mārgā}}}\\+}
\tl{
\pada{yāva%d
	\app{\lem[wit={ceteri},alt={bindur}]{\skm{d }bindu\skp{r}} % bindu P11, vāyuḥ V19ac
		\rdg[wit={N19}]{bandhaṃ}
		\rdg[wit={V15}]{bandho}
		}%r 
	\app{\lem[wit={ceteri},alt={na bhavati}]{\skm{r }na bhavati}
		\rdg[wit={J10}]{bhavati na}}
	\app{\lem[wit={ceteri}]{dṛḍhaḥ} % dṛḍha V19
		\rdg[wit={N3,Epsi,P11}]{dṛḍhaṃ}% +F
		\rdg[wit={J5}]{sthiraḥ}}
	prāṇa\app{\lem[wit={Alpha,J7,Epsi,J10,Pi,Jyo}]{vāta}% prāṇa om. J5, prā<<ṇa>> V3
		\rdg[wit={N23,Delta,V15}]{vātaḥ}
		\rdg[wit={N19}]{vātaṃ}}%
	\app{\lem[wit={Gamma,C6}]{prabaddhaḥ}
		\rdg[wit={N3}]{prabodhaḥ}
		\rdg[wit={J5}]{prakṛddhaḥ}
		\rdg[wit={G4}]{prabaddhaṃ}
		\rdg[wit={Delta,J10}]{prabuddhaḥ}
		\rdg[wit={G11,V15,P11}]{prabandhaḥ}% ḥ om. P11
		\rdg[wit={G5}]{prabandham}
		\rdg[wit={N19}]{na bandhanaḥ}
		\rdg[wit={V3}]{prabodhakaḥ}
		\rdg[wit={Jyo}]{prabandhāt}% +F(twice),M1
		}/}\\+}
\tl{
\pada{\app{\lem[wit={G5,Zeta2,P11,C6}]{yāvad vyomnā}
		\rdg[wit={N3,G4,G11}]{yāvad yomnā}
		\rdg[wit={J5}]{yād vyemnā}
		\rdg[wit={N23}]{yāva\,\_\,mnaḥ}% yāva _ mnaḥ N23
		\rdg[wit={J7,Delta,J10}]{yāvad vyomnaḥ}% dyo° J7; rather °nmaḥ J10
		\rdg[wit={V3}]{yāvad byomna}
		\rdg[wit={Jyo}]{yāvad dhyāne}
		}
	\app{\lem[wit={ceteri}]{sahajasadṛśaṃ}% °sadṛsāṃ J5
		\rdg[wit={N23}]{sahajasaṃśaṃ}
		\rdg[wit={Epsi}]{sadṛśasahajā}% °sahajaṃ F
		} jāyate naiva
	\app{\lem[wit={ceteri}]{tattvaṃ}
		\rdg[wit={V15,J10,V3}]{cittaṃ}}}\\+}
\tl{
\pada{tāva\app{\lem[wit={ceteri},alt={sarvaṃ}]{\skm{t }sarvaṃ}
		\rdg[wit={G11}]{satvaṃ}
		\rdg[wit={J10,V3,Jyo}]{jñānaṃ}
		} vadati % vadaṃti V3
	\app{\lem[wit={N3,J5,J7,E2,G5,Zeta2,J10,C6}]{yad idaṃ}% +C7
		\rdg[wit={N23,P11}]{yadi}% post=\texteng{(daṃ om. by haplography)}
		\rdg[wit={V19,Jyo}]{tad idaṃ}% +K3
		\rdg[wit={G11}]{yadi tat}% yadi taṃ F
		\rdg[wit={V3}]{satataṃ}}
	\app{\lem[wit={ceteri}]{dambha}
		\rdg[wit={G11,N19}]{ḍaṃbha}
		\rdg[wit={G5}]{naiva}
		}mithyā%
	\app{\lem[wit={ceteri}]{pralāpaḥ}% °lāpa J5,N19, prallāpaḥ V3
		\rdg[wit={C6}]{pralābhaḥ}}//\versenr}%
	%\myfn{In \getsiglum{Jyo}, this verse is found at the end of the chapter.}
	\label{yavan}		
	\orgvnr{57}\\!}
\end{tlg}
%</vs57>
\commcite\newpage

%========== Passage B ==========


%<*vs58>
\begin{tlg}[hp04_058]
\tl{
\pada{\app{\lem[wit={ceteri}]{jñātvā}
		\rdg[wit={V15}]{suṣu}
		\rdg[wit={C6}]{jitvā}
		}
	\app{\lem[wit={N3,J5,J10,Jyo}]{suṣumṇāsadbhedaṃ}
		\rdg[wit={N23}]{suṣu<<m>>ṇāṃmedehi}
		\rdg[wit={J7,Delta}]{suṣumṇābhedaṃ hi}% sukhu° V19
		\rdg[wit={G11,Pi}]{suṣumṇāsaṃbhedaṃ}% +F
		\rdg[wit={G5}]{suṣumnāntarbhedaṃ}
		\rdg[wit={N19}]{suṣumṇāṃ saśvedaṃ}
		\rdg[wit={V15}]{mnāṃtagataṃ mārgaṃ}
		}}
\pada{\app{\lem[wit={ceteri}]{kṛtvā vāyuṃ}% kṛtvān V3; vāyu V19, pāyu N19
		\rdg[wit={J5}]{tvāpa vāyuṃ}
		\rdg[wit={V15}]{vāyuṃ kṛtvā}
		} ca
	\app{\lem[wit={ceteri}]{madhyagam}
		\rdg[wit={P11}]{madhyamaḥ}}/}\\+}
\tl{
\pada{\app{\lem[wit={N3}]{kṛtvāsāv aindave sthāne}
		\rdg[wit={V3}]{kṛtvāsāv aidavasthāne}
		\rdg[wit={J5}]{kṛtvā tām aidave tthāne}
		\rdg[wit={G4}]{[dh]ṛ\,..\,[sāv a]ṃdra\,..\,[sthā]ne}% dhṛtvā taṃ caindavasthāne F
		\rdg[wit={N23}]{nītvā tāv iṃdavasthāne}% kṛtvā tāv iṃdavasthāne K1
		\rdg[wit={J7}]{nītvā tāvad avasthāne}% kṛtvā tāvad vindusthairyaṃ B2
		\rdg[wit={Delta}]{nītvā tām anavasthāne}
		\rdg[wit={G11}]{sthitvā savaindave sthāne}
		\rdg[wit={G5}]{sthitvāsāv aindavasthāne}
		\rdg[wit={N19}]{sthitvā sāṃcaiṃdave sthāne}
		\rdg[wit={V15}]{samāvasthā sthito yogī}
		\rdg[wit={J10}]{sthitvā sadaiṃdave sthāne}
		\rdg[wit={P11}]{kṛtvāsav aidavai sthānair}
		\rdg[wit={C6}]{hṛtvā mamedaṃ ca sthānaṃ}
		\rdg[wit={Jyo}]{sthitvā sadaiva susthāne}
		}}
\pada{\app{\lem[wit={Alpha,Epsi,N19,Pi}]{ghrāṇa}
		\rdg[wit={Gamma,Delta,V15,J10}]{prāṇa}% +N24
		\rdg[wit={Jyo}]{brahma}
		}%
	\app{\lem[wit={Alpha,J7,J10,C6,V3,Jyo}]{randhre}
		\rdg[wit={N23,Delta,Epsi,Zeta2}]{randhraṃ}
		\rdg[wit={P11}]{randhra}
		}
	\app{\lem[wit={N3,G4,Epsi,Zeta2,J10,Pi,Jyo}]{nirodhayet}% +C7
		\rdg[wit={J5}]{niyojayet}% +F
		\rdg[wit={Gamma,Delta}]{nirundhayet}
		}//\versenr}\label{jnatva}
		\orgvnr{58}\\!}
\end{tlg}
%</vs58>
\commcite\newpage
% V15 totally different: suṣumnāṃtagataṃ mārgaṃ vāyuṃ kṛtvā ca madhyagaṃ/ samāvasthā sthito yogī prāṇaraṃdhraṃ nirodhayet//
% K1 kṛtvā tāv iṃdavasthāne prāṇa°
% B2 kṛtvā tāvad bindusthairyaṃ prāṇa°
% J11 sthitvā sā caiṃdavasthāne ghrāṇa°


% between V3_4.175-176
%<*vs59a>
\begin{ava}[hp04_059a]
\app{\lem[nolem]{}
	\rdg[wit={Alpha,C6,V3}]{\only}}%
\app{\lem[wit={N3,G4,C6}]{tathā ca vasiṣṭhaḥ} % <<ḥ>> N3, vasiṣṭāḥ P7, ca vasiṣṭaḥ C6
		\rdg[wit={J5}]{tathā vaśiṣṭhavacanaṃ}
		\rdg[wit={V3}]{tatvāva || ☼ ||}/}
\end{ava}
%</vs59a>
%<*vs59>
\begin{tlg}[hp04_059]
\tl{\app{\lem[nolem]{}
	\rdg[wit={G11,Zeta2,J10,Jyo}]{\om}}%
\pada{iḍāyāṃ % idā° N23
	\app{\lem[wit={N3,J5,Gamma,Delta,P11,C6}]{piṅgalāyāṃ ca}
		\rdg[wit={V3}]{piṅgalāyāṃśca}}}
\pada{carataś candrabhāskarau/}\\+}
\tl{
\pada{candras tāmasa ity uktaḥ}
\pada{sūryo
	\app{\lem[wit={N3,J5,J7,Delta,Pi}]{rājasa}
		\rdg[wit={N23}]{\textapp{breaks off after} rā}
		} ucyate//\versenr}
		%\myfn{\getsiglum{N23} is lost after \textit{sūryo rā} in pāda d.}
		\orgvnr{59}\\!}
\end{tlg}
%</vs59>


%<*vs60>
\begin{tlg}[hp04_060]
\tl{\app{\lem[nolem]{}
	\rdg[wit={G5}]{\om}
	\rdg[wit={N23}]{\unavbl}}%
\pada{\app{\lem[alt={tāv eva ... sakalaṃ}]{\skp{tāv eva ... sakalaṃ}}%wit={N3,J5,J7,Delta1,Pi},
	\rdg[wit={G11}]{sūryaś candraḥ sadā dhatte}% om. G5; sūryaṃ candraṃ M3
	\rdg[wit={N19}]{sūryacandrau sadā dhatte}% dhattaḥ J11,E4,N26
	\rdg[wit={V15,Jyo}]{sūryācandramasau dhattaḥ}% +M1; sūryacandramasā yuktaṃ F
	\rdg[wit={J10}]{sūryācandramasau kṛtvā}
	}%
	\app{\lem[wit={N3,J5,J7,E2,P11,C6}]{tāv eva}% tau eva C6
		\rdg[wit={V19}]{tā eva}
		\rdg[wit={V3}]{tāṃve}} 
	\app{\lem[wit={N3,Delta1,P11,V3}]{dhattaḥ}% dhataḥ N3, dharttaḥ P11
		\rdg[wit={J5}]{dhanva}
		\rdg[wit={J7}]{dattaḥ}% +K3,C7
		\rdg[wit={C6}]{vahataḥ}}
	\app{\lem[wit={N3,J5,J7,Delta1,P11,V3}]{sakalaṃ}% ṃ om. V3
%		\rdg[wit={J5}]{śakalaṃ}
		\rdg[wit={C6}]{sarvaṃ}
		}} % nested!!
\pada{\app{\lem[nolem]{\skp{pāda b}}
	\rdg[wit={J10}]{\om}}%
	\app{\lem[wit={J7,Delta,G11,V15,P11,Jyo}]{kālaṃ}
		\rdg[wit={N3,J5,C6}]{kāla}
		\rdg[wit={N19}]{kālāṃ} % kālāṃśatri N19
		\rdg[wit={J10,V3}]{\om}}
	\app{\lem[wit={G11,Jyo}]{rātriṃdivātmakam}
		\rdg[wit={N3,J5,J7,Zeta2,P11,C6}]{rātridivātmakam}% °ka J5; +V6,F
		\rdg[wit={G4}]{rātriṃ divākaraṃ}
		\rdg[wit={Delta}]{rātrindinātmakaṃ}
%		\rdg[wit={N19}]{śatridivātmakaṃ}
		\rdg[wit={V3}]{rātridivātmakaṃ yogavit}
		}/}\\+}
\tl{
\pada{\app{\lem[nolem]{\skp{pāda c}}
	\rdg[wit={J10}]{\om}}%
	\app{\lem[wit={N3,J7,Delta,G11,V15,P11,Jyo}]{bhoktrī} % bho<<ktrī>> N3
		\rdg[wit={J5}]{bhoktu}
		\rdg[wit={G4}]{[bho]gī}
		\rdg[wit={N19}]{bhoktī}
		\rdg[wit={C6}]{bhoktṛ}
		\rdg[wit={V3}]{bhoktā}
		}
	suṣumṇā kālasya}
\pada{\app{\lem[nolem]{\skp{pāda d}}
	\rdg[wit={J10}]{\om}}%
	\app{\lem[wit={ceteri},alt={guhyam etad}]{guhyam eta\skp{d}}
		\rdg[wit={V19}]{guptam etad}
		\rdg[wit={E2}]{sattvam etad}% +C7
		}%
	\app{\lem[wit={ceteri},alt={udāhṛtam}]{\skm{d }udāhṛtam}
	\rdg[wit={J5}]{udīritaṃ}}//\versenr}
	\orgvnr{60}\\!}
\end{tlg}
%</vs60>

\medskip
\avacite{59a}
\commciterange{59}{59}\teimute{\vspace{-2.5ex}}
\commciterange{60}{60}\newpage


%<*vs61a>
\begin{ava}[hp04_061a]
\app{\lem[nolem]{}
	\rdg[wit={Epsi}]{\textapp{found after 3.93*7}}
	\rdg[wit={Zeta2,J10,Jyo}]{\om}
	}%
\app{\lem[wit={Alpha,V19,G5,C6,V3}]{tathā hi}% +K3,C7
	\rdg[wit={J7,E2}]{tathā}
	\rdg[wit={G11}]{athā hi}
	\rdg[wit={P11}]{tathāpi hi}
	}
\app{\lem[wit={N3,J5,Delta}]{saubhadraṃ nāma}% rāma J5
	\rdg[wit={J7}]{saubhadranāmā}
	\rdg[wit={G11}]{sobhadrā nāma}
	\rdg[wit={G5}]{saubhadranāmaś ca}
	\rdg[wit={P11}]{saubhadreryān nāma}
	\rdg[wit={C6}]{saubhadreyanāma}
	\rdg[wit={V3}]{saubhadreyaṃ nāma}
	}
\app{\lem[wit={N3,Delta,Epsi,Pi}]{ślokacatuṣṭayam}% śloṣa P11
	\rdg[wit={J5}]{ślokam eva catuṣṭayaṃ}
	\rdg[wit={J7}]{ślokacatuṣṭayam āha}}/%
	\myfn{\getsiglum{G11,G5} have this set of verses as \manuref{3.93*7--10} in a different order (cf. the introduction p.\,\hpref{77}). Their readings are also reported here.}
\end{ava}
%</vs61a>
%<*vs61>
\begin{tlg}[hp04_061]
\tl{\app{\lem[nolem]{}
	\rdg[wit={N23}]{\unavbl}
	\rdg[wit={Epsi}]{\textapp{found as 3.93*8}}
	\rdg[wit={Zeta2,J10,Jyo}]{\om}
	}%
\pada{\app{\lem[wit={J5,J7,Delta,Epsi,Pi}]{ṣaṭcakraṃ}
	\rdg[wit={N3}]{ṣaḍraktaṃ}} ṣoḍaśādhāraṃ}
\pada{\app{\lem[wit={J7,Delta,Epsi,V3},alt={tridhā lakṣ(y)aṃ}]{tridhā lakṣyaṃ}% lakṣaṃ V3,J7,G11
		\rdg[wit={N3,J5}]{tridhā bhajyaṃ}% m. om. J5
%		\rdg[wit={C7}]{tridhā yuktaṃ}
		\rdg[wit={P11}]{tridhākṣa ca}
		\rdg[wit={C6}]{trilakṣyaṃ ca}} guṇatrayam/}\\+}
\tl{
\pada{\app{\lem[wit={N3,J5,G5,Pi}]{śeṣaṃ tu} % m vs n!
		\rdg[wit={J7,Delta}]{śeṣas tu}% +F
		\rdg[wit={G11}]{śeṣaṃ tat}}
	\app{\lem[wit={ceteri}]{grantha}% ceteri = N3,J5,Pi,J7,V19,K3; gratha N3
		\rdg[wit={G11,C6}]{granthi}}% +C7
	\app{\lem[wit={N3,Epsi,Pi}]{vistāraṃ}
		\rdg[wit={J5}]{vistāra}
		\rdg[wit={J7,Delta}]{vistāras}% +F
		}}
\pada{\app{\lem[wit={N3,J5,J7,V19,Epsi,P11,V3}]{trikūṭaṃ}% °kuṭaṃ G11
		\rdg[wit={E2}]{trirūpaṃ}% +K3,C7
		\rdg[wit={C6}]{trikoṭi}
		} paramaṃ padam//\versenr}
	\orgvnr{61}\\!}
\end{tlg}
%</vs61>

\avacite{61a}
\commcite\newpage


%<*vs62>
\begin{tlg}[hp04_062]
\tl{\app{\lem[nolem]{}
	\rdg[wit={N23}]{\unavbl}
	\rdg[wit={J7,Zeta2,J10,C6,Jyo}]{\om}
	\rdg[wit={Epsi}]{\textapp{found as 3.93*7}}
	}%
\pada{kuṇḍalī kuṭilākārā} % kuṭulākāraṃ G4, kuḍalākārā J5, °kāra G5
\pada{sarpavat parikīrtitā/}\\+} % sarvasthat G4, sarvavat G11
\tl{
\pada{sā śakti\app{\lem[wit={N3,J5,Epsi,V3},alt={cālitā}]{\skm{ś }cālitā}
	\rdg[wit={G4}]{cāri\,..}
	\rdg[wit={V19}]{kīlitā}% +K3
	\rdg[wit={E2}]{kelitā}% +C7
	\rdg[wit={P11}]{calitā}
	} yena}
\pada{sa \app{\lem[wit={Delta,Epsi}]{mukto}% +G7 (eye-skip to the next verse?)
	\rdg[wit={N3,J5,P11,V3}]{yogī}} % +F ##
	nātra saṃśayaḥ//\versenr} % śaṃsayaḥ V19
	\orgvnr{62}\\!} % +B2
\end{tlg}
%</vs62>
\commcite%\newpage
% V3 has a gap indicated before the following verse.


%<*vs63>
\begin{tlg}[hp04_063]
\tl{\app{\lem[nolem]{}
	\rdg[wit={N23}]{\unavbl}
	\rdg[wit={J7,Delta,Zeta2,J10,Jyo}]{\om}
	\rdg[wit={Epsi}]{\textapp{found as 3.93*9}}}%
\pada{\app{\lem[wit={ceteri}]{yadā}\rdg[wit={G5}]{yathā}} 
	\app{\lem[wit={ceteri},alt={kūṭaṃ tri°}]{kūṭaṃ tri\skp{°}}
		\rdg[wit={C6}]{kūṭasti}}kūṭasthaṃ} % 
\pada{cittaṃ  % citta V3
	\app{\lem[wit={N3}]{citraṃ}
		\rdg[wit={J5}]{cittaṃ}
		\rdg[wit={Epsi}]{yatra}
		\rdg[wit={Pi}]{tatra}% +F
		} 
	\app{\lem[wit={ceteri}]{nirantaram}
		\rdg[wit={Epsi}]{nirajñanaṃ}}/}\\+}
\tl{
\pada{\app{\lem[wit={ceteri}]{kuṇḍalyās tu} % kuṃḍilyās P11
		\rdg[wit={G11}]{kuṇḍalyāpta}
		\rdg[wit={G5}]{kuṇḍalinyāḥ}}
	\app{\lem[wit={N3,J5,Epsi,P11,V3},post=\texteng{(°na˟ \getsiglum{N3})}]{prayogeṇa}%  + + geṇa G4, °na G11
		\rdg[wit={C6}]{prabodhena}}}
\pada{sa mukto nātra saṃśayaḥ//\versenr} 
	\orgvnr{63}\\!}
\end{tlg}
%</vs63>
\commcite\newpage


%<*vs64>
\begin{tlg}[hp04_064]
\tl{\app{\lem[nolem]{}
	\rdg{\texteng{=\,3.93*10}}
	\rdg[wit={N23}]{\unavbl}}%
\pada{\app{\lem[wit={N3,J5,J7,Delta,Pi,Jyo}]{dvāsaptatisahasrāṇi}
		\rdg[wit={G4,Epsi,Zeta2},alt={dvisaptati°}]{dvisaptatisahasrāṇi}% +F
		\rdg[wit={J10}]{\om}}} % °śrāṇi V3,V19
\pada{\app{\lem[wit={Alpha,J7,Epsi,V15,Pi,Jyo},post=\texteng{(nāḍi° \getsiglum{J5,P11})}]{nāḍīdvārāṇi}% nāḍi J5,P11
		\rdg[wit={V19}]{nāḍīnāṃdeda}
		\rdg[wit={E2}]{nāḍīnāṃ deha}% +K3,C7
		\rdg[wit={N19}]{nāḍīdvāre ca}
		\rdg[wit={J10}]{datvā kārāpi}}\marmas
	\app{\lem[wit={ceteri}]{pañjare}
		\rdg[wit={N3}]{paṃkaje}
		\rdg[wit={G4}]{maṃjarī}}/}\\+}
\tl{
\pada{suṣumṇā śāmbhavī śaktiḥ} % sāṃbhavī J10; śakti V19,V3
\pada{\app{\lem[wit={N3,E2,Epsi,N19,Pi,Jyo}]{śeṣās tv eva}% aiva P11; +C7
		\rdg[wit={J5}]{śeṣāsvevaṃ}
		\rdg[wit={J7,V19,V15}]{śeṣāś caiva}
		\rdg[wit={J10}]{śeṣās tv evaṃ}
		}
	\app{\lem[wit={ceteri}]{nirarthakāḥ} % nina° J10; °kā J5,V3
		\rdg[wit={N19}]{nivarttakāḥ}
		}//\versenr}%
	\myfn{\getsiglum{G11} includes this passage both in chapter three (3.93*10) and here, while \getsiglum{G5} includes it only in chapter three.}
	\orgvnr{64}\\!}
\end{tlg}
%</vs64>
\commcite%\newpage


%<*vs65>
\begin{tlg}[hp04_065]
\tl{\app{\lem[nolem]{}
	\rdg[wit={N23}]{\unavbl}
	\rdg[wit={E2}]{\om}}%
\pada{vāyuḥ % vāyu P11,V19
	\app{\lem[wit={N3,J5,Epsi,Zeta2,J10,C6,Jyo}]{paricito}
		\rdg[wit={V3}]{paricipta}
		\rdg[wit={J7}]{sa parito}
		\rdg[wit={V19}]{saṃparito}% +C7
		\rdg[wit={P11}]{parivṛtto}}
	\app{\lem[wit={N3,J7,V19,G11,Zeta2,P11,C6},alt={yatnād}]{yatnā\skp{d}}% yannād P11
		\rdg[wit={J5,G5,J10,Jyo}]{yasmād}% +F
%		\rdg[wit={C7}]{yadvad}
		\rdg[wit={V3}]{nādād}}}%
\pada{\app{\lem[wit={V19,Epsi,Zeta2,J10,Pi,Jyo},alt={agninā}]{\skm{d }agninā} % agnidā N19
		\rdg[wit={N3}]{yaṣṭinā}
		\rdg[wit={J5}]{yadasthā}
		\rdg[wit={J7}]{ṛgvinā}
		} saha
	\app{\lem[wit={Epsi,Jyo}]{kuṇḍalīm}% +C7
		\rdg[wit={N3,J5,J7,V19,Zeta2,J10,Pi}]{kuṇḍalī}}/}\\+}
\tl{
\pada{\app{\lem[nolem]{\skp{pāda c}}
	\rdg[wit={J10}]{\om}}%
	bodhayitvā suṣumṇāyāṃ} % °yā N3
\pada{\app{\lem[nolem]{\skp{pāda d}}
	\rdg[wit={J10}]{\om}}%
\app{\lem[wit={ceteri},alt={praviśed}]{praviśe\skp{d}}
		\rdg[wit={J5}]{praviśad}
		\rdg[wit={G5}]{praviśyad}
		\rdg[wit={J10}]{\om}
		\rdg[wit={V3}]{praveśad}
		}%
	\app{\lem[wit={G4,Epsi,V15,Pi,Jyo},alt={anirodhataḥ}]{\skm{d }anirodhataḥ} % °ta P11,V3; +M1
		\rdg[wit={N3,J5,J7,V19}]{avirodhataḥ} % aviśeṣataḥ V19ac; +C7
		\rdg[wit={N19}]{atirodhataḥ}
		\rdg[wit={J10}]{\om}}//\versenr}
	\orgvnr{65}\\!}
\end{tlg}
%</vs65>
\commcite\newpage


%<*vs66>
\begin{tlg}[hp04_066]
\tl{\app{\lem[nolem]{}
	\rdg[wit={N23}]{\unavbl}}%
\pada{\app{\lem[nolem]{\skp{pāda a}}
	\rdg[wit={J10}]{\om}}%
suṣumṇā\app{\lem[wit={G4,J7,E2,G11,C6,V3,Jyo}]{vāhini}
		\rdg[wit={N3,J5,G5,Zeta2,P11}]{vāhinī}
		\rdg[wit={V19}]{hini}
		} 
		prāṇe} % prāṇo J5
\pada{\app{\lem[nolem]{\skp{pāda b}}
	\rdg[wit={J10}]{\om}}%
\app{\lem[wit={G4,J7,Delta,Epsi,V15,C6,V3,Jyo}]{sidhyaty eva}% °teva G11
		\rdg[wit={N3}]{siddhyety eva}
		\rdg[wit={J5}]{siddhity eva}
%		\rdg[wit={C7}]{siddhyatīva}
		\rdg[wit={N19,P11}]{siddhaty eva}
		} 
	manonmanī/}
	\\+}
\tl{
\pada{\app{\lem[wit={Alpha,J7,Pi},alt={anyathā vividhā°}]{anyathā vividhā\skp{°}}% anātha J5
		\rdg[wit={Delta}]{anye ye vividhā}
%		\rdg[wit={C7}]{anye ca vividhā}
		\rdg[wit={G11}]{prāṇe suṣumnāṃ saṃ}
		\rdg[wit={G5,Zeta2}]{anyathā tv itare}
		\rdg[wit={J10}]{atha cittāntare}
		\rdg[wit={Jyo}]{anyathā tv itarā}
		}%
	\app{\lem[wit={N3,E2,C6,Jyo},alt={°bhyāsāḥ}]{\skp{°}bhyāsāḥ}% bhyāsā<<ḥ>> C7
		\rdg[wit={J5,G5,N19,P11}]{bhyāsāt}
		\rdg[wit={G4,J7,V19,V3}]{bhyāsā}
		\rdg[wit={G11}]{prāpte}
		\rdg[wit={V15,J10}]{bhyāsa}
		}}
\pada{\app{\lem[wit={N3,J5,J7,Epsi,Pi,Jyo}]{prayāsāyaiva}% +C7
		\rdg[wit={V19}]{prāyāsāś caiva}
		\rdg[wit={E2}]{prayāsāyai}
		\rdg[wit={N19}]{prayāsā eka}
		\rdg[wit={V15}]{prayāsā eva}
		\rdg[wit={J10}]{pratyāśā jīva}
		}
	\app{\lem[wit={ceteri}]{yoginām}
		\rdg[wit={J5,J10,V3}]{yoginā}
		\rdg[wit={N19}]{yoginī}}//\versenr}
	\orgvnr{66}\\!}
\end{tlg}
%</vs66>
\commcite%\newpage


%<*vs67>
\begin{tlg}[hp04_067]
\tl{\app{\lem[nolem]{}
	\rdg[wit={N23}]{\unavbl}}%
\pada{pavano badhyate 
	\app{\lem[wit={ceteri}]{yena}
	\rdg[wit={J5}]{deva}
	}}
\pada{\app{\lem[wit={ceteri}]{manas tenaiva badhyate} % <va>dhyate N19
		\rdg[wit={J10}]{tenaiva badhyate manaḥ}}/}\\+}
\tl{
\pada{\app{\lem[nolem]{\skp{pāda c}}
	\rdg[wit={J5,J7,G5,J10}]{\om}}%
\app{\lem[wit={N3,G11,Zeta2,P11,V3,Jyo}]{manaś ca}
		\rdg[wit={Delta}]{manas tu}% +F
		\rdg[wit={C6}]{manas tad}
		} 
		badhyate yena}
\pada{\app{\lem[nolem]{\skp{pāda d}}
	\rdg[wit={J5,J7,G5,J10}]{\om}}%
\app{\lem[wit={ceteri}]{pavanas tena} % <pa>vanas V19
		\rdg[wit={V3}]{pavanamana}
		} 
		badhyate//\versenr}
	\orgvnr{67}\\!}
\end{tlg}
%</vs67>
\commcite\newpage

%<*vs68>
\begin{tlg}[hp04_068]
\tl{\app{\lem[nolem]{}
	\rdg[wit={N23}]{\unavbl}
	\rdg[wit={V19}]{\textapp{found after \ref{dugdha}}}}%
\pada{\app{\lem[wit={ceteri}]{hetu}
%		\rdg[wit={C7}]{deha}
		\rdg[wit={J5}]{heta}
		\rdg[wit={G4}]{eta}
		\rdg[wit={G5}]{ajī}}%
	\app{\lem[wit={N3,G4,E2,J10,Jyo}]{dvayaṃ tu}% +C7
		\rdg[wit={J5}]{dvāv api}
		\rdg[wit={J7,G11,P11,V3}]{dvayaṃ hi}
		\rdg[wit={V19,C6}]{dvayaṃ ca}
		\rdg[wit={G5}]{vayat hi}
		\rdg[wit={Zeta2}]{dvayasya}
		}
	\app{\lem[wit={ceteri}]{cittasya}
		\rdg[wit={J7,Delta}]{manaso}}}
\pada{vāsanā ca samīraṇaḥ/}\\+} % ḥ om. V3; vāsanācca samīraṇa<<ḥ>> N3, vāsanaṃ ca samīraṇaṃ J5
\tl{
\pada{tayo\app{\lem[wit={ceteri},alt={vinaṣṭa ekasmin}]{\skm{r }vinaṣṭa ekasmi}% tayo J5,J7; ekasminn J7, ekasmi V19, yeka° G5
		\rdg[wit={G11}]{vinaṣṭa etasmin}
		\rdg[wit={C6}]{vinaṣṭas tv ekaś ca hy}}}%n 
\pada{\app{\lem[wit={N3},alt={drutaṃ dvāv api naśyataḥ},postwit=\texteng{(druttaṃ)}]{\skm{n }drutaṃ dvāv api naśyataḥ}% +M3
%		\rdg[wit={N3}]{druttaṃ dvāv api naśyataḥ}
		\rdg[wit={J5}]{dṛtaṃ vāvati nasyataḥ}
		\rdg[wit={G4}]{dhṛtaṃ dvāv api naśyataḥ}
		\rdg[wit={J7,E2,J10,C6}]{ubhāv api vinaśyataḥ}
		\rdg[wit={V19}]{svabhāvo pi vinaśyataḥ}
		\rdg[wit={G11}]{nṛtaṃ dvāv api naśyati}
		\rdg[wit={G5}]{yatta dvāv api naśyataḥ}
		\rdg[wit={Zeta2,P11,V3,Jyo}]{tau dvāv api vinaśyataḥ} % °ta P11,V3, °tiḥ N19
		}//\versenr}%
	\orgvnr{68}\\!}
\end{tlg}
%</vs68>
\commcite\newpage


%<*vs69>
\begin{tlg}[hp04_069]
\tl{\app{\lem[nolem]{}
	\rdg[wit={N23}]{\unavbl}
	\rdg[wit={V19}]{\textapp{found after \ref{dugdha} together with the previous verse}}}%
\pada{\app{\lem[nolem]{\skp{pāda a}}
	\rdg[wit={J10}]{\om}}%
mano yatra \app{\lem[wit={ceteri}]{vilīyeta} % °līyena N19
		\rdg[wit={V3}]{vilīyate}
		}}
\pada{\app{\lem[nolem]{\skp{pāda b}}
	\rdg[wit={J10}]{\om}}%
\app{\lem[wit={ceteri},alt={pavanas}]{pavana\skp{s}}
		\rdg[wit={Epsi,Zeta2}]{mārutas}
		}s tatra līyate%
\app{\lem[alt={\post līyate}]{\skp{\post līyate}}
	\rdg[wit={V15}]{ekatra[m]iśritau \add}}/}\\+}
\tl{
\pada{\app{\lem[nolem]{\skp{pāda c}}
	\rdg[wit={J5,Zeta2}]{\om}}%
\app{\lem[wit={N3,J7,C6,Jyo}]{pavano līyate yatra}
		\rdg[wit={Delta}]{pavano yatra līyeta}
		\rdg[wit={Epsi}]{māruto yatra līyeta}
		\rdg[wit={J10}]{yatraiva līyate vāyur}
		\rdg[wit={P11,V3}]{pavano yatra līyate}% līyaṃte P11
		}}
\pada{\app{\lem[nolem]{\skp{pāda d}}
	\rdg[wit={J5,Zeta2}]{\om}}%
mana\app{\lem[wit={N3,Delta,Epsi,J10,Pi},alt={tatraiva līyate}]{\skm{s }tatraiva līyate}% +N2
		\rdg[wit={J7,Jyo}]{tatra vilīyate}% +V6
		}//\versenr}
	\orgvnr{69}\\!}
\end{tlg}
%</vs69>
\commcite%\newpage


%<*vs70>
\begin{tlg}[hp04_070]
\tl{\app{\lem[nolem]{}
	\rdg[wit={N23}]{\unavbl}}%
\pada{dugdhāmbuvat saṃmilitau % dugdhābu° N3; saṃmilaṃ J5, <<li>> C6, °tāv J10,Jyo, °to N19
	\app{\lem[wit={N3,J5,Epsi,Zeta2,Pi}]{sadaiva}
		\rdg[wit={G4}]{sadeva}
		\rdg[wit={J7,Delta}]{tathaiva}
		\rdg[wit={J10,Jyo}]{ubhau tau}}}\\+}
\tl{
\pada{tulyakriyau % kriyo N19, tulyasya kriyā J5, tvalya C6
	\app{\lem[wit={ceteri}]{mānasamārutau} % māruto N19; 2 x mānasa J10
		\rdg[wit={Epsi,P11,C6}]{mārutamānasau}
		\rdg[wit={V3}]{\illeg}}
	\app{\lem[wit={N3,G4,Epsi,Zeta2,J10,P11,Jyo}]{hi}
		\rdg[wit={J5,J7,Delta,C6}]{ca}
		\rdg[wit={V3}]{\illeg}
		}/}\\+}
\tl{
\pada{\app{\lem[wit={ceteri},alt={yāvan manas}]{yāvan mana\skp{s}}
		\rdg[wit={J10,Jyo}]{yato marut}}%
	\app{\lem[wit={ceteri},alt={tatra}]{\skm{s }tatra}
		\rdg[wit={J5}]{caiva}}
	\app{\lem[wit={ceteri},alt={marut}]{maru\skp{t}} % murut G11, marat V15
		\rdg[wit={J10,Jyo}]{manaḥ}
		\rdg[wit={C6}]{\_\,sat}}%
	\app{\lem[wit={ceteri},alt={pravṛttir}]{\skm{t}pravṛtti\skp{r}} % vṛtti J7
		\rdg[wit={N19}]{pravṛddhitti}
		\rdg[wit={C6}]{pravṛtta}% varti P7
		}-}\\+}
\tl{
\pada{\app{\lem[nolem]{\skp{pāda c}}
		\rdg[wit={Zeta2}]{\om}}%
	\app{\lem[wit={Alpha,J7,Delta,Epsi,Pi},alt={yāvan}]{\skm{r }yāva\skp{n}}
		\rdg[wit={J10,Jyo}]{yato}
		}%
	\app{\lem[wit={N3,J5,J7,Delta,Epsi,P11,C6},alt={maruc cāpi}]{\skm{n }maruc cāpi}
		\rdg[wit={J10,Jyo}]{manas tatra}
		\rdg[wit={V3}]{marut tatra}
		}
	\app{\lem[wit={N3pc,J7,E2,Epsi,C6,V3}]{manaḥ}
		\rdg[wit={N3ac,J5,V19,P11}]{mana}
		\rdg[wit={J10,Jyo}]{marut}% or: marun
		}%
	\app{\lem[wit={N3,J7,Delta,Epsi,P11,V3,Jyo}]{pravṛttiḥ} % °vartiḥ P7
		\rdg[wit={J5}]{pravittato}
		\rdg[wit={J10}]{nivṛttiḥ}
		\rdg[wit={C6}]{pravṛttaḥ}
		}//\versenr}\label{dugdha}
	\orgvnr{70}\\!}
\end{tlg}
%</vs70>
\commcite\newpage


%<*vs71>
\begin{tlg}[hp04_071]
\tl{\app{\lem[nolem]{}
	\rdg[wit={N23}]{\unavbl}
	\rdg[wit={V19}]{\textapp{\emph{pāda}s ab and cd are transposed}}}%
\pada{\app{\lem[wit={ceteri}]{tatraika}% tatr<<aika>> N3
		\rdg[wit={N3ac}]{tatra}
		\rdg[wit={G5}]{tathaika}
		\rdg[wit={Zeta2}]{atraika}
		\rdg[wit={J10}]{ekasya}}nāśād aparasya % nā<śā>d J5,N19, nāsād J10
	\app{\lem[wit={N3,J5,J7,E2,G5,Zeta2,J10,C6,Jyo},alt={nāśa(ḥ)}]{nāśa\skp{(ḥ)}}
%		\rdg[wit={J5,Zeta2}]{nāśaḥ}
		\rdg[wit={V19}]{nāśam}
		\rdg[wit={G11}]{nāśā}
%		\rdg[wit={J10}]{nāśas}
		\rdg[wit={P11}]{nāśe}
		\rdg[wit={V3}]{nāśo}
		}}\\+}
\tl{
\pada{\app{\lem[wit={N3,J5,J7,G5,N19,P11,Jyo},alt={ekapravṛtter}]{ekapravṛtte\skp{r}}
		\rdg[wit={Delta,G11,V15}]{ekapravṛttāv}
		\rdg[wit={J10}]{tatraikavṛtter}
		\rdg[wit={C6}]{ekapravṛtte}
		\rdg[wit={V3}]{e\,..\,..\,..\,..}
		}%
	\app{\lem[wit={ceteri},alt={aparapravṛttiḥ}]{\skm{r }aparapravṛttiḥ} % °ttir N19, °tt<<i>>ḥ N3
		\rdg[wit={G11}]{itarapravṛttiḥ}
		\rdg[wit={J10}]{aparasya vṛttiḥ}
		\rdg[wit={C6}]{ca parapravṛttiḥ}
		\rdg[wit={V3}]{..\,..\,..\,..\,..\,ttiḥ}
		}%
	\app{\lem[alt={\post pravṛttiḥ}]{\skp{\post pravṛttiḥ}}
	\rdg[wit={V15}]{ekasya nā<śā>d aparasya nāśaḥ \textapp{(alternative reading for \emph{pāda} a)} \add}}/}%
	\\+}
\tl{
\pada{\app{\lem[wit={N3,E2,G5,V15,P11,Jyo},alt={adhvastayoś/r}]{adhvastayo\skp{ś/r}}
		\rdg[wit={J5}]{adhastasar}
		\rdg[wit={J7}]{adhyastayor}
		\rdg[wit={V19,J10,C6}]{adhastayoś/r}
		%\rdg[wit={E2,V15}]{adhvastayor}
		\rdg[wit={G11}]{adhvaścayoś}
		\rdg[wit={N19}]{addhastayoś}
		%\rdg[wit={J10,C6}]{adhastayoś}
		\rdg[wit={V3}]{atastayoś}
		}%
	\app{\lem[wit={Alpha,Epsi,N19,J10,Pi,Jyo},alt={cendriya}]{\skm{ś }cendriya}% veddriya J5, caiṃdriya P11, cenniya G11, caidriya N19
		\rdg[wit={J7,Delta,V15}]{indriya}}varga% 
	\app{\lem[wit={N3,G4},alt={buddhir}]{buddhi\skp{r}}
		\rdg[wit={J5,C6}]{śuddhir}
		\rdg[wit={J7,E2}]{vṛddhir}% +C7
		\rdg[wit={V19,Epsi,Zeta2,J10,Jyo}]{vṛttiḥ}% vṛttir V19, vṛtti V15
		\rdg[wit={P11}]{baṃdhir}
		\rdg[wit={V3}]{vudhir}
		}\marma-}\\+}
\tl{
\pada{\app{\lem[wit={N3,G4,Delta,G5,V15,Pi},alt={vidhvastayor}]{\skm{r }vidhvastayo\skp{r}}
		\rdg[wit={J5}]{adhastayor}
		\rdg[wit={J7}]{vivṛddhayor}% °yo J7
		\rdg[wit={G11}]{nidhvastayo}
		\rdg[wit={N19}]{addhvastayor}
		\rdg[wit={J10}]{vijñātayor}
		% vṛttiḥ|(gap for about one hemistich)raddhvasta° N19
		\rdg[wit={Jyo}]{pradhvastayor}
		}%
	\app{\lem[wit={ceteri},alt={mokṣapadasya}]{\skm{r }mokṣapadasya}
		\rdg[wit={J7},alt={°pradasya}]{mokṣapradasya}
%		\rdg[wit={C7},alt={°pathasya}]{mokṣapathasya}
		}
		siddhiḥ//\versenr}\label{tatraika}
	\orgvnr{71}\\!} % siddhi J5,P11,V19
\end{tlg}
%</vs71>
\commcite\newpage


%====== Passage B and yāvanna-stanza collated with more mss (V19,N19,J10)

%<*vs72>
\begin{tlg}[hp04_072]
\tl{\app{\lem[nolem]{}
	\rdg[wit={N23}]{\unavbl}
	\rdg[wit={Jyo}]{\om}
	}%
	\pada{\app{\lem[wit={ceteri}]{vāyu}
		\rdg[wit={V19,V15}]{vāyur}}%
	\app{\lem[wit={Epsi}]{mārge tv asaṃcāre}
		\rdg[wit={Alpha,J7,Pi}]{mārgeṇa saṃcāre}
		\rdg[wit={Delta}]{mārgeṇa saṃcārī}
		\rdg[wit={N19}]{mārge tha saṃcāre}
		\rdg[wit={V15}]{mārge py asaṃcāre}
		\rdg[wit={J10}]{mārge ca saṃcāre}
		}}
\pada{\app{\lem[wit={N3,J7,Delta,V3}]{sakalāṃ}
		\rdg[wit={J5,Epsi,Zeta2,C6}]{sakalaṃ}
		\rdg[wit={G4}]{sakalā}
		\rdg[wit={J10}]{sa phalaṃ}
		\rdg[wit={P11}]{saṃkalpāt}
		}
	\app{\lem[wit={Alpha,Epsi,V15,J10,P11}]{labhate}
		\rdg[wit={J7,Delta}]{bhramate}
		\rdg[wit={N19,C6}]{labhyate}
		\rdg[wit={V3}]{carate}
		}\marmas
	\app{\lem[wit={N3,G4,J7,Delta,G11,P11}]{mahīm}
		\rdg[wit={J5}]{mahiḥ}
		\rdg[wit={G5,Zeta2}]{mahaḥ}
		\rdg[wit={J10}]{mahān}
		\rdg[wit={C6,V3}]{mahī}
		}/}\\+}
\tl{
\pada{\app{\lem[wit={Alpha,Delta,Epsi},post=\texteng{(tathā<<ṣṭa>> \getsiglum{N3})}]{tathāṣṭa}
		\rdg[wit={J7}]{na tathā}
		\rdg[wit={Zeta2,J10}]{tato 'ṣṭa}
		\rdg[wit={P11}]{aṣṭadhā}
		\rdg[wit={C6,V3}]{athāṣṭa}
		}guṇam aiśvaryaṃ} % ṃ om. J7
\pada{\app{\lem[wit={N3,G4,J7,Delta,Pi}]{satyaṃ satyaṃ varānane}
		\rdg[wit={Epsi,Zeta2,J10}]{satyam ity āha śaṃkaraḥ}
		\rdg[wit={J5}]{labhate sakalān varān}}//\versenr}
	\label{vayumargena}
	\orgvnr{72}\\!}
\end{tlg}
%</vs72>
\commcite\newpage


%<*vs73a>
\begin{ava}[hp04_073a]
\app{\lem[nolem]{}
	\rdg[wit={Epsi,Zeta2,J10,V3,Jyo}]{\om}}%
\app{\lem[wit={N3,P11,C6}]{tathā}
	\rdg[wit={J5}]{tathā ca}
	\rdg[wit={G4}]{tathāha}
	\rdg[wit={J7,Delta}]{\om}} 
	viśvarūpācāryaḥ/ % °yāḥ J7
\end{ava}
%</vs73a>
%<*vs73>
\begin{tlg}[hp04_073]
\tl{\app{\lem[nolem]{}
		\rdg[wit={N23}]{\unavbl}
		\rdg[wit={Zeta2,J10,V3}]{\om}
		}%
	\pada{\app{\lem[wit={Alpha,Delta,C6,Jyo}]{yadā saṃkṣīyate}
		\rdg[wit={J7,P11}]{yadā sa kṣīyate}
		\rdg[wit={G11}]{\om}
		\rdg[wit={G5}]{yadā saṃlīyate}} 
		prāṇo} % =P7; prāṇaṃ J5, prāṇe C6
\pada{mānasaṃ \app{\lem[wit={Alpha,Epsi,P11,C6}]{ca vilīyate}% vīlī° J5; +C7
		\rdg[wit={J7,Jyo}]{ca pralīyate}
		\rdg[wit={V19}]{pravilīyate}
		\rdg[wit={E2}]{saṃpralīyate}}/}\\+}
\tl{
\pada{\app{\lem[wit={ceteri}]{tadā}
		\rdg[wit={Epsi}]{tayoḥ}} % tadvā N3ac
	\app{\lem[wit={ceteri}]{samarasatvaṃ}
		\rdg[wit={J5},post={\unm}]{samarasaikatvaṃ}} % = VM6! J5 is hypermetrical, since it has yat too.
	\app{\lem[wit={N3,J5,J7,E2,Epsi,C6},alt={yat}]{ya\skp{t}}% +F
		\rdg[wit={G4,V19}]{yaḥ}
%		\rdg[wit={C7}]{hi}
		\rdg[wit={P11,Jyo}]{ca}}}%
\pada{\app{\lem[wit={N3,G4,J7,V19,G11,C6},alt={samādhiḥ so 'bhidhīyate}]{\skm{t }samādhiḥ so'bhidhīyate} % samādhi G4,V19
		\rdg[wit={J5}]{samādhiś ca vilīyate}
		\rdg[wit={E2,G5}]{samādhiḥ sābhidhīyate}
		\rdg[wit={P11}]{samādhī sau bhidhīyate}
		\rdg[wit={Jyo}]{samādhir abhidhīyate}
		}//\versenr}
		\label{visvarupa} 
	\orgvnr{73}\\!}
\end{tlg}
%</vs73>

\avacite{73a}
\commcite\newpage


%<*vs74>
\begin{tlg}[hp04_074]
\tl{\app{\lem[nolem]{}
	\rdg[wit={N23}]{\unavbl}
	\rdg[wit={V3}]{\om}
	}%
\pada{\app{\lem[wit={N3pc,J7,Delta,C6,Jyo}]{manaḥ}
		\rdg[wit={N3ac,J5,G4,Epsi,Zeta2,J10,P11}]{mana}}%
\app{\lem[wit={N3,J5,J7,Epsi,N19,J10,P11,C6,Jyo}]{sthairye}
		\rdg[wit={G4,V19}]{sthairya}
		\rdg[wit={E2}]{sthairyāt}
		\rdg[wit={V15}]{sthairyaḥ}}
	\app{\lem[wit={ceteri}]{sthiro}
		\rdg[wit={G4,G11,V15}]{sthito}} vāyu}% vāyaḥ G11
\pada{s tato \app{\lem[wit={N3pc,G4,J7,E2,V15,Jyo}]{binduḥ}
		\rdg[wit={N3ac,J5,V19,Epsi,N19,J10,P11,C6}]{bindu}}
	\app{\lem[wit={ceteri}]{sthiro bhavet}
%		\rdg[wit={C7}]{sthito bhavet}
		\rdg[wit={G4}]{tato layaḥ}}/}\\+}
\tl{
\pada{\app{\lem[wit={ceteri}]{bindu}
		\rdg[wit={J7}]{binduḥ}}%
	\app{\lem[wit={N3,E2,G5,P11,C6},alt={sthairyodayāt}]{sthairyodayā\skp{t}}% +C7
		\rdg[wit={J5}]{sthairyo sthiro}
		\rdg[wit={G4,N19}]{sthairyodayā}
		\rdg[wit={J7}]{sthairyād athā}
		\rdg[wit={V19}]{sthairyād yathā}
		\rdg[wit={G11}]{sthairyoyadāt}
		\rdg[wit={V15}]{sthairye dayā}
		\rdg[wit={J10}]{sthairyād dayā}
		\rdg[wit={Jyo}]{sthairyāt sadā}
		}%
	\app{\lem[wit={N3,P11},alt={putra}]{\skm{t }putra}
		\rdg[wit={J5}]{vāyu}
		\rdg[wit={G4}]{tatra}
		\rdg[wit={J7}]{panna}% pannaṃ V6,YCM; panna J7,N2
		\rdg[wit={V19},alt={\lacuna}]{\_\,\_\,\_}
		\rdg[wit={E2,Zeta2}]{satyaṃ}% +Ten,C7,G3
		\rdg[wit={Epsi}]{samyak}% +M3
		\rdg[wit={J10,Jyo}]{satvaṃ}
		\rdg[wit={C6}]{mūtra}
		}}
\pada{piṇḍasthairyaṃ prajāyate//\versenr} % ddhaḍa J5; sthairya P11
	\label{manahsthairye}
	\orgvnr{74}\\!}
\end{tlg}
%</vs74>
\commcite\newpage


%%%%%%%%%%%%%%%%%%%%%%%
%<*vs75>
\begin{tlg}[hp04_075]
\tl{\app{\lem[nolem]{}
	\rdg[wit={N23}]{\unavbl}
	\rdg[wit={N19,Jyo}]{\om}}%
\pada{dṛṣṭiḥ sthirā yasya % dṛṣṭi N3ac,J5,G4,G5,P11,V19,V3,J10
	\app{\lem[wit={ceteri}]{vinaiva}% Alpha,G11,G5,V15,J10,Pi
%		\rdg[wit={C7}]{vinā ca}
		\rdg[wit={J7,Delta}]{vināpi}}
	\app{\lem[wit={N3,G4,G5,V15,Pi},alt={dṛśyād}]{dṛśyā\skp{d}}
		\rdg[wit={J5}]{dṛśyavān}
		\rdg[wit={J7,Delta,G11,J10}]{dṛśyaṃ}
		}-}\\+}
\tl{d
\pada{vāyuḥ sthiro yasya % vāyu N3ac,J5,G4,P11,J7,V19,G11,G5,V15,J10
	\app{\lem[wit={ceteri}]{vinā prayatnāt} % prayatnataḥ J10ac
		\rdg[wit={J7}]{vināpi yatnaṃ}}/}\\+}
\tl{
\pada{cittaṃ sthiraṃ yasya
	\app{\lem[wit={N3pc,G4,Epsi,V15,C6,V3},alt={vināvalambāt}]{vināvalambā\skp{t}}% °laṃbā C6
		\rdg[wit={N3ac}]{vināvalambanāt}
		\rdg[wit={J5,Delta}]{vināvalaṃbanaṃ}% +N2
		\rdg[wit={J7}]{vinā prayatnāt}
%		\rdg[wit={C7}]{vinā balaṃ ca}
		\rdg[wit={J10}]{vināvalaṃnaṃ}
		\rdg[wit={P11}]{vinā vilambāt}
		}-}\\+}
\tl{t
\pada{sa eva yogī % yeva J5,P11
	\app{\lem[wit={ceteri}]{sa guruḥ}
		\rdg[wit={J10}]{sadguruḥ}}
	\app{\lem[wit={ceteri}]{sa sevyaḥ} % sevya P11
		\rdg[wit={J7,V19}]{sa śiṣyaḥ}}//\versenr}
	\orgvnr{75}\\!}
\end{tlg}
%</vs75>
\commcite\newpage

%%%%%%%%%%

%<*vs76>
\begin{tlg}[hp04_076]
\tl{\app{\lem[nolem]{}
	\rdg[wit={N23}]{\unavbl}
	\rdg[wit={N19,Jyo}]{\om}
	}%
\pada{praveśe nirgame % praveśa G4; prathame niyame J5
	\app{\lem[wit={ceteri}]{vāme}
		\rdg[wit={G4}]{vāma}
		\rdg[wit={V15}]{cāpi}
		\rdg[wit={P11}]{vāpi}
		}}
\pada{dakṣiṇe \app{\lem[wit={Alpha,G11,P11}]{cordhvam apy adhaḥ}% adha G4
		\rdg[wit={J7,Delta}]{cordhvamadhyagaḥ} % cordha V19
%		\rdg[wit={C7}]{cordhvamadhyamaḥ}
		\rdg[wit={G5,V15,J10}]{cordhvamadhyataḥ}
		\rdg[wit={C6}]{cordhvage 'py adhaḥ}
		\rdg[wit={V3}]{tanirodhataḥ}
		}/}\\+}
\tl{
\pada{\app{\lem[nolem]{\skp{pāda c}}
		\rdg[wit={G5}]{\om}}%
	\app{\lem[wit={ceteri}]{na yasya}
		\rdg[wit={C6}]{layasya}}
	\app{\lem[wit={ceteri}]{vāyur vahati}% vahari G4
		\rdg[wit={V3}]{vahate vāyu}}}
\pada{\app{\lem[nolem]{\skp{pāda d}}
		\rdg[wit={G5}]{\om}}%
	sa mukto nātra saṃśayaḥ//\versenr} % yukto? V19
	\label{pravese}
	\orgvnr{76}\\!}
\end{tlg}
%</vs76>
\commcite%\newpage


%<*vs77>
\begin{tlg}[hp04_077]
% J7 breaks off at \textit{rājayo} in pāda c%
\tl{\app{\lem[nolem]{}
	\rdg[wit={N23}]{\unavbl}
	}%
\pada{sarve \app{\lem[wit={N3,J5,V15,J10,P11b,C6,V3,Jyo}]{haṭhalayopāyā}% °pāya J5
		\rdg[wit={J7,E2}]{layahaṭhābhyāsā} % +V6,°haṭha°J7
		\rdg[wit={V19}]{haṭhālayābhyāsā}
		\rdg[wit={G11}]{layahaṭhopāyā}
		\rdg[wit={G5}]{layaparopāyā}
		\rdg[wit={N19}]{haṭhalayoyāgā}
		\rdg[wit={P11a}]{haṭhalayā bhāvyā}
		}}
\pada{\app{\lem[wit={N3,J5,J7,Delta1,P11b,C6b,Jyo}]{rājayogasya siddhaye} % +J7
		\rdg[wit={Epsi,Zeta2,J10}]{rājayogāya kevalaṃ} % yogaya N19
		\rdg[wit={P11a}]{rājayogapadāvadhi} %
		\rdg[wit={C6a}]{°padāvadhiḥ} %
		\rdg[wit={V3a}]{°padāvadhiṃ}
		\rdg[wit={V3b}]{°phalāvadhi}
		}/}\\+}
\tl{
\pada{\app{\lem[wit={ceteri}]{rājayoga}
		\rdg[wit={G4}]{rajayogaṃ}
		\rdg[wit={J7},post=\textapp{(then lost)}]{rājayo\skp{ (then lost)}}
		\rdg[wit={E2}]{rājayoge}% +C7
		}%
	\app{\lem[wit={ceteri}]{samārūḍhaḥ}% rūḍha V15, rūḍhā V3
	\rdg[wit={J5}]{padaprāptaḥ}
	\rdg[wit={P11a,C6a,V3a}]{padaṃ prāpya}
	}}
\pada{\app{\lem[wit={ceteri}]{puruṣaḥ kālavañcakaḥ}
		\rdg[wit={P11a,C6a}]{jāyate 'sau nirañjanaḥ}
		\rdg[wit={V3a}]{jāyate so nirañjana}
		}//\versenr}\label{kalavancaka}% °caka V3
\myfn{This verse appears twice in the \textpi\ manuscripts. The first instance (a) is as equivalent of X4.116, and the second (b) is as the semi-final verse of this chapter (4.77 in the α recension). \mydelim 
\xexclude{After this verse, \getsiglum{V19,C7} (not \getsiglum{E2}) have two additional verses:\vspaceinnote\\
\devnote{\myvin iḍā bhagavatī gaṅgā piṅgalā yamunā nadī/
vijñeyā taddvayor madhye suṣumṇā ca}\,(\getsiglum{V19}; \devnote{tu} \getsiglum{C7}) \devnote{sarasvatī//} (cf. 3.94*1)\\
\devnote{\myvin triveṇīsaṃgamo yatra tīrtharājaḥ sa ucyate/
tatra snānaṃ prakurvīta} (\getsiglum{V19}; \devnote{tasmiṃs tīrthavare snātvā} \getsiglum{C7}) \devnote{sarvapāpaiḥ pramucyate//}
}%
\xinclude{After this verse, \getsiglum{E4} has the verses on \emph{kālajñāna}, \emph{videhamukti}, and \emph{kālavañcaka}. See Appendix B.}}
	\orgvnr{77}\\!} % ḥ om. V3
\end{tlg}
%</vs77>
\commcite\newpage


% \startaltrecension
% \begin{alttlg}[hp04_077_1]
% \tl{
% \pada{iḍā bhagavatī gaṅgā}
% \pada{piṅgalā
	% \app{\lem[wit={C7}]{yamunā}
	% \rdg[wit={V19}]{jamunā}} nadī/}\\+}
% \tl{
% \pada{\app{\lem[wit={C7}]{vijñeyā}
	% \rdg[wit={V19}]{vidheyā}} taddvayor madhye} % tadvayor V19
% \pada{suṣumṇā \app{\lem[wit={C7}]{tu} % so auch in L1
	% \rdg[wit={V19}]{ca}} sarasvatī//\versenr} 
	% \anm{cf. 3.94*1}\\!}
% \end{alttlg}

% \begin{alttlg}[hp04_077_2]
% \tl{
% \pada{triveṇīsaṃgamo yatra}
% \pada{tīrtharājaḥ sa ucyate/}\\+}
% \tl{
% \pada{\app{\lem[wit={C7}]{tasmiṃs tīrthavare snātvā} % +L1,K2
	% \rdg[wit={V19}]{tatra snānaṃ prakurvīta}
	% }}
% \pada{sarvapāpaiḥ pramucyate//\versenr}\\!}
% \end{alttlg}
% \endaltrecension


%<*vs78>
\begin{tlg}[hp04_078]
\tl{\app{\lem[nolem]{}
	\rdg[wit={Gamma}]{\unavbl}
	\rdg[wit={Zeta2,J10,Jyo}]{\om}}%
\pada{iti 
	\app{\lem[wit={Delta,Pi}]{tu}
		\rdg[wit={N3}]{<<tu>>}
		\rdg[wit={J5,Epsi}]{\om}% +F
		\rdg[wit={G4}]{śrī}}
	\app{\lem[wit={ceteri}]{sakalayoga}
		\rdg[wit={G11}]{sakalasuyoga}
		\rdg[wit={G5}]{sakala}}śāstra% 
	\app{\lem[wit={N3pc,E2,G5,C6}]{sindhoḥ}% +C7
		\rdg[wit={N3ac}]{siddhāḥ}
		\rdg[wit={J5}]{sindhauḥ}
		\rdg[wit={G4}]{\om}
		\rdg[wit={V19}]{sindhau}
		\rdg[wit={G11}]{siddhoḥ}
		\rdg[wit={P11}]{siddheḥ}
		\rdg[wit={V3}]{siddhyaiḥ}
		}}\\+}
\tl{
\pada{\app{\lem[wit={N3,J5,Delta,Epsi,P11,C6},alt={parimathitād}]{parimathitā\skp{d}}
		\rdg[wit={G4}]{mathitā pari}
		\rdg[wit={V3}]{paripaṭhitā}
		}%
	\app{\lem[wit={N3ac,J5,Delta,Epsi},alt={avakṛṣṭa}]{\skm{d }avakṛṣṭa}
		\rdg[wit={N3pc,C6}]{avakṛṣya}
		\rdg[wit={G4}]{sāra}
%		\rdg[wit={C7}]{apakṛṣṭa}
		\rdg[wit={P11}]{avakṛṣṇa}
		\rdg[wit={V3}]{kṛṣṭa}
		}% \marma
	\app{\lem[wit={Alpha,E2,Epsi,C6,V3}]{sāra}% +C7
		\rdg[wit={V19}]{sarva}
		\rdg[wit={P11}]{sārā}
		}bhūtam/}\\+} % bhūtaḥ P11
\tl{
\pada{\app{\lem[wit={N3,G4,Delta,G5,V3}]{anubhavata}
		\rdg[wit={J5}]{anubhavān}
		\rdg[wit={G11,P11}]{anubhava}
		\rdg[wit={C6}]{anubhavatu}
		}
		haṭhāmṛtaṃ % <<ha>>ṭhā° V19, ayamṛtaṃ J5
	\app{\lem[wit={N3,G4,C7,V3}]{yamīndrā}% °rāḥ N3pc, yamiṃdrā N3ac, yamīndro F
		\rdg[wit={J5}]{yogīdrā}
		\rdg[wit={V19,Epsi,P11}]{yatīndrā}
		\rdg[wit={E2}]{ya\,(text stopps here)}
		\rdg[wit={C6}]{mayedaṃ}
		}}\\+}
\tl{
\pada{yadi bhavatā%m
	\app{\lem[wit={Alpha,Delta2,P11,C6},alt={ajarāmaratva}]{\skm{m }ajarāmaratva}% ajjarā° P11
		\rdg[wit={Epsi}]{ajarāmṛtatva}
		\rdg[wit={V3}]{ajarājaraṃtvaṃ}
		}%
	\app{\lem[wit={N3,J5,Delta2,P11}]{vāñchā}
		\rdg[wit={G4}]{vāṃcchāṃ}
		\rdg[wit={G11}]{vāṃcha}
		\rdg[wit={G5}]{va[t]saḥ}
		\rdg[wit={C6}]{vāṃchāḥ}
		\rdg[wit={V3}]{vā}
		}//\versenr}%
	\myfnx{In place of this verse, the \texteta\ group manuscripts have the following as the final verse:
	\vspaceinnote\\
	\devnote{\myvin vidyātīrthe jagati vibudhāḥ sādhavaḥ satyatīrthe, %
 	gaṅgātīrthe malinamanaso yogino jñānatīrthe/\\ %
	\myvin dhārātīrthe dharaṇipatayo dānatīrthe dhanāḍhyāḥ, %
	lajjātīrthe kulayuvatayaḥ pātakaṃ kṣālayanti//} (metre: \emph{mandākrāntā})}
	\orgvnr{78}\\!}
\end{tlg}
%</vs78>
\commcite%\newpage


\iffalse
\startaltrecension
\begin{alttlg}[hp04_078_1]
%<*vs78-1>
\tl{\pada{vidyātīrthe \app{\lem[resp=emend]{jagati}\rdg[wit={J10}]{yagati}} vibudhāḥ sādhavaḥ satyatīrthe}\\+}
\tl{\pada{gaṅgātīrthe malinamanaso yogino jñānatīrthe/}\\+}
\tl{\pada{dhārātīrthe dharaṇipatayo dānatīrthe dhanāḍhyāḥ}\\+}
\tl{\pada{lajjātīrthe kulayuvatayaḥ pātakaṃ kṣālayanti//\versenr}
\sgwit{J10}\\!}
%</vs78-1>
\end{alttlg}
\endaltrecension
\fi


%<*vscol>
\begin{col}[hp04_col]
\app{\lem[nolem]{}
	\rdg[wit={Gamma,E2}]{\unavbl}
	\rdg[wit={N19}]{\om}}%
iti \app{\lem[wit={N3,J5,C7,V15,J10,V3,Jyo}]{śrī}
	\rdg[wit={G4,V19,Epsi,P11,C6}]{\om}}%
\app{\lem[alt={\post śrī}]{\skp{\post śrī}}
	\rdg[wit={N3}]{sadguru \add}
	\rdg[wit={J5}]{madguru \add}
	\rdg[wit={Epsi,V15,Jyo}]{sahajānandasaṃtānacintāmaṇinā \add}}% saṃtāra G11
\app{\lem[wit={J5,C6,V3}]{svātmārāmayogīndra} % yogiṃdra V3
	\rdg[wit={N3}]{svātmārāmayogendra}
	\rdg[wit={G4,J10}]{ātmārāmayogīṃdra}
	\rdg[wit={Delta2,Epsi}]{\om}
	\rdg[wit={V15}]{svātmārāmayogīṃdreṇa}
	\rdg[wit={P11}]{°yo° \textapp{(sic!)}}
	\rdg[wit={Jyo}]{svātmārāmayogīndra}
	}%
\app{\lem[wit={ceteri}]{viracitāyāṃ}
	\rdg[wit={N3}]{pravaracitāyāṃ}
	\rdg[wit={V19,P11}]{\om}
	} haṭhapradīpikāyāṃ
\app{\lem[alt={\ante caturtho°}]{\skp{\ante caturtho°}}
	\rdg[wit={V15}]{nādopāsanaṃ nāma \add}
	\rdg[wit={V3}]{siddhāntamuktāvalī nāma \add}
	\rdg[wit={Jyo}]{samādhilakṣaṇaṃ nāma \add}
	}%
\app{\lem[wit={Alpha,Epsi,V15,Pi,Jyo}]{caturthopadeśaḥ}
	\rdg[wit={V19}]{caturtha upadeśaḥ}
	\rdg[wit={C7}]{caturtho\{\{dhyā\}\}yam upadeśaḥ}
	\rdg[wit={J10}]{caturthodhyāyaḥ}}//~4//
%\myfn{%
%The colophon is found only in \getsiglum{Alpha,V19,C7,G11,V15,J10,Pi}. 
%\getsiglum{N23,J7,K3} have lost their last folios.
%\getsiglum{N19} has no colophon.
%\getsiglum{N19} has no colophon. Its last verse is \ref{ajananta}, which just fills fol. no. 346, and from the next folio another text begins. The colophon of \getsiglum{Jyo} reads: \devnote{iti śrīsvātmārāmayogīṃdraviracitāyāṃ haṭhayogapradīpikāyāṃ nāma caturtho'dhyāyaḥ} (Wai) or \devnote{iti śrīsajahānandasantānacintāmaṇisvātmārāmayogīṃdraviracitāyāṃ haṭhayogapradīpikāyāṃ samādhilakṣaṇaṃ nāma caturthopadeśaḥ samāptaḥ} (Tue)}
\end{col}
%</vscol>

\colcite

% N3 śrīsahajasadguruātmārāmārpaṇa saṃpūrṇaṃ |
% P11 (immeadiately followed by another text)
% C6 oṃ || pustakam idaṃ jājīrāmavyāsasya || śrīrāmaḥ
% V3 subhaṃ būyāt @ idaṃ graṃthasaṃṣyā || 502 || @ @ samvat || 16 || 93 || varṣe || || || liṣitaṃ jamadagnigirisanyāsī svātmapaṭhanārthaṃ || oṃ namaḥ śivāya ||
% V19
% C7
% G11 śrī śivāya namaḥ || śrī śaṃkarācāryagurucaraṇāravindābhyāṃ namaḥ || gaṇapati svahasta likhitaṃ | śrī gurubhyo namaḥ || oṃ | śubhaṃ |
% J10 saṃvat || 16 || 83 ||


\end{ekdosis}
\end{otherlanguage}
\end{document}
